\documentclass[11pt,a4paper]{article}
\usepackage[T1]{fontenc}
\usepackage[utf8]{inputenc}
\usepackage{lmodern}
\usepackage[margin=26mm,headheight=15pt]{geometry}
\usepackage{amsmath,amssymb,amsthm,mathtools}
\usepackage{microtype}
\usepackage{xcolor}
\usepackage{booktabs,longtable,array}
\usepackage{enumitem}
\usepackage{fancyhdr}
\usepackage{aliascnt}
\usepackage{xurl}
\usepackage{hyperref}
\usepackage[nameinlink,noabbrev]{cleveref}
\usepackage{bookmark}
% Captionless longtables rewind the table counter; give every longtable its
% own hyperref anchor so no PDF destination is duplicated.
\newcounter{gzlt}
\AddToHook{env/longtable/begin}{\stepcounter{gzlt}\renewcommand{\theHtable}{lt.\arabic{gzlt}}}
\definecolor{ink}{HTML}{17324D}
\definecolor{linkblue}{HTML}{245A81}
\hypersetup{colorlinks=true,linkcolor=linkblue,citecolor=linkblue,urlcolor=linkblue,
 pdftitle={Gamma and Zeta over the Surcomplex Numbers},
 pdfauthor={AI-assisted research report prepared for Vladimir Reshetnikov},
 pdfsubject={Merged report from six manuscripts: finite lifts, Dirichlet algebra at infinity,
 Stirling and Hurwitz, phases, the horizontal tube, obstructions, transfer and RH}}
\urlstyle{same}
\setlength{\emergencystretch}{3em}
\setlength{\parskip}{4pt}
\setlength{\parindent}{15pt}
\setlist{itemsep=0.15em,topsep=0.3em,leftmargin=1.8em}
\pagestyle{fancy}
\fancyhf{}
\fancyhead[L]{\small\textsc{Gamma and zeta over the surcomplex numbers}}
\fancyhead[R]{\small September 2026}
\fancyfoot[C]{\thepage}
\renewcommand{\headrulewidth}{0.3pt}
\setcounter{tocdepth}{2}
\numberwithin{equation}{section}
\newtheorem{theorem}{Theorem}[section]
\newaliascnt{lemma}{theorem}\newtheorem{lemma}[lemma]{Lemma}\aliascntresetthe{lemma}
\newaliascnt{proposition}{theorem}\newtheorem{proposition}[proposition]{Proposition}
\aliascntresetthe{proposition}
\newaliascnt{corollary}{theorem}\newtheorem{corollary}[corollary]{Corollary}
\aliascntresetthe{corollary}
\theoremstyle{definition}
\newaliascnt{definition}{theorem}\newtheorem{definition}[definition]{Definition}
\aliascntresetthe{definition}
\newaliascnt{example}{theorem}\newtheorem{example}[example]{Example}\aliascntresetthe{example}
\newaliascnt{question}{theorem}\newtheorem{question}[question]{Question}\aliascntresetthe{question}
\theoremstyle{remark}
\newaliascnt{remark}{theorem}\newtheorem{remark}[remark]{Remark}\aliascntresetthe{remark}
\newaliascnt{warning}{theorem}\newtheorem{warning}[warning]{Warning}\aliascntresetthe{warning}
\crefname{theorem}{Theorem}{Theorems}
\crefname{lemma}{Lemma}{Lemmas}
\crefname{proposition}{Proposition}{Propositions}
\crefname{corollary}{Corollary}{Corollaries}
\crefname{definition}{Definition}{Definitions}
\crefname{example}{Example}{Examples}
\crefname{question}{Question}{Questions}
\crefname{remark}{Remark}{Remarks}
\crefname{warning}{Warning}{Warnings}
\crefname{section}{Section}{Sections}
\crefname{subsection}{Section}{Sections}
\crefname{appendix}{Appendix}{Appendices}
\crefname{part}{Part}{Parts}
\crefname{table}{Table}{Tables}
\crefname{equation}{}{}
\creflabelformat{equation}{(#2#1#3)}
\newcolumntype{P}[1]{>{\raggedright\arraybackslash}p{#1}}
% ---------------------------------------------------------------- scalars
\newcommand{\No}{\mathbf{No}}
\newcommand{\SC}{\mathbf{SC}}
\newcommand{\Oz}{\mathbf{Oz}}
\newcommand{\R}{\mathbb R}
\newcommand{\C}{\mathbb C}
\newcommand{\Q}{\mathbb Q}
\newcommand{\Z}{\mathbb Z}
\newcommand{\N}{\mathbb N}
\newcommand{\OO}{\mathcal O}
\newcommand{\mm}{\mathfrak m}
\newcommand{\TT}{\mathbb T}
\newcommand{\eps}{\varepsilon}
\renewcommand{\Re}{\operatorname{Re}}
\renewcommand{\Im}{\operatorname{Im}}
\newcommand{\st}{\operatorname{st}}
\newcommand{\pp}{\operatorname{pp}}
\newcommand{\fin}{\operatorname{fin}}
\newcommand{\supp}{\operatorname{supp}}
\newcommand{\lc}{\operatorname{lc}}
\newcommand{\cis}{\operatorname{cis}}
\newcommand{\Arg}{\operatorname{Arg}}
\newcommand{\Log}{\operatorname{Log}}
\newcommand{\Res}{\operatorname{Res}}
\newcommand{\FP}{\operatorname{FP}}
\newcommand{\ord}{\operatorname{ord}}
\newcommand{\Frac}{\operatorname{Frac}}
\newcommand{\ev}{\operatorname{ev}}
\newcommand{\RH}{\mathsf{RH}}
% ---------------------------------------------------------------- exponentials and phases
\newcommand{\Exp}{\operatorname{Exp}}
\newcommand{\Ex}{\operatorname{Exp}_{\mm}}
\newcommand{\Lgm}{\operatorname{Log}_{\mm}}
\newcommand{\Efp}{E_{\mathrm{fp}}}
\newcommand{\Sin}{\operatorname{Sin}}
\newcommand{\Cos}{\operatorname{Cos}}
\newcommand{\cI}{\mathcal I}
% ---------------------------------------------------------------- sums
\newcommand{\Sh}{\mathop{\textstyle\sum}\nolimits^{\mathrm H}}
\newcommand{\Ph}{\mathop{\textstyle\prod}\nolimits^{\mathrm H}}
\newcommand{\Scoef}{\mathop{\textstyle\sum}\nolimits^{\mathrm{coef}}}
% ---------------------------------------------------------------- functions
\newcommand{\LSt}{L_{\mathrm{St}}}
\newcommand{\GSt}{\Gamma_{\mathrm{St}}}
\newcommand{\SSt}{S_{\mathrm{St}}}
\newcommand{\psiSt}{\psi_{\mathrm{St}}}
\newcommand{\zD}{\zeta_{\mathrm D}}
\newcommand{\zH}{\zeta_{\mathrm H}}
\newcommand{\zHh}{\widehat\zeta_{\mathrm H}}
\newcommand{\Gp}{\Gamma^{\#}}
\newcommand{\zp}{\zeta^{\#}}
\newcommand{\xip}{\xi^{\#}}
\newcommand{\DD}{\mathcal D}
\newcommand{\Arith}{\mathcal A}
\newcommand{\FE}{\mathcal X}
% ---------------------------------------------------------------- domains
\newcommand{\Tp}{\mathcal T_{+}}
\newcommand{\Tm}{\mathcal T_{-}}
\newcommand{\Th}{\mathcal T_{\mathrm h}}
\newcommand{\Wing}{\mathcal W}
\newcommand{\Sinf}{\mathcal S_{\infty}}
\newcommand{\Om}{\Omega_{\infty}}
\newcommand{\Vs}{\mathcal V}
\newcommand{\DJs}{\mathcal D_{J}}
\newcommand{\Ssym}{\mathcal S_{\mathrm{sym}}}
\newcommand{\SA}{\mathcal S_{A}}
% ---------------------------------------------------------------- misc
\DeclareUrlCommand\rl{\urlstyle{tt}}
\newcommand{\src}[1]{\textup{[#1]}}
\newcommand{\merge}{\textup{[merge]}}
\newcommand{\pinA}{\texttt{804c63c}}
\newcommand{\pinB}{\texttt{905dd19}}
\newcommand{\bigO}{\mathrm{O}}

\begin{document}
\hypersetup{pageanchor=false}
\begin{titlepage}
\vspace*{12mm}
{\color{ink}\Huge\bfseries Gamma and Zeta over the\\[4pt] Surcomplex Numbers\par}
\vspace{8mm}
{\Large Finite lifts, the Dirichlet algebra at infinity, Stirling and
Hurwitz, phases on the horizontal tube, obstructions at infinite height,
and transfer of the Riemann hypothesis\par}
\vspace{9mm}
{\large Merged research report from six manuscripts\\[3pt]
AI-assisted, prepared for Vladimir Reshetnikov\par}
\vspace{4mm}
September 2026
\vspace{9mm}
\begin{abstract}
Six manuscripts, written independently in September 2026, extend the
Gamma and Riemann zeta functions to the surcomplex numbers $\SC=\No[i]$.
This report merges them after a claim-by-claim comparison. No two of them
contradict each other once four things are fixed: the domain, the
exponential (phase), the summation notion and, for deformation results,
the deformation class. The report states these conventions once and says,
at every point where two sources seem to disagree, which of them separates
the statements.

At finite arguments, Taylor--Laurent lifting preserves the exact zero and
pole divisors. The finite Riemann hypothesis is therefore equivalent to the
classical one. One unified stability theorem contains the deformation
criteria of 04, 05 and 06: every finite zero of every J-symmetric
real-infinitesimal deformation of the completed zeta function lies on the
critical line if and only if RH holds and every nontrivial zero is simple.
At positive infinite real part and finite imaginary part, every arithmetic
function, of any growth, is realized injectively by a strongly summable
Dirichlet series. Zeta has an exact Euler product there and no zeros. Its
fibres over every nonzero infinitesimal value are found explicitly, and
its jets are algebraically independent at dilations and at real shifts over
fields containing Stirling Gamma values. An infinite-parameter Hurwitz
function recovers the Stirling logarithm, and every phase gives Gamma, zeta
and $\xi$ on the whole horizontal tube $\{\Im s\text{ finite}\}$, with
reflection and the functional equation imposed as definitions on the left
half of the tube (not a continuation or uniqueness theorem). The
$\xi$ divisor there is exactly the classical nontrivial divisor.

At infinite arguments, the report proves four obstructions. Reflection
forces Gamma poles on the negative infinite periods of the phase, on the
real axis. On the vertical strip of infinite height, reflection is
incompatible with the Stirling modulus except at two exceptional abscissae
for each height, and the Stirling Gamma itself, for every phase, fails
reflection at every point of the strip. Coherent Dirichlet
continuation replicates $\zp$ on resonant galaxies. Weak symmetric
completions do not determine infinite zeros. Elementary extensions
transfer RH and the classical zero theorems, and conditionally on RH and
simplicity the nonreal de Bruijn--Newman zeros at negative infinitesimal time
escape to infinite modulus. Several results needed for the reconciliation
were established during the merge; they are marked as such and listed in
the appendix.
\end{abstract}
\vfill
\noindent\small
\textbf{Status.} An unrefereed, AI-assisted research report. Complete
arguments are given for the derived results, relative to the cited
classical and surreal foundations. No theorem here has been verified in
Lean or any other proof assistant. It does not prove the classical Riemann
hypothesis, and it does not construct a canonical or unique global zeta or
Gamma function on all of $\No[i]$. Priority is not certified for any
consequence. The 2026 proportion results cited here are preprints.
\end{titlepage}
\hypersetup{pageanchor=true}
\begingroup
\footnotesize
\setlength{\parskip}{0pt}
\tableofcontents
\endgroup
\newpage

\section{Scope, provenance and main results}\label{gz:sec:intro}

\subsection{The question has several answers}\label{gz:sec:question}
Replacing $\C$ by $\SC=\No[i]$ in a classical formula does not define Gamma
or zeta. The classical formulas use limits, integrals, infinite products,
analytic continuation, complex exponentiation and arithmetic counting.
These operations need not have unique extensions to a non-Archimedean
field, so the extension is part of the definition. Four constructions are
kept separate throughout this report.
\begin{center}\small
\begin{tabular}{P{0.19\textwidth}P{0.34\textwidth}P{0.37\textwidth}}
\toprule
Construction & Domain and mechanism & What it establishes\\
\midrule
Canonical finite lift & Finite arguments; classical germs evaluated at infinitesimals & Exact classical divisors; a finite RH equivalent to RH; a robust RH equivalent to RH plus simplicity\\[2pt]
Series at infinity & Positive infinite real part; strong Hahn sums (Dirichlet, Stirling, Hurwitz) & An injective arithmetic algebra, Euler products, zero-freeness, inverse fibres, independence, exact Stirling and Lerch identities\\[2pt]
A named global phase & Horizontal tube, right wing and galaxies; a phase $\Psi_\chi$ and a continuation rule & Gamma, zeta and $\xi$ on the tube; forced infinite divisors; obstructions at infinite height\\[2pt]
Elementary extension & A set-sized elementary model embedded in $\No[i]$ & Transfer of RH and zero theorems at infinite height, with transported operations\\
\bottomrule
\end{tabular}
\end{center}

\subsection{Six sources}\label{gz:sec:sources}
The report is built from six manuscripts, numbered 01--06 in order of
arrival. All six date from September 2026 (06 gives only the month; the
others are dated 22 September) and were written without knowledge of one
another. Each cites a pinned commit of this repository.
\begin{center}\small
\begin{tabular}{@{}lP{0.40\textwidth}lP{0.30\textwidth}@{}}
\toprule
No. & Title (subtitle shortened) & Pin & Distinctive contribution\\
\midrule
01 & \emph{Gamma and Zeta over the Surcomplex Numbers: strong summation, a Hurwitz--Gamma bridge, phase choices, and exact inverse-zeta fibers} (31 pp.) & \pinA & inverse-zeta fibres (\cref{gz:sec:inverse}); the twisted family $E_\lambda$; one-sided left zeta\\
02 & \emph{Gamma and Zeta at Surcomplex Scales} (29 pp.) & \pinA & independence at dilations with $x$, $\LSt(x)$, $\GSt(x)$ (\cref{gz:sec:dilation,gz:sec:leading}); phase families $E_\theta$\\
03 & \emph{Gamma and Zeta on the Surcomplex Horizontal Tube} (34 pp.) & \pinA & Gamma, zeta, $\xi$ on the whole tube for every phase (\cref{gz:part:tube}); independence at shifts (\cref{gz:sec:shift})\\
04 & \emph{Gamma, Zeta, and the Riemann Hypothesis over Surcomplex Numbers} (31 pp.) & \pinB & two-sign criterion; protected counts; flat Hurwitz modification; transported divisors; heat\\
05 & \emph{Gamma and Zeta over Surcomplex Fields} (28 pp.) & \pinB & reflection defect at $\tfrac12+iT$; ratio cocycle; formal-series deformations; transfer\\
06 & \emph{Gamma and Zeta over the Surcomplex Numbers: canonical finite lifts, infinite Dirichlet series, phase obstructions, and the Riemann hypothesis} (28 pp.) & \pinB & the base text; robust RH; forced poles and resonant galaxies; weak completions; heat\\
\bottomrule
\end{tabular}
\end{center}
The pins are \pinA{} $=$ \texttt{804c63c2c675be9fd4e5086460bfa795f17a6d86}
and \pinB{} $=$ \texttt{905dd19954db43ac81f76f72036520d0eadc5710}. The
six packages were placed in this directory in commit \texttt{e4f8848}.
Manuscript 06 is the base text. It is the only source covering both the
phase thread (Dirichlet algebra, Stirling and Hurwitz, phases, obstructions)
and the RH thread (finite and robust RH, transfer, heat), and its
exponential conventions are the most general. Where another source is
strictly more general, its version is used: 03 for global objects on the
tube, 01 and 02 for the Hurwitz parameter domain, 05 for formal-series
deformations and the Stirling cocycle, and 04 for the one-parameter
converse of the stability theorem.
\Cref{gz:app:provenance} records what each source contributed, where every
result of every source is printed, where the merge had to choose, and every
error found and fixed. The code and data of all six are shipped under the
prefixes \rl{01-inverse-zeta-}, \rl{02-scale-separation-},
\rl{03-horizontal-tube-}, \rl{04-rh-transfer-}, \rl{05-surcomplex-fields-}
and \rl{06-phases-and-rh-}; their manuscripts are not.

\subsection{How results are attributed}\label{gz:sec:howtoread}
Every result carries the sources that prove it in square brackets, for
example \src{01, 03}. A result proved by several sources is printed once,
with one proof. A genuinely different proof from another source is kept
as a marked \emph{second route}. The tag \merge{} marks a statement
established while reconciling the sources, not taken from any one of them.
Each such statement has a complete proof here, and it was checked again
for this report, re-derived in the same AI-assisted process; it is not
refereed. Classical theorems and results of other reports in this
collection are cited, not reproved, unless a source gives a different
proof. \emph{Derived theorem} describes a proof status. It does not
certify historical novelty.

\subsection{The main results}\label{gz:sec:main}
\begin{enumerate}[label=(\arabic*),leftmargin=2.2em]
\item \emph{Finite divisor rigidity and finite RH}
 (\cref{gz:thm:divisor,gz:thm:finiteRH}). The canonical Taylor--Laurent
 lift of a classical meromorphic function has exactly the classical zeros
 and poles on the finite plane $\OO$, with no infinitesimally displaced
 zeros. Hence $\RH_{\mathrm{fin}}$ is equivalent to RH.
\item \emph{One stability theorem} (\cref{gz:thm:stability}). RH together
 with simplicity of every nontrivial zero is equivalent to the statement
 that all finite zeros of every J-symmetric formal deformation of $\xi$, at
 every real infinitesimal parameter, lie on the critical line. It is also
 equivalent to the same statement for $\xip\pm\eta$ with one positive
 infinitesimal $\eta$. The criteria of 04, 05 and 06 are special cases.
\item \emph{The Dirichlet algebra on $\Tp$} (\cref{gz:thm:dirichlet,gz:thm:euler}).
 On $\Tp=\{\Re s>\R,\ \Im s\text{ finite}\}$, every complex arithmetic
 function gives a strongly summable Dirichlet series, and the result is an
 injective algebra homomorphism from the Dirichlet-convolution algebra. So
 $\zD$ has an exact Euler product and no zeros.
\item \emph{Inverse zeta} (\cref{gz:thm:fibre,gz:thm:arithmetic}). For
 every nonzero infinitesimal $u$, the solutions of $\zD(s)=1+u$ in $\Tp$
 form one explicit strongly summable branch for each ordinary integer.
\item \emph{Algebraic independence} (\cref{gz:thm:indep-main,gz:thm:gamma-zeta-independence}).
 At a positive infinite real $x$, the values $x$, $\LSt(x)$, $\GSt(x)$ and
 all $\zD^{(j)}(kx)$ are algebraically independent over $\C$. At a positive
 infinite monomial $X$, the shifted jets $\zD^{(k)}(X+a)$ are independent
 over a field containing all real powers of $X$, its full slow Hahn field,
 the individual $e^{-\lambda X}$, and every $\GSt(X+b)$ with $b\in\R$.
\item \emph{Stirling and Hurwitz at infinity}
 (\cref{gz:thm:stirling,gz:thm:translation,gz:thm:hurwitz-identities,gz:thm:lerch}).
 These give exact recurrence, Gauss, translation and cocycle identities,
 Bernoulli values, and the Lerch identity
 $\partial_s\zH(s,A)|_{s=0}=\LSt(A)-\tfrac12\log2\pi$.
\item \emph{Phases on the horizontal tube} (\cref{gz:thm:gamma-tube,gz:thm:zeta-tube,gz:thm:xi-tube}).
 For every phase $\Psi_\chi$ there are Gamma, zeta and $\xi$ on
 $\Th=\{\Im s\text{ finite}\}$ with recurrence, reflection, Gauss and the
 functional equation (imposed as definitions on $\Tm$; not a continuation
 or uniqueness theorem). Gamma has poles on the negative infinite part of the
 period group $\cI(\Psi_\chi)$, zeta has extra zeros on its double, and the
 $\xi$ divisor is exactly the classical nontrivial divisor.
\item \emph{Obstructions at infinite arguments}
 (\cref{gz:thm:forced,gz:thm:strip,gz:thm:resonance,gz:thm:weak}). There
 are four. Reflection forces Gamma poles at negative infinite real points
 for every phase. On the vertical strip of infinite height, reflection is
 incompatible with the Stirling modulus except at two exceptional abscissae
 $\sigma_\pm(T)$ for each height, and $E\circ\LSt$ fails reflection at
 every point of the strip, for every exponential $E$ extending $\Efp$.
 Coherent Dirichlet continuation on resonant galaxies
 replicates $\zp$ and breaks the pointwise functional equation, for every
 phase with a resonance and every Gamma factor field-valued at $-2-iT$.
 Local analyticity and the two $\xi$ symmetries do
 not determine the infinite zeros.
\item \emph{Transfer} (\cref{gz:thm:RHstar,gz:thm:escape}). RH in an
 elementary extension is equivalent to RH. The classical zero-free
 regions, counts, proportions and density bounds transfer with their
 quantifiers. Under RH and simplicity, the nonreal zeros of the transferred
 heat family at negative infinitesimal time have infinite modulus.
\end{enumerate}

\subsection{Relation to the collection}\label{gz:sec:collection}
\emph{The collection had no zeta material before these sources.} A search of
\texttt{docs/} at the placement commit finds no Riemann zeta function,
Dirichlet series, Euler product, Hurwitz zeta or de Bruijn--Newman constant.
Every source confirmed this independently. Everything about zeta in this
report is therefore new to the collection.

\emph{Gamma functions over the surreals} \cite{RGamma} constructs the
Taylor--Stirling baseline $L_0$ (\rl{sec:baseline}, \rl{eq:L0-finite},
\rl{eq:L0-small}, \rl{eq:L0-infinite}). It proves the recurrence, formal
uniqueness and Gauss formula (\rl{lem:formal-shift},
\rl{eq:baseline-recurrence}, \rl{lem:periodic}, \rl{prop:gauss0}), the
complex log-Gamma on its tube $\{x>0,\ y\text{ finite}\}$
(\rl{thm:complex-log}, \rl{eq:complex-infinite},
\rl{eq:log-infinite-complex}), the phase lemma and the exact finite-phase
domain $\Omega$ (\rl{lem:phase}, \rl{thm:phase-domain},
\rl{eq:omega-domain-intro}), and the convex gauge family with the sharp
classification (\rl{cor:headline}, \rl{eq:explicit-gauge}, \rl{thm:main},
\rl{thm:independent}). These results are cited here, not claimed. Where a
source reproves one of them by a different route, the route is printed and
marked. The report's warning at its phase domain declines to choose a
global exponential: ``Such choices define a different problem.'' This
report takes up that different problem. It names the phase
(\cref{gz:sec:phases}), shows that outside $\Om$ every value of Gamma is
phase data (\cref{gz:prop:phase-locus}), and gives the vertical-strip
obstruction (\cref{gz:thm:strip}). The warning was a non-claim of that
report, not an open question in it, so nothing there is re-scoped.

\emph{Trigonometry on the surcomplex plane} \cite{RTrig} contains the
Ehrlich--Kaplan canonical exponential $\Exp$ with kernel $2\pi i\Oz$
(\rl{trigonometry:def:globaltrig}, \rl{trigonometry:thm:globalexp}), the
zero classes of $\Sin$ (\rl{trigonometry:cor:globalzeros}), the
classification of all phases $\Psi_\chi$ (\rl{trigonometry:thm:characters},
\rl{trigonometry:per:prop:exponentials}), the family at $\omega$
(\rl{trigonometry:ex:phases}), unavoidable infinite periods
(\rl{trigonometry:thm:infiniteperiods}, \rl{trigonometry:per:thm:cofinal})
and the normalized period group $\cI(\Psi)$
(\rl{trigonometry:per:thm:dictionary}). Sources 01, 02, 04 and 05 do not
cite these results. 03 cites Ehrlich--Kaplan but not the report. 06 cites
Ehrlich--Kaplan and the trigonometry report as a whole, for $\cis$, the
classification of phases and their infinite periods, but none of its labels.
Four identifications follow. 01's $E_0$ and 02's $E_{\ell=0}$ are $\Exp$;
03's admissible phases are the $\Psi_\chi$; 03's period group $Z_\phi$ is
$\cI(\Psi_\chi)$; and the forced-pole theorems of 03 and 06 are
consequences of \rl{trigonometry:thm:infiniteperiods}.

\emph{Surcomplex analysis} \cite{RAnalysis} contains the halo
(\rl{a:def:halo}), the canonical lift of holomorphic germs
(\rl{c:p4:def-lift}), the no-displaced-zeros proposition for holomorphic
germs (\rl{c:p4:liftzeros}), coherent sections
(\rl{c:p4:def-HU}), the canonical exponential (\rl{e:def-canonicalexp},
\rl{e:thm-exp}), the twisted family $E_\lambda$ (\rl{e:thm-twisted}) and the
derivation obstruction (\rl{b:derivation-obstruction}). 01's family
$E_\lambda$ is literally \rl{e:thm-twisted}, with the same symbol and
formula. According to the ledger \texttt{docs/FORMALIZATION.md}, some
imported infrastructure is Lean-proved: the finite-variable and univariate
clauses of the substitution lemma \rl{a:cor:complexsub}, near-one binomial
powers, and the fine derivative of evaluated ordinary power series. No
theorem of this report is itself formalized.

\emph{Stale statements corrected.} Sources 01--03 were written at \pinA{}
and 04--06 at \pinB; both pins are ancestors of the placement commit.
Since both pins the gamma-functions and analysis reports are unchanged. The
trigonometry report had no section on noncanonical exponentials at either
pin (\rl{trigonometry:per:}); it now has one, and this report cites it. The
sources' statements that the collection has no zeta material were true at
their pins and remain true apart from this report. 06 itself adds that an
empty search is not evidence of absence. The sources' descriptions of the
neighbouring reports are accurate, and no source quotes a label or line
number of the collection inaccurately. The one stale claim concerns a
source's \emph{own} construction: the remark after 01's left-zero theorem,
corrected in \cref{gz:rem:stale01}.

\section{Conventions and the table of notions}\label{gz:sec:conventions}

This section fixes one meaning for every symbol and word that the sources
use in different ways. \Cref{gz:tab:rename} lists every renaming.

\subsection{Scalars, orders and normal forms}\label{gz:sec:scalars}
We work in NBG with global choice, or in a universe-relative version of it.
Every normal form, support, family and workspace used in a construction
is set-sized, no sum over a proper class occurs, and every derivative
order, multiplicity and summation index is an ordinary integer. The
surcomplex field is $\SC=\No[i]=\{x+iy:x,y\in\No\}$ (01, 02 and 04 write
$\No(i)$, which is the same field). The letter $K$ is reserved for set-sized
Hahn workspaces $K_\Gamma=\C((t^\Gamma))$, as in \texttt{docs/NOTATION.md};
04's workspace is written $K_\Lambda$.

The modulus is $|x+iy|=(x^2+y^2)^{1/2}$. Put
$\OO=\{z:|z|<r\text{ for some }r\in\R_{>0}\}$ and
$\mm=\{z:|z|<r\text{ for all }r\in\R_{>0}\}$, with real versions $\OO_\R$
and $\mm_\R$. Every $z\in\OO$ is $z=\st z+\eps$ with $\st z\in\C$ and
$\eps\in\mm$. We write $x>\R$ for ``$x>r$ for every $r\in\R$'',
$a\prec b$ for $a/b\in\mm$, and $a\sim b$ for $a/b\in1+\mm$. The
\emph{halo} of an ordinary set $U\subseteq\C$ is $U^\#=\st^{-1}(U)$
(\rl{a:def:halo}). The \emph{monad} of a point $b$ is $b+\mm$, and its
\emph{galaxy} is $b+\OO$, also for infinite $b$.

Monomials follow the collection's convention: $z=\sum_\gamma c_\gamma t^\gamma$
with $t^\gamma=\omega^{-\gamma}$, well-ordered support, and
$v(z)=\min\supp z$, so that larger valuation means smaller magnitude.
Every real surreal splits as $x=\pp(x)+\fin(x)$. Here $\pp(x)$ is the
\emph{purely infinite part}, the terms with positive Conway exponent, and
$\fin(x)\in\OO_\R$ is the rest. Note that $\fin$ keeps the infinitesimal
terms and is not the standard part. The purely infinite surreals form the
real vector space $\Pi$ (03's $\mathbb J$, 02's $\No^\uparrow$), and the
omnific integers are $\Oz=\Pi\oplus\Z$ (\cite{RFound}, \rl{found:eq:omnific}).
The unit circle is $\TT(\No)=\{w\in\SC:w\bar w=1\}$, and $S^1(\C)$ is
the ordinary circle inside it.

The real surreal exponential $\exp_{\No}$ is Gonshor's, with inverse
$\log$ on $\No_{>0}$ \cite{Gonshor,vdDE}. We use the standard fact that
$\exp_{\No}$ maps $\Pi$ isomorphically onto the group of Conway monomials
\cite[Proposition 3.2]{EK}. In particular $e^{-X\log n}$ is a monomial for
every $X\in\Pi$ and ordinary $n\ge1$.

\subsection{Three summation notions}\label{gz:sec:sumnotions}
\begin{enumerate}[label=(\alph*)]
\item An \emph{ordinary sum} in $\C$.
\item A \emph{strong Hahn sum} $\Sh$ (\cref{gz:def:strong}). A strong
 product $\Ph$ means the strong sum of its expanded family of finite choices.
\item A \emph{coefficientwise} sum $\Scoef$: ordinary complex sums inside
 each Hahn coefficient, followed by a strong sum.
\end{enumerate}
The three are different operations. At ordinary $\sigma>1$ the family
$(n^{-\sigma})$ is not strongly summable, since every member sits at
exponent zero, but its ordinary sum is $\zeta(\sigma)$. At a positive
infinite $X$ the family $(n^{-X})$ is strongly summable. Every use of (c)
in this report is labelled: the coefficientwise Dirichlet and Mellin
formulas at finite $s$ (\cref{gz:sec:coefficientwise}), and the
galaxy-wise summation of \cref{gz:thm:resonance}. Fine-topological limits
are never used.

\subsection{Exponentials and phases}\label{gz:sec:expconv}
The following table is the most important convention of the report. The
sources use six names for four objects, and the word ``canonical'' in two
senses.
\begin{center}\small
\begin{longtable}{@{}P{0.20\textwidth}P{0.30\textwidth}P{0.42\textwidth}@{}}
\toprule
Symbol here & Definition & Source names\\
\midrule
\endhead
$\Ex$, $\Lgm$ & strong exponential on $\mm$, logarithm on $1+\mm$ & 01, 03 $\exp_\mm$; 04 $\exp^\#$; 05 $\operatorname{Exp}_{\mathrm{fin}}$ on $\mm$; 02 $\operatorname{Exp}$ on $\mm$\\
$\cis$ & $\cis(r+\eps)=e^{ir}\Ex(i\eps)$ on $\OO_\R$, kernel $2\pi\Z$ & all six; 04 $\cis_{\rm fin}$; \rl{trigonometry:thm:polar}\\
$\Efp$ (finite-phase exponential) & $\Efp(x+iy)=\exp_{\No}(x)\cis(y)$ for $x\in\No$, $y\in\OO_\R$; kernel $2\pi i\Z$ & 01 $\operatorname{Exp}_{\rm can}$ (``the canonical exponential''); 02 $E_{\mathrm f}$; 05 $\operatorname{Exp}_{\rm fp}$; 06 $E_{\rm fp}$; gamma-functions $E$ (\rl{eq:finite-exp})\\
$\Psi_\chi$, $E_\chi$ (a phase) & $\Psi_\chi(p+f)=\chi(p)\cis f$ for a character $\chi:\Pi\to\TT(\No)$; $E_\chi(x+iy)=\exp_{\No}(x)\Psi_\chi(y)$ & 03 admissible phase $\phi$, $E_\phi$; 06 $U$; \rl{trigonometry:thm:characters}\\
$\Exp$, $\Phi$ (canonical exponential) & $\chi=1$: $\Exp(x+iy)=\exp_{\No}(x)\cis(\fin y)$; kernel $2\pi i\Oz$ & Ehrlich--Kaplan; 01 $E_0$; 02 $E_{\ell=0}$; 03 $E_{\phi_0}$; 06 $E_0$ (``phase-forgetting''); \rl{trigonometry:thm:globalexp}, \rl{e:thm-exp}\\
$E_\lambda$ & $\chi_\lambda(p)=e^{i\lambda c_\omega(p)}$, $c_\omega$ the coefficient of $\omega$ & 01 $E_\lambda$; \rl{e:thm-twisted}; \rl{trigonometry:ex:phases} ($\chi_\tau$)\\
$E_{\lambda,M}$ & $\chi(p)=e^{i\lambda c_M(p)}$ for a fixed positive infinite monomial $M$ & 03 $\phi_\lambda$\\
$E_\theta$ & $\chi(p)=\cis(\theta(p))$, $\theta:\Pi\to\R$ real-linear & 02 $E_\ell$ (renamed: 02's $\ell$ collides with trigonometry's $\ell=c_\omega$)\\
$E$ & an additive-to-multiplicative $E:\SC\to\SC^\times$ extending $\Efp$; the modulus law is \emph{not} assumed & 06 Proposition 6.1\\
$\cI(\Psi)$ & $\{x\in\No:\Psi(2\pi x)=1\}$, the normalized period group; $\cI(\Phi)=\Oz$ & 03 $Z_\phi$; \rl{trigonometry:per:thm:dictionary}. Not trigonometry's $Z_E$, which is the kernel-multiplier ring\\
${}^*\exp$, ${}^*\zeta,\dots$ & operations transported from an elementary extension & 04 $\exp_*$; 05 ${}^*\exp$; 06 $M^*$\\
\bottomrule
\end{longtable}
\addtocounter{table}{-1}
\end{center}
Every $E$ with $|E(z)|=\exp_{\No}(\Re z)$ that agrees with $\Efp$ at finite
imaginary arguments is $E_\chi$ for exactly one $\chi$
(\rl{trigonometry:per:prop:exponentials}). All the families above are
phases in this sense, and $E_\lambda=E_{\lambda,\omega}$.

\emph{Two meanings of ``canonical''.} ``Canonical finite lift'' refers to
the explicit Taylor--Laurent rule (\cref{gz:def:lift}). It is not a
uniqueness claim among all functions agreeing with $f$ on $\C$. The
\emph{canonical exponential} $\Exp$ is canonical in the sense of
Ehrlich and Kaplan \cite[\S11]{EK}: it comes from the unique initial
integer part $\Oz$. It is not unique under the analytic axioms, since every
$E_\chi$ satisfies them (\rl{trigonometry:thm:characters}), and initiality
of the kernel-multiplier ring does not force it
(\rl{trigonometry:per:thm:robustness}). 01 calls $\Efp$ ``the canonical
exponential'', and 02 calls its $\ell=0$ member merely one choice. Both
usages are replaced here by the table above. The attribute
``phase-forgetting'' (06) is kept only as a gloss for the property
$\Exp(iY)=1$ for all $Y\in\Pi$.

\emph{Powers.} For a positive surreal $a$ and $s\in\SC$ we write
$a^{-s}=\Efp(-s\log a)$ whenever $(\Im s)\log a$ is finite; this value does
not depend on the phase. Otherwise $a^{-s}=E_\chi(-s\log a)$, and the phase
$\chi$ is always displayed.

\emph{Logarithm.} $\Log z=\log|z|+i\Arg z$ with finite $\Arg z$. 01 and 05
take $\Arg\in(-\pi,\pi]$, and 06 takes $(-\pi/2,\pi/2)$ on $\Re z>0$. The
two conventions agree wherever both are used, namely on $\Re z>0$. There
$\Log(x+iy)=\log x+\Lgm(1+iy/x)$ when $y/x\in\mm$, and $E_\chi(\Log z)=z$
for every phase, since $\Arg z$ is finite.

\subsection{Domains}\label{gz:sec:domains}
\begin{center}\small
\begin{longtable}{@{}P{0.14\textwidth}P{0.38\textwidth}P{0.40\textwidth}@{}}
\toprule
Symbol & Definition & Source names\\
\midrule
\endhead
$\OO$ & finite plane & all\\
$U^\#$, $c+\mm$, $b+\OO$ & halo, monad, galaxy & 02, 04 $U^\#$; 06 galaxy\\
$\Tp$ & $\{\Re s>\R,\ \Im s\in\OO_\R\}$ & 01 $\mathcal D_\infty$; 02 $\mathcal R_\infty=\mathcal T_\infty$; 03 $\mathcal T_+$; 04 Dirichlet chart; 05 $\mathcal D_{\mathrm R}$; 06 $\mathcal D_+$\\
$\Tm$ & $\{\Re s<-\R,\ \Im s\in\OO_\R\}$ & 01 $\mathcal D_{-\infty}$; 03 $\mathcal T_-$\\
$\Th$ & $\Tm\sqcup\OO\sqcup\Tp=\{\Im s\in\OO_\R\}$, the horizontal tube & 03 $\mathcal T$. \emph{Not} gamma-functions' tube $\mathcal T=\{x>0,\ y\in\OO_\R\}$\\
$\Wing$ & $\{\Re s>\R\}$, any imaginary part & 01 $\mathcal R_\infty$; 02 ``$\Re s>\R$''; 03 ``canonical wing''\\
$\Sinf$ & $\{\Re z>0,\ |z|\text{ infinite}\}$, the Stirling half-plane at infinity & 05 ``right infinite domain''; 06 ``infinite right half-plane''\\
$\Om$ & $\{x+iy:x>\R,\ y\in\OO_\R,\ y\log x\in\OO_\R\}=\Omega\cap\{x>\R\}$ & 01 ``horn''; 02 $\Omega_\infty$; gamma-functions $\Omega$\\
$\SA$ & $\{s\in\OO:(\Im s)\log A\in\OO_\R\}$, for real $A>\R$ & 02 $\mathcal S_a$\\
$\Vs$ & $\{s:0<\Re s<1,\ \Re s\text{ finite},\ \Im s\notin\OO_\R\}$, the vertical strip of infinite height & 04 ``infinite critical heights''; 05 $\tfrac12+iT$; 06 (galaxies $\pm iT+\OO$ meet $\Vs$ but are not contained in it)\\
\bottomrule
\end{longtable}
\addtocounter{table}{-1}
\end{center}
Note that $\Tp\subset\Wing\subset\Sinf$, that $\Th\cap\Sinf=\Tp$, and
that $\Vs$ is disjoint from $\Th$ and contained in $\Sinf$. 05 used ``right
infinite'' for both $\Sinf$ and $\Tp$; the two names are now separated.

\subsection{Functions and normalizations}\label{gz:sec:fnconv}
\begin{center}\small
\begin{longtable}{@{}P{0.20\textwidth}P{0.42\textwidth}P{0.30\textwidth}@{}}
\toprule
Symbol & Definition, domain & Source names\\
\midrule
\endhead
$f^\#$; $\Gp,\zp,\xip$ & Taylor--Laurent lift on $\OO$ (\cref{gz:def:lift}) & 01, 05 $f^\sharp$; 03 $f^{\mathrm H}$\\
$\xi$ & $\tfrac12s(s-1)\pi^{-s/2}\Gamma(s/2)\zeta(s)$, entire, $\xi(0)=\xi(1)=\tfrac12$ & all six\\
$\Xi$ & $\Xi(w)=\xi(\tfrac12+iw)$ & 04\\
$H_t$ & Rodgers--Tao heat family, $H_0(z)=\tfrac18\xi(\tfrac12+\tfrac{iz}2)=\tfrac18\Xi(z/2)$ & 04, 05, 06\\
$\SSt$, $\LSt$ & $\SSt(z)=\Sh_{k\ge1}b_kz^{1-2k}$, $b_k=\frac{B_{2k}}{2k(2k-1)}$; $\LSt(z)=(z-\tfrac12)\Log z-z+\tfrac12\log2\pi+\SSt(z)$ on $\Sinf$ & 01, 05, 06 $L_{\rm St}$; 02 $L_\infty$; 03 $L$, $S$; 04 $L_S$; gamma-functions $L_0$, $\widetilde L_0$\\
$\GSt$ & $\exp_{\No}\LSt(x)$ at real $x>\R$ & 01 $\Gamma_{\rm St}$; 02 $\Gamma_\infty$; 04 $\Gamma_S$; gamma-functions $\Gamma_0$\\
$\Gamma_\chi$, $\Gamma_{\Exp}$ & $E_\chi\circ\LSt$ on $\Sinf$, extended to $\Th$ in \cref{gz:sec:gammatube} & 01 $\Gamma_{{\rm St},\lambda}$; 02 $\Gamma_\ell$; 03 $\Gamma^+_\phi$, $\Gamma_\phi$\\
$\Gamma_{\mathrm{fp}}$ & $\Efp\circ\LSt$ on $\Om$ & 01 $\Gamma_{\rm St}$ on the horn; 02 $\Gamma_{\rm can}$\\
$\Gamma_{E,\mathrm{St}}$ & $E\circ\LSt$ on $\Sinf$, $E$ as in the last table & 06\\
$\psiSt$ & $\Log z-\frac1{2z}-\Sh_{k\ge1}\frac{B_{2k}}{2k}z^{-2k}$ & 01 $\psi_{\rm St}$; 02 $\psi_\infty$; 03 $\psi_+$\\
$\zD$ & $\Sh_{n\ge1}n^{-s}$ on $\Tp$ & 01 $Z_\infty$; 02 $\zeta_\infty$; 03 $\zeta_+$; 04, 05, 06 $\zeta_{\mathrm D}$\\
$\zeta_\chi$, $\xi_\chi$ & zeta and $\xi$ on $\Th$ for the phase $\chi$ & 03 $\zeta_\phi$, $\xi_\phi$; 01 $Z_{-,\lambda}$ on $\Tm$\\
$\zeta^{\Wing}_\chi$ & $\Sh_nE_\chi(-s\log n)$ on $\Wing$, where summable & 02 $\zeta_\ell$; 03 canonical wing\\
$\zH(s,A)$ & $\frac{A^{1-s}}{s-1}+\frac{A^{-s}}2+\Sh_{k\ge1}\frac{B_{2k}}{(2k)!}(s)_{2k-1}A^{1-s-2k}$ & 02 $\mathcal H_\infty$; 01 $H$; 03 $\mathcal Z_\phi$; 04 $Z_H$; 06 $Z_\infty(s,A)$\\
$\zHh(s,A)$ & $A^{s-1}\zH(s,A)=\frac1{s-1}+\frac1{2A}+\Sh_{k\ge1}\frac{B_{2k}}{(2k)!}(s)_{2k-1}A^{-2k}$ & 01 $\widehat H$; \textbf{03 $H$}\\
$Q_{a,b}$, $R_{a,b}$ & ratio cocycle (\cref{gz:thm:cocycle}); $Q_a:=Q_{a,0}$ & 05; 03 $R_a$; 02 eq.\ 8.6\\
\bottomrule
\end{longtable}
\addtocounter{table}{-1}
\end{center}
The Bernoulli numbers satisfy $u/(e^u-1)=\sum_nB_nu^n/n!$, so
$B_1=-\tfrac12$, and $ue^{au}/(e^u-1)=\sum_nB_n(a)u^n/n!$. The Stirling
constant is $\tfrac12\log2\pi$ in every source. The Hurwitz letter is the
only normalization trap: \emph{03's $H$ is 01's $\widehat H$}, not 02's
$\mathcal H_\infty$ (which is $\zH$). Neither $H$ nor $Z_\infty$ is used
for Hurwitz here, since $H$ is the heat family and $Z_\infty$ is 01's
$\zD$. Likewise $\Gamma_0$ is not used: it means the gauge $a=0$ in
gamma-functions and the phase $\ell=0$ in 02.

\begin{table}[!tbp]
\caption{Renamings of source symbols. Normalizations are unchanged unless stated.}\label{gz:tab:rename}
\small
\begin{tabular}{@{}lP{0.86\textwidth}@{}}
\toprule
Source & Renamings\\
\midrule
01 & $K=\No(i)\to\SC$; $\mathcal O_K,\mathfrak m_K\to\OO,\mm$; $\operatorname{Exp}_{\rm can}\to\Efp$; $E_0\to\Exp$ ($E_\lambda$ kept); $\sin_\lambda\to\Sin_{\chi_\lambda}$; horn $\to\Om$; $\mathcal D_\infty\to\Tp$; $\mathcal D_{-\infty}\to\Tm$; $\mathcal R_\infty\to\Wing$ (for $\LSt$: $\Sinf$); $Z_\infty(s)\to\zD(s)$; $Z_{-,\lambda}\to\zeta_{\chi_\lambda}$ on $\Tm$; $\widehat H\to\zHh$, $H(s,A)\to\zH(s,A)$; $H_A(S)\to\zH(S,A)$ (a formal germ); $P_\lambda$ kept; $A_\omega$ kept (a constant, not a Hurwitz parameter); $\Gamma_{\mathrm{St},\lambda}\to\Gamma_{\chi_\lambda}$; $\Gamma_{\rm St}$ on the horn $\to\Gamma_{\rm fp}$; $R(z)\to\SSt(z)$; $c_\Gamma\to\tfrac12\log2\pi$; $f^\sharp\to f^\#$; $[\omega]b\to c_\omega(b)$; $\mathcal E_s\to\DD_s$.\\
02 & $\No(i)\to\SC$; $\OO_\C,\mm_\C\to\OO,\mm$; $E_{\mathrm f}\to\Efp$; $\No^\uparrow\to\Pi$; $\ell\to\theta$, $E_\ell\to E_\theta$, $\pi_\ell\to\pi_\theta$, $\Gamma_\ell\to\Gamma_{\chi_\theta}$, $\zeta_\ell\to\zeta^\Wing_{\chi_\theta}$; $\Gamma_0,\zeta_0\to\Gamma_{\Exp},\zeta^\Wing_{\Exp}$; $\mathcal R_\infty=\mathcal T_\infty\to\Tp$; $\Omega_\infty$ kept; $\mathcal S_a\to\SA$; $\zeta_\infty\to\zD$; $L_\infty,\Gamma_\infty,\psi_\infty\to\LSt,\GSt,\psiSt$; $\mathcal H_\infty\to\zH$; $\Gamma_{\rm can}\to\Gamma_{\rm fp}$; $\mathcal F_x\to\mathcal F_x^{\rm Dir}$; $\mathcal D_s$, $\mathcal A$ kept.\\
03 & $\SC=\No[i]$ kept; $\mathbb J\to\Pi$; $\phi\to\Psi_\chi$, $E_\phi\to E_\chi$; $\phi_0\to\Phi$ ($E_{\phi_0}\to\Exp$); $\phi_\lambda\to\Psi_{\chi_{\lambda,M}}$; $Z_\phi\to\cI(\Psi_\chi)$; $\sin_\phi,\cos_\phi\to\Sin_\chi,\Cos_\chi$; $\mathcal T,\mathcal T_\pm,\OO_\C\to\Th,\mathcal T_\pm,\OO$; $f^{\mathrm H}\to f^\#$; $\zeta_+\to\zD$; $L,S\to\LSt,\SSt$; $\Gamma^+_\phi\to\Gamma_\chi$ on $\Tp$; $g_\phi\to g_\chi=1/\Gamma_\chi$; $\chi_\phi(s)\to\FE_\chi(s)$ (the letter $\chi$ is the character); $R_a\to Q_a$; $H(s,A)\to\zHh$, $\mathcal Z_\phi(s,A)\to\zH(s,A)$ with $A^w=E_\chi(w\log A)$; $F_X\to F^{\rm Hahn}_X$; $M_X,E_X,E_X^\Gamma,T_X$ kept; $\psi_+\to\psiSt$.\\
04 & $\No(i)\to\SC$; $\mathbb K=\C((t^\Lambda))\to K_\Lambda$; $\cis_{\rm fin}\to\cis$; Dirichlet chart $\to\Tp$; $L_S,\Gamma_S\to\LSt,\GSt$; $Z_H\to\zH$; $\exp_*,\Gamma_*,\zeta_*,\xi_*,N_*\to{}^*\exp,{}^*\Gamma,{}^*\zeta,{}^*\xi,{}^*N$; $\Lambda_{\rm DN}$ kept; $\mathcal E_q$ kept; $\chi$ of its Proposition 9.1 is a set-sized character in the sense of \rl{trigonometry:thm:characters}.\\
05 & $\mathbb K\to\SC$; $\operatorname{Exp}_{\rm fin}\to\Ex$ on $\mm$ and $\cis$ on $\OO_\R$; $\operatorname{Exp}_{\rm fp}\to\Efp$; ``right infinite domain'' $\to\Sinf$; $\mathcal D_{\mathrm R}\to\Tp$; $\mathcal E_s\to\DD_s$; gauge $\Gamma_\lambda=e^{h_\lambda}\Gamma\to e^{h^{05}_\lambda}\GSt$ with $h^{05}_\lambda=\lambda e^{-\pp}$; $\Xi$ of its Theorem 9.1 $\to$ a symmetric function $F$; $\mathcal A,{}^*\mathcal A\to\mathfrak A,\mathfrak M^*$; the four-point set $A$ in its proof of Theorem 9.1 $\to\mathcal Q$.\\
06 & $\mathbb K\to\SC$; $E_0\to\Exp$; $U\to\Psi_\chi$; $\operatorname{Sin}_U\to\Sin_\chi$; $\mathcal D_+\to\Tp$; $Z_\infty(s,A)\to\zH(s,A)$; $\Gamma_{E,\rm St}$ kept; gauge $h_\lambda\to h_{a_\lambda}$ with $a_\lambda(m)=\lambda e^{-m}$ (gamma-functions \rl{cor:headline}); class $\mathcal S\to\Ssym$; $\mathcal M^*,\xi^*\to\mathfrak M^*,{}^*\xi$; parametric $\kappa$ of its Corollary 10.2 $\to c_*$.\\
Local & Letters that are reserved above keep one meaning; local objects get their own: the Lagrange function $\Theta$ and Rodgers--Tao kernel $\Phi_{\rm RT}$ (not the phase $\Phi$); the de Bruijn--Newman constant $\Lambda_{\rm DN}$ (not von Mangoldt's $\Lambda$ or the group of $K_\Lambda$); word length $|\nu|$, tuple length $d$ and a finite index set $\mathcal K$ (not a workspace $K$); $F(x)$ written out in \cref{gz:lem:primeindependence}; the Platt--Trudgian height $T_{\rm PT}$ and the factor $P_U$ of \cref{gz:thm:weak} (not the heat family $H$); monomials $\mathfrak t_n$, the map $\nu(P)$ and the prime projection $\operatorname{pr}$ with derivation $\delta_Z$ (not M\"obius $\mu$ or zeros $\rho$); the Hartogs number $\mathfrak h$ and constants $\kappa_0,\kappa_{\rm s},c,c_*$; the field map $\iota_c$ and the $\Q$-linear maps $\mathfrak s,\mathfrak s'$ (not $\tau$ or $\sigma$); the imaginary part $y$ in \cref{gz:thm:fibre}; the index letters $i,b$ (not $\ell$); the structure $\mathfrak A$ and the set $\mathcal Q$ (not $\Arith$); the Hermitian matrix $\Delta$ (not $E$).\\
\bottomrule
\end{tabular}
\end{table}

\subsection{Tempting false readings}\label{gz:sec:false}
Each line pairs a false reading with the true statement it resembles.
\begin{itemize}
\item \emph{False:} $\Exp$ is the unique global exponential. \emph{True:}
 $\Exp$ is canonical in the integer-part sense. Every $E_\chi$ satisfies the
 same group, modulus and local laws (\cref{gz:sec:expconv}).
\item \emph{False:} Gamma is undefined outside $\Om$. \emph{True:} the
 phase-free exponential $\Efp\circ\LSt$ is defined exactly on $\Om$.
 Outside $\Om$ every phase gives a value, and every such value is phase
 data (\cref{gz:prop:phase-locus}).
\item \emph{False:} zeta has a single pole at $1$ everywhere. \emph{True:}
 on $\Th$, $\zeta_\chi$ has a single pole at $1$ (\cref{gz:thm:zeta-tube}).
 On resonant galaxies outside $\Th$, coherent Dirichlet continuation has
 poles at $1\pm iT$ (\cref{gz:thm:resonance}).
\item \emph{False:} an infinite period is a resonance. \emph{True:} every
 phase has infinite periods (\cref{gz:thm:periods}). A simultaneous
 logarithmic resonance is stronger. It exists for every phase with
 ordinary-circle values, and fails for some other phases
 (\cref{gz:prop:resonance}).
\item \emph{False:} reflection holds for Stirling Gamma off the real axis.
 \emph{True:} it holds on $\Th$ by construction (\cref{gz:thm:gamma-tube}).
 On the vertical strip a general prescription with Stirling moduli
 contradicts it except at two exceptional abscissae for each height, and
 $E\circ\LSt$ itself fails it at every point of the strip, for every $E$
 extending $\Efp$ (\cref{gz:thm:strip}).
\item \emph{False:} $\zH(s,A)=\sum_{n\ge0}(A+n)^{-s}$ as a strong sum.
 \emph{True:} that family is not strongly summable (\cref{gz:warn:hurwitz-raw}).
\item \emph{False:} a strong sum is a limit of partial sums. \emph{True:}
 no set-indexed sequence that is not eventually constant converges in the
 fine topology (\cref{gz:rem:fine}).
\end{itemize}

\section{How the six sources fit together}\label{gz:sec:fit}

A claim-by-claim comparison of the six sources found no genuine
contradiction. Every apparent conflict comes from a different domain,
exponential, summation notion or deformation class. \Cref{gz:tab:fit}
lists each apparent conflict and what separates the two sides. The same
reconciliation is also printed at the point of use.
\begin{center}\small
\begin{longtable}{@{}P{0.30\textwidth}P{0.44\textwidth}P{0.18\textwidth}@{}}
\caption{Apparent conflicts and their resolution.}\label{gz:tab:fit}\\
\toprule
Apparent conflict & Resolution & Where\\
\midrule
\endfirsthead
\toprule
Apparent conflict & Resolution & Where\\
\midrule
\endhead
Divisor preservation (all six) & Identical statements; printed once & \cref{gz:thm:divisor}\\
Three stability criteria: 04 linear J-symmetric, one fixed $\eta>0$; 05 formal series, both symmetries, $\tau>0$; 06 linear, both symmetries, both signs & Subclasses of one class $\DJs$ (06's lies in 05's; 04's and 05's are incomparable); one theorem, with the one-$\eta$ converse & \cref{gz:thm:stability}\\
04: exact and shadow finite RH are ``inequivalent in strength'' & The exact form implies the shadow form, and is equivalent to it together with simplicity of every nontrivial zero. They can differ in truth value only if RH holds and some nontrivial zero is multiple, which is not known to happen & \cref{gz:rem:shadow}\\
02 Warning 6.4: zero-freeness is confined to $\Tp$; 02 Prop.\ 11.3 ($E_\theta$) and 03 Prop.\ 13.1 ($\Exp$ only): zero-free on $\Wing$ & The warning concerns the raw $\zD$; the wing sums use a named phase & \cref{gz:sec:wing}\\
03: Gamma with reflection on $\Th$; 06: no Gamma with only ordinary poles satisfies reflection & The same fact: poles are forced on $\cI(\Psi_\chi)$, and 03's pole set satisfies 06's requirement & \cref{gz:thm:forced}\\
03: $\zeta_\chi$ has a single pole at $1$; 06: poles at $1\pm iT$ & Different domains ($\Th$ versus resonant galaxies outside it) and different axioms & \cref{gz:thm:resonance}\\
03: reflection holds; 05: reflection fails at $\tfrac12+iT$ & $\tfrac12+iT\notin\Th$; 05's obstruction concerns $\Vs$, and 06's $\Gamma_{E,\rm St}$ fails reflection at every point of $\Vs$ & \cref{gz:thm:strip}\\
01: $Z_{-,\lambda}(-\omega)$ zero or not, ``nothing canonical decided''; 03: $\zeta_{\Exp}(-\omega)=0$ & Phase choice. With Ehrlich--Kaplan's $\Exp$ ($\lambda=0$) the zero is present; no functional equation forces it & \cref{gz:prop:moving,gz:rem:stale01}\\
03: poles on $\Oz_{\le0}$; 04, 06: poles at internal integers $-{}^*\N$ & Different constructions: the period group of a phase versus the transported integer sort & \cref{gz:sec:transported-divisors}\\
``Canonical exponential'' means $\Efp$ in 01 and $\Exp$ in 03, 06; 02: ``no canonical choice'' & Canonical in the integer-part sense is not unique & \cref{gz:sec:expconv}\\
Gamma non-uniqueness: 01, 02, 03, 06 (phase); 04, 05 (flat gauges); 06, gamma-functions (convex gauges) & Three independent mechanisms; the flat gauges of 04 and 05 break Gauss, the convex gauges keep it & \cref{gz:sec:nonunique}\\
Hurwitz $H$ in 01, 02, 03 & 03's $H$ is 01's $\widehat H$; 02's $\mathcal H_\infty$ is 01's $H$; all are $\zH$ or $\zHh$ & \cref{gz:sec:fnconv}\\
Hurwitz domains $\SA$ (01, 02), all finite $s$ (03), germs (04, 06) & $\zH$ is phase-free exactly on $\SA$; elsewhere it depends on the phase and all identities hold for every phase & \cref{gz:sec:hurwitz}\\
Two distribution identities (01, 03) & The same identity after $A\mapsto A/q$ & \cref{gz:thm:hurwitz-identities}\\
02: independence at dilations over $\C$ and $\mathcal F_x^{\rm Dir}$; 03: at shifts over $E_X^\Gamma$ & Incomparable families and base fields; both kept; the combined statement is open & \cref{gz:sec:indep-compare}\\
05 ``right infinite'' twice & $\Sinf$ versus $\Tp$ & \cref{gz:sec:domains}\\
$\Arg$ in $(-\pi,\pi]$ (01, 05) or $(-\pi/2,\pi/2)$ (06) & The conventions coincide on $\Re z>0$ & \cref{gz:sec:expconv}\\
05's cocycle $R_{a,b}$ as a Gamma ratio & A formal identity; the classical reading holds only away from the negative real direction & \cref{gz:rem:cocycle-gloss}\\
Heat variables: $\Xi(w)=\xi(\tfrac12+iw)$ (04 stability) versus $H_0$ (all) & $H_0(z)=\Xi(z/2)/8$ & \cref{gz:sec:heat}\\
04: the heat theorem's nonreal zero under RH; 05, 06: unconditional & Hypotheses attached clause by clause & \cref{gz:thm:escape}\\
Platt--Trudgian simplicity: 05 uses it; 06 ``no simplicity inferred'' & The published statement includes simplicity; 06's argument does not need it, and \cref{gz:cor:verified} (from 05) uses it & \cref{gz:sec:verified}\\
$N(\sigma,T)$ counts $|\Im\rho|\le T$ (Guth--Maynard, 04) or $0<\Im\rho\le T$ (05, 06) & The one-sided count is smaller, so the bounds hold for it & \cref{gz:sec:density}\\
Guth--Maynard: arXiv (04, 06) versus \emph{Annals} (05) & Published: \emph{Ann.\ of Math.}\ 203 (2026) & \cref{gz:sec:density}\\
Proportions $0.417293$, $0.407511$, $>5/12$, $C_0$ & Different quantities (on the line versus simple on the line) and statuses (published versus preprint) & \cref{gz:sec:proportions}\\
$\Gamma_0$ in 02 and in gamma-functions & Renamed; neither is used & \cref{gz:sec:fnconv}\\
01's $E_0$, 02's $E_{\ell=0}$, 03's $E_{\phi_0}$, 06's $E_0$ & All equal $\Exp$ & \cref{gz:sec:expconv}\\
06: the transported phase is not $\Exp$ & A first-order two-frequency test separates them & \cref{gz:prop:twofreq}\\
06 Remark 7.4: an infinite period need not give a resonance (unproved) & True: \cref{gz:prop:resonance}(c) constructs, using choice (a basis of $\R$ over $\Q(\log3/\log2)$), a phase with no resonance & \cref{gz:prop:resonance}\\
01: $\Efp(\LSt(z))$ defined iff $y\log x$ finite; $\Exp$ is defined everywhere & The statement concerns $\Efp$; the phase-free locus of all phases is exactly $\Om$ & \cref{gz:prop:phase-locus}\\
\bottomrule
\end{longtable}
\end{center}

\part{Foundations}\label{gz:part:foundations}

\section{Strong summation and infinitesimal evaluation}\label{gz:sec:summation}

\subsection{Strong summability}
Normal forms are Conway's \cite{Conway,Gonshor}.
\begin{definition}[Strong summability \src{01--06}]\label{gz:def:strong}
A set-indexed family $(f_j)_{j\in I}$ of Hahn series (or of normal forms) is
\emph{strongly summable} if $\bigcup_j\supp f_j$ is well ordered and every
exponent lies in the support of only finitely many $f_j$. Its strong sum
$\Sh_jf_j$ is formed by finite coefficient addition at each exponent.
\end{definition}

The second condition is substantive: a convergent ordinary series of
nonzero complex constants is not strongly summable, because every member
contributes at exponent zero. For $\sum_{n\ge1}2^{-n}$ the ordinary value
$1$ comes from ordinary analysis, not from \cref{gz:def:strong}. Conversely
$\Sh_{n\ge0}n!\,\eps^n$ exists for every $\eps\in\mm$, although the
classical series has radius zero \src{01, 02}. The converse fails: nothing
here licenses substituting a finite non-infinitesimal value into an
arbitrary formal series \src{02}.

\begin{lemma}[Neumann's support lemma \cite{Neumann}]\label{gz:lem:neumann}
Let $S$ be a well-ordered subset of the positive cone of an ordered abelian
group. The additive monoid $\langle S\rangle$ is well ordered, and each of
its elements has only finitely many representations as a finite sum of
elements of $S$, over all word lengths. If $A,B$ are well ordered, so is
$A+B$, and every element of $A+B$ has finitely many representations
$a+b$.
\end{lemma}
The lemma is imported; the collection proves it in the analysis report
(\rl{a:sec:neumann}) and in Lean (\texttt{Surreal/HahnSeries/Neumann.lean},
as recorded in \texttt{docs/FORMALIZATION.md}).

\begin{lemma}[Infinitesimal substitution \src{01--06}]\label{gz:lem:substitution}
Let $\eps_1,\dots,\eps_d\in\mm$ and $F\in\C[[X_1,\dots,X_d]]$, with no
condition on the growth of its coefficients. Then $F(\eps_1,\dots,\eps_d)$
is defined by a strongly summable family. Evaluation respects sums and
products, and it commutes with formal composition when every inner series
has zero constant term. A series with nonzero constant term evaluates to a
unit of $\OO$. The same holds for a Laurent series with finitely many
negative powers in one nonzero infinitesimal.
\end{lemma}
\begin{proof}
Let $S$ be the union of the supports of the nonzero $\eps_j$, a finite
union of well-ordered subsets of the positive cone. Every exponent in
$\eps^\alpha$ is a sum of $|\alpha|$ elements of $S$, so
$\supp\eps^\alpha\subseteq\langle S\rangle$. By \cref{gz:lem:neumann}, a
fixed exponent has finitely many representations over all word lengths.
Assigning each letter of a word to one of the $d$ variables leaves finitely
many possibilities, so each exponent receives finitely many
contributions. Products and compositions are regrouped exponent by
exponent, and each coefficient identity is a finite algebraic identity.
The constant term of the evaluation is $F(0)$. A finite principal part
multiplies the evaluation by a fixed nonzero power of $\eps$.
\end{proof}
The finite-variable clauses are the analysis report's \rl{a:cor:complexsub},
whose finite-variable part is recorded as Lean-proved in
\texttt{docs/FORMALIZATION.md}. The strong exponential and logarithm
\[
\Ex(u)=\Sh_{n\ge0}\frac{u^n}{n!}\quad(u\in\mm),\qquad
\Lgm(1+u)=\Sh_{n\ge1}\frac{(-1)^{n+1}u^n}{n}
\]
are mutually inverse bijections $\mm\leftrightarrow1+\mm$. They satisfy
the exponential laws and commute with conjugation. For $a\in\C$ the
near-one power is $(1+u)^a=\Ex(a\Lgm(1+u))$, equivalently the binomial
series.

\begin{lemma}[Mixed supports \src{02, 03, 04}]\label{gz:lem:mixed}
Let $(\mathfrak t_n)$ be distinct monomials whose valuations form a well-ordered
set $B$, let $\eps_1,\dots,\eps_r\in\mm$, and let $F_n\in\C[[T_1,\dots,T_r]]$
be arbitrary. Then the family $\bigl(\mathfrak t_n\,[T^\alpha]F_n\,\eps^\alpha\bigr)_{n,\alpha}$
is strongly summable.
\end{lemma}
\begin{proof}
With $S$ as in the previous proof, $\supp\eps^\alpha$ lies in the
$|\alpha|$-fold sums of elements of $S$, hence in $\langle S\rangle$.
Every exponent of the family lies in $B+\langle S\rangle$, which is well
ordered by \cref{gz:lem:neumann}. A fixed exponent has finitely many
decompositions $\beta+\gamma$ with $\beta\in B$ and $\gamma\in\langle S\rangle$.
Distinctness of the $\mathfrak t_n$ gives at most one $n$ for each $\beta$, and for
each $\gamma$ only finitely many words, hence finitely many $\alpha$.
Complex coefficients do not change supports.
\end{proof}
This spells out the ``words'' argument of 02's Lemma 2.2(b), which treated
$\eps^\alpha$ as a family of monomials; the fix is the inclusion
$\supp\eps^\alpha\subseteq\langle S\rangle$ just used (error 8 of
\cref{gz:tab:errors}).

\begin{lemma}[Locally finite real exponents \src{01}]\label{gz:lem:realpowers}
Let $A\subseteq\R_{\ge0}$ meet every bounded interval in a finite set, let
$0<u\in\mm_\R$, and let $c_a\in\C$. Then $(c_au^a)_{a\in A}$ is strongly
summable, where $u^a=\exp_{\No}(a\log u)$. The same holds for a family with
finitely many members at each exponent. For a nonzero $w\in\mm$ put
$w^a=\Efp(a\Log w)=|w|^a\cis(a\Arg w)$; the family $(c_aw^a)_{a\in A}$ is
again strongly summable.
\end{lemma}
\begin{proof}
Write $u=c\,\mathfrak t(1+\eta)$ with real $c>0$, a Conway monomial
$\mathfrak t<1$ and $\eta\in\mm_\R$. Then $u^a=c^a\mathfrak t^a(1+\eta)^a$.
The set $\{\mathfrak t^a:a\in A\}$ is reverse well ordered, with the order
type of $A$. The binomial expansion of $(1+\eta)^a$ has support in
$\langle\supp\eta\rangle$, independently of $a$. \Cref{gz:lem:neumann}
gives both conditions of \cref{gz:def:strong}, and finite multiplicities
change nothing. For complex $w$ write $\Arg w=\theta_0+\delta$ with
$\theta_0\in\R$ and $\delta\in\mm_\R$. Then
$w^a=|w|^ae^{ia\theta_0}\Ex(ia\delta)$, and the extra factor has support in
$\langle\supp\delta\rangle$, independently of $a$.
\end{proof}

\begin{lemma}[Evaluation of weighted series \merge]\label{gz:lem:evalhom}
Let $(x_i)_{i\in I}$ be a set of formal variables with real weights
$w_i>0$ such that for every $R$ only finitely many finitely supported
multi-indices $\nu$ have $w(\nu)=\sum_i\nu_iw_i\le R$. Let $\C[[x]]$ be the
ring of all formal series $F=\sum_\nu f_\nu x^\nu$. Fix constants
$c_i\in\C$ and $0<u\in\mm_\R$, or a nonzero $u\in\mm$ with the powers of
\cref{gz:lem:realpowers}. Then
\[
\ev(F)=\Sh_\nu f_\nu c^\nu u^{w(\nu)}
\]
is defined, and $\ev$ is a ring homomorphism. It is strongly additive: if
$(F_j)$ is a family in which each monomial $x^\nu$ has nonzero
coefficient in only finitely many $F_j$, then the expanded family of all
$\ev$-terms is strongly summable and $\ev(\sum_jF_j)=\Sh_j\ev(F_j)$. In
particular $\ev((1+W)^a)=(1+\ev W)^a$ for every $W$ with zero constant
term and every $a\in\C$.
\end{lemma}
\begin{proof}
The exponents $w(\nu)$ form a set meeting each bounded interval finitely,
with finitely many $\nu$ at each exponent, so \cref{gz:lem:realpowers}
gives the first claim. The product of two strongly summable families is
strongly summable, and its strong sum is the product of the sums. Grouping
the product family by $\nu=\nu'+\nu''$, a finite grouping, gives
$\ev(FG)=\ev(F)\ev(G)$. For a formally summable family $(F_j)$, the terms
at a fixed exponent come from finitely many $\nu$ and, for each $\nu$,
from finitely many $j$. So the doubly indexed family is strongly summable
and may be regrouped either way. Finally $(1+W)^a=\sum_j\binom ajW^j$ is
formally summable because $W^j$ has only monomials of weight at least
$j\min_iw_i$. Applying the two previous properties and
\cref{gz:lem:substitution} to the infinitesimal $\ev W$ gives
$\ev((1+W)^a)=\Sh_j\binom aj(\ev W)^j=(1+\ev W)^a$.
\end{proof}
This is the lemma that 01's existence argument for the near-one solution
of the inverse-zeta equation uses without statement (error 3 of
\cref{gz:tab:errors}). It is used in \cref{gz:lem:Vunique}.

\begin{remark}[Strong sums are not fine limits \src{01, 03, 05, 06}]\label{gz:rem:fine}
The fine topology on $\SC$ has balls of every positive surreal radius
(\rl{a:def:fine}). A set-indexed net that is not eventually constant cannot
converge in it, since a positive surreal smaller than every nonzero distance
in a set of values exists by the cut property \src{01, 03}. In particular
a sequence of ordinary complex partial sums converges to a finite
surcomplex number only if it is eventually constant. Fine convergence
forces the standard parts to converge ordinarily, and then an infinitesimal
radius separates every other term \src{05}. A fixed set-sized Hahn field
with its own valuation topology need not behave this way. No strong sum in
this report is a fine limit, and convergence in one topology is never used
to justify summation in another. The same holds for strong products.
\end{remark}

\subsection{Derivatives}\label{gz:sec:derivatives}
All six sources differentiate \emph{in the argument}: $f'(a)$ is the
coefficient of $h$ in the local expansion of $f(a+h)$, and $D_z\omega=0$.
This is not the Berarducci--Mantova derivation of scalars \cite{BMder},
which has $\partial\omega=1$. The two notions are incompatible, not just
normalized differently (\rl{b:derivation-obstruction}): $\Exp(i\omega)=1$
cannot be differentiated with $\partial\omega=1$. Every local germ in this
report is a series in $\SC[[h]]$ that is strongly evaluable on some fine
ball. By the analysis report's germ theorem (\rl{b:realization},
\rl{b:germalgebra}), its formal derivative is its fine derivative on a
small enough ball, and sums, products, quotients and chain rules hold.
This settles the outline argument of 03's Section 2.3 (error 19). The
admissible ball can be much smaller than $\mm$: near an infinite scalar
$A$, an increment may need $h\log A\in\mm$ \src{02, 03}. No source, and no
statement here, asserts compatibility with the Berarducci--Mantova
derivation.

\section{Exponentials, phases and period groups}\label{gz:sec:phases}

\subsection{Finite angles and the finite-phase exponential}\label{gz:sec:cis}
\begin{proposition}[Polar form and finite-phase exponential \src{01, 02, 03, 05, 06}]\label{gz:prop:polar}
The map $\cis(r+\eps)=e^{ir}\Ex(i\eps)$ ($r\in\R$, $\eps\in\mm_\R$) is a
homomorphism from $(\OO_\R,+)$ onto $\TT(\No)$ with kernel $2\pi\Z$. Every
$w\in\TT(\No)$ is $\cis\theta$ for a finite $\theta$, unique modulo $2\pi\Z$.
The finite-phase exponential $\Efp(x+iy)=\exp_{\No}(x)\cis(y)$, for
$x\in\No$ and $y\in\OO_\R$, is a homomorphism with kernel $2\pi i\Z$ that
commutes with conjugation and satisfies $|\Efp(z)|=\exp_{\No}(\Re z)$.
Every $z\ne0$ has a finite principal argument, and $\Efp(\Log z)=z$.
\end{proposition}
\begin{proof}
If $w\bar w=1$, then $w$ is finite and $|\st w|=1$. Choose an ordinary
argument $\theta_0$ of $\st w$ and put $v=e^{-i\theta_0}w\in1+\mm$. Then
$\Lgm v+\overline{\Lgm v}=\Lgm(v\bar v)=0$, so $\Lgm v=i\delta$ with
$\delta\in\mm_\R$, and $w=\cis(\theta_0+\delta)$. A kernel element of
$\cis$ has standard part in $2\pi\Z$, and its infinitesimal part is zero
because $\Ex$ is injective. The statements about $\Efp$ follow from the
real exponential laws. For $z\ne0$ apply the first part to $z/|z|$.
\end{proof}
This is the trigonometry report's \rl{trigonometry:thm:polar}. The
finite-phase exponential is the gamma-functions report's $E$
(\rl{eq:finite-exp}) and agrees with the trigonometry report's strip
exponential (\rl{trigonometry:thm:stripexp}). It exponentiates exactly the
arguments with finite imaginary part.

\subsection{Phases}\label{gz:sec:phasedef}
\begin{definition}[Phase \src{03, 06}]\label{gz:def:phase}
A \emph{phase} is an additive-to-multiplicative homomorphism
$\Psi:(\No,+)\to\TT(\No)$ that agrees with $\cis$ on $\OO_\R$. Its
exponential is $E_\Psi(x+iy)=\exp_{\No}(x)\Psi(y)$, with
$\Sin_\Psi z=(E_\Psi(iz)-E_\Psi(-iz))/(2i)$ and
$\Cos_\Psi z=(E_\Psi(iz)+E_\Psi(-iz))/2$.
\end{definition}
Since $\No=\Pi\oplus\OO_\R$, every phase is
$\Psi_\chi(p+f)=\chi(p)\cis f$ for a unique character
$\chi:(\Pi,+)\to\TT(\No)$, and every character occurs
(\rl{trigonometry:thm:characters}). We write $E_\chi$, $\Sin_\chi$,
$\Cos_\chi$. 03's admissible phases and 06's phase homomorphisms $U$ are
exactly the $\Psi_\chi$.

\begin{proposition}[Exponential calculus \src{01, 02, 03}]\label{gz:prop:Echi}
Each $E_\chi$ is a surjective homomorphism $(\SC,+)\to(\SC^\times,\cdot)$
with $|E_\chi(z)|=\exp_{\No}(\Re z)$ that commutes with conjugation. It
agrees with $\Efp$ at finite imaginary part, and $E_\chi(z+h)=E_\chi(z)\Ex(h)$
for $h\in\mm$. So $E_\chi$ is locally analytic with $E_\chi'=E_\chi$, and
$\Sin_\chi$, $\Cos_\chi$ satisfy the addition identities and have the
usual derivatives. Conversely, every homomorphism $E$ that agrees with
$\Efp$ at finite imaginary arguments and satisfies
$|E(z)|=\exp_{\No}(\Re z)$ is some $E_\chi$.
\end{proposition}
\begin{proof}
The group law and modulus are immediate. Conjugation follows from
$\Psi(-y)=\Psi(y)^{-1}=\overline{\Psi(y)}$. For surjectivity write
$w=|w|\cis\theta$ by \cref{gz:prop:polar}. An infinitesimal increment lies
in the finite summand, which gives the local law. For the converse, the
modulus law makes $y\mapsto E(iy)$ a phase
(\rl{trigonometry:per:prop:exponentials}).
\end{proof}

\subsection{The canonical exponential and the explicit families}\label{gz:sec:canonical}
\begin{definition}[The canonical exponential]\label{gz:def:Exp}
The character $\chi=1$ gives $\Phi(y)=\cis(\fin y)$ and
\[
\Exp(x+iy)=\exp_{\No}(x)\cis(\fin y),\qquad \ker\Exp=2\pi i\Oz .
\]
This is the canonical surcomplex exponential of Ehrlich and Kaplan
\cite[Theorem 11.1]{EK}, obtained by reducing angles modulo the initial
integer part $\Oz$. In the collection it is \rl{trigonometry:def:globaltrig},
\rl{trigonometry:thm:globalexp}, and in the analysis report
\rl{e:def-canonicalexp}, \rl{e:thm-exp}. It satisfies $\Exp(iY)=1$ for
every $Y\in\Pi$, and $\{z:\Sin z=0\}=\pi\Oz$ (\rl{trigonometry:cor:globalzeros}).
\end{definition}
06 and 03 use $\Exp$ and credit Ehrlich--Kaplan. 01's $E_0$ and 02's
$E_{\ell=0}$ are also $\Exp$, although neither source says so (errors 1 and
14). 04 mentions Ehrlich--Kaplan without using it, and 05 uses only $\Efp$
(error 38). With $\Exp$ fixed, $\Gamma_{\Exp}=\Exp\circ\LSt$ is the
canonical-exponential Gamma on all of $\Sinf$ (\cref{gz:sec:gammachi}). It
is canonical in the integer-part sense, not by uniqueness
(\cref{gz:sec:expconv}).

\begin{proposition}[Explicit phase families \src{01, 02, 03}]\label{gz:prop:families}
For a positive infinite monomial $M$ let $c_M(p)\in\R$ be the coefficient
of $M$ in $p\in\Pi$; for example $c_\omega$ is the coefficient of $\omega$.
\begin{enumerate}[label=(\alph*)]
\item \src{01} For $\lambda\in\R$, the character $\chi_\lambda(p)=e^{i\lambda c_\omega(p)}$
 gives $E_\lambda(a+ib)=\exp_{\No}(a)\cis(P_\lambda(b))$ with
 $P_\lambda(b)=\fin b+\lambda c_\omega(b)$. Then $E_0=\Exp$ and
 $E_\lambda(i\omega)=e^{i\lambda}$.
\item \src{03} For fixed $M$, $\chi_{\lambda,M}(p)=e^{i\lambda c_M(p)}$ gives
 $E_{\lambda,M}$. Its normalized period group is
 $\cI=\{x:\fin x+\lambda c_M(x)\in\Z\}$.
\item \src{02} For a real-linear $\theta:\Pi\to\R$, $\chi_\theta=\cis\circ\theta$
 gives $E_\theta(a+ib)=\exp_{\No}(a)\cis(\fin b+\theta(\pp b))$. Coefficient
 functionals $\theta=r\,c_M$ ($r\in\R$) are explicit and need no choice.
\end{enumerate}
In each case the character takes values in the ordinary circle $S^1(\C)$.
\end{proposition}
\begin{proof}
$c_M$ and $\fin$ are additive and real-linear, so each formula defines a
character. The value at $i\omega$ follows from $\fin\omega=0$ and
$c_\omega(\omega)=1$. For (b), $\Psi(2\pi x)=\cis(2\pi(\fin x+\lambda c_M(x)))$.
\end{proof}
(a) is literally the analysis report's twisted family \rl{e:thm-twisted},
with the same symbol and formula, and the trigonometry report's family at
$\omega$ (\rl{trigonometry:ex:phases}). All three families are subfamilies
of \rl{trigonometry:thm:characters}.

\begin{proposition}[Prescribed prime phases on a set-sized subspace \src{04}]\label{gz:prop:primephases}
Let $T$ be positive infinite, and let $V\subset\No$ be a set-sized
$\Q$-subspace containing $\R$ and every $T\log p$ for ordinary primes $p$.
Let $V_0$ be its subspace of finite elements. For arbitrary ordinary angles
$\vartheta_p$ there is a homomorphism $\chi:V\to\TT(\No)$ with
$\chi|_{V_0}=\cis$, $\chi(T\log p)=e^{i\vartheta_p}$, and
$\chi(y+h)=\chi(y)\cis h$ for $h\in V_0$.
\end{proposition}
\begin{proof}
The classes of the $T\log p$ in $V/V_0$ are $\Q$-independent: a nonzero
rational combination of the $\log p$ is a nonzero real by unique
factorization, and multiplying it by $T$ gives an infinite element. Define
a $\Q$-linear map on $V_0\oplus\operatorname{span}_\Q\{T\log p\}$ fixing
$V_0$ and sending $T\log p$ to $\vartheta_p$. Extend it by a basis of the
quotient (choice) to $P:V\to V_0$, and set $\chi=\cis\circ P$.
\end{proof}
These are local-data extensions. 04 does not claim they satisfy
Ehrlich--Kaplan's initial-integer-part characterization, and they do not
contradict it. The proposition shows that finite-angle data and finite
translation laws leave the phases at infinite arguments undetermined.

\subsection{Period groups and forced infinite periods}\label{gz:sec:periods}
\begin{theorem}[Period groups \src{03}]\label{gz:thm:periods}
For a phase $\Psi=\Psi_\chi$ let $\cI(\Psi)=\{x\in\No:\Psi(2\pi x)=1\}$.
\begin{enumerate}[label=(\alph*)]
\item $\cI(\Psi)$ is an additive subgroup with $\cI(\Psi)\cap\OO_\R=\Z$.
 Every coset $x+\OO_\R$ meets it in exactly one coset of $\Z$. In
 particular it has positive and negative infinite elements.
\item $\Sin_\chi(\pi z)=0\iff z\in\cI(\Psi)$, and
 $\Sin_\chi(\pi z/2)=0\iff z\in2\cI(\Psi)$. These zeros are real and
 simple, and $\Cos_\chi(\pi x)=\pm1$ for $x\in\cI(\Psi)$.
\item $\cI(\Phi)=\Oz$, and for 01's family, $\cI(\Psi_{\chi_\lambda})=\{x:P_\lambda(x)\in\Z\}$.
\end{enumerate}
\end{theorem}
\begin{proof}
(a) The finite part is the kernel statement of \cref{gz:prop:polar}. For
$x\in\No$ choose a finite $\theta$ with $\cis\theta=\Psi(2\pi x)$; then
$x-\theta/2\pi\in\cI(\Psi)$ differs from $x$ by a finite amount. Two
elements of $\cI(\Psi)$ in one finite coset differ by a finite period,
hence by an ordinary integer. (b) $\Sin_\chi(\pi z)=0$ means
$E_\chi(2\pi iz)=1$. For $z=x+iy$ the modulus is $\exp_{\No}(-2\pi y)$,
which equals $1$ only if $y=0$, and then $x\in\cI(\Psi)$. There
$E_\chi(i\pi x)^2=1$, so $\Cos_\chi(\pi x)=\pm1$ and the derivative
$\pi\Cos_\chi(\pi x)$ is nonzero. Replace $z$ by $z/2$ for the second
statement. (c) $\Phi(2\pi x)=\cis(2\pi\fin x)$, so $\cI(\Phi)=\{x:\fin x\in\Z\}=\Pi\oplus\Z=\Oz$;
the second claim is \cref{gz:prop:families}(b) with $M=\omega$.
\end{proof}
Part (a) is \rl{trigonometry:thm:infiniteperiods} in normalized form, and
$\cI(\Psi)$ is the normalized period group of
\rl{trigonometry:per:thm:dictionary}. (03 writes $Z_\phi$; the name
$Z_E$ in the trigonometry report denotes the kernel-multiplier ring, a
different object.) By \rl{trigonometry:per:thm:cofinal}, every phase on
$\No$ even has cofinally large periods all of whose rational multiples are
periods. (b) for $\Phi$ is \rl{trigonometry:cor:globalzeros}.

\subsection{Resonances}\label{gz:sec:resonance}
\begin{proposition}[Two-frequency test \src{06}]\label{gz:prop:twofreq}
The statement
$\cis(t\log2)=1\wedge\cis(t\log3)=1\Rightarrow t=0$ holds for ordinary real
$t$. It remains true in every elementary extension of the real field with
its cosine and sine. It fails for the phase $\Phi$ of $\Exp$: every nonzero
$T\in\Pi$ satisfies both antecedents.
\end{proposition}
\begin{proof}
Over $\R$, $t\log2=2\pi k$ and $t\log3=2\pi l$ give $3^k=2^l$, so $k=l=0$ by
unique factorization. The implication is a first-order sentence in the
named functions, so it transfers. For $\Phi$, $\Phi(T\log n)=\cis(\fin(T\log n))=1$.
\end{proof}
So $\Exp$ cannot be the full complex exponential of an elementary
extension. The structure behind this is the following notion.

\begin{definition}[Resonance]\label{gz:def:resonance}
A nonzero $T\in\No$ is a \emph{resonance} of the phase $\Psi$ if
$\Psi(T\log n)=1$ for every ordinary integer $n\ge1$.
\end{definition}
Every nonzero $T\in\Pi$ is a resonance of $\Phi$. A resonance is a nonzero
element of $\bigcap_{n\ge2}(2\pi/\log n)\cdot\cI(\Psi)$, and it is much
rarer than a period. The following three facts are established in this
merge. They settle the last sentence of 06's Remark 7.4 (``The existence of
an infinite period of $U$ alone does \emph{not} imply such a simultaneous
logarithmic resonance'', where 06's $U$ is $\Psi_\chi$), which 06 stated
without proof (error 39).

\begin{proposition}[Existence and failure of resonances \merge]\label{gz:prop:resonance}
\leavevmode
\begin{enumerate}[label=(\alph*)]
\item If $\chi(\Pi)\subseteq S^1(\C)$, then $\Psi_\chi$ has resonances in
 $\Pi$ of arbitrarily large modulus. This covers $\Exp$, every $E_\lambda$,
 every $E_{\lambda,M}$ and every $E_\theta$.
\item If $\Psi(p+f)=\cis(\varphi(p)+f)$ with $\varphi:\Pi\to\OO_\R$
 real-linear, then $T=p-\varphi(p)$ is a resonance for every nonzero
 $p\in\Pi$. Its galaxy is $iT+\OO=ip+\OO$.
\item Assume choice for a basis of $\R$ over $\Q(\log3/\log2)$. There is a
 phase $\Psi$ with no nonzero $T\in\No$ such that $\Psi(T\log2)=\Psi(T\log3)=1$.
 In particular $\Psi$ has no resonance. It has infinite periods, as every
 phase does.
\end{enumerate}
\end{proposition}
\begin{proof}
(a) Let $B\in\No$. Let $\mathfrak h$ be the Hartogs number of the set
$S^1(\C)^{\mathbb P}$, where $\mathbb P$ is the set of ordinary primes;
$\mathfrak h$ is an ordinal that does not inject into that set. Choose a
positive surreal $g$ with $\omega^g>\omega(|B|+1)$, and let $G\subset\Pi$
be the real span (finite combinations) of the monomials
$\omega^{g+\beta}$, $\beta<\mathfrak h$. A nonzero element of $G$ has a leading
term $r\omega^{g+\beta}$ with $r\in\R^\times$ and $\beta\ge0$, so its
modulus exceeds $\tfrac12|r|\omega^g>\tfrac12|r|\omega(|B|+1)>|B|$. The map $G\to S^1(\C)^{\mathbb P}$,
$P\mapsto(\chi(P\log p))_p$, is a group homomorphism, because $\chi$ is a
character and $P\mapsto P\log p$ is additive into $\Pi$. It cannot be
injective, since $\beta\mapsto\omega^{g+\beta}$ injects $\mathfrak h$ into $G$.
So there is a nonzero $P\in G$ with $\chi(P\log p)=1$ for every prime $p$.
For every $n=\prod p^{v_p}$, $\chi(P\log n)=\prod_p\chi(P\log p)^{v_p}=1$.
Since $P\log n\in\Pi$, $\Psi_\chi(P\log n)=\chi(P\log n)=1$, and $\pm P$ is
a positive resonance with $|P|>|B|$. The families of
\cref{gz:prop:families} take values in $S^1(\C)$.

(b) For every $n$, $T\log n=p\log n-\varphi(p)\log n$ with $p\log n\in\Pi$ and
$\varphi(p)\log n$ finite. By real-linearity,
$\Psi(T\log n)=\cis(\varphi(p\log n)-\varphi(p)\log n)=\cis(0)=1$. Also
$T-p=-\varphi(p)$ is finite.

(c) Let $c=\log3/\log2$. It is irrational by unique factorization, and
transcendental by the Gelfond--Schneider theorem, since $2^c=3$. Hence
$\iota_c:\Q(c)\to\Q(c)$, $g(c)\mapsto g(c^2)$, is an injective field
homomorphism. Choose a basis $(e_i)$ of $\R$ over $\Q(c)$ and define
$\mathfrak s'(\sum_ig_i(c)e_i)=\sum_i\iota_c(g_i(c))e_i$. Then $\mathfrak s'$ is
$\Q$-linear and injective, and $\mathfrak s'(cr)=c^2\mathfrak s'(r)$. Hence
$\mathfrak s'(cr)-c\mathfrak s'(r)=(c^2-c)\mathfrak s'(r)\ne0$ for $r\ne0$. Put
$\mathfrak s(x)=\log2\cdot\mathfrak s'(x/\log2)$. For $P=\sum_yr_y\omega^y\in\Pi$
(all $y>0$) define
\[
\nu(P)=\sum_y\mathfrak s(r_y)\,\omega^{-1/y}\in\mm_\R .
\]
Since $y\mapsto-1/y$ is increasing on positive surreals, the exponents form
a reverse well-ordered set of negative surreals, so this is a normal form.
The map $\nu$ is additive because $\mathfrak s$ is. Hence
$\Psi(P+f)=\cis(\nu(P)+f)$ is a phase, and it has infinite periods by
\cref{gz:thm:periods}.

Suppose $T=P+f\ne0$ ($P\in\Pi$, $f\in\OO_\R$) satisfies
$\Psi(T\log2)=\Psi(T\log3)=1$. Then $\nu(P\log2)+f\log2\in2\pi\Z$ and
$\nu(P\log3)+f\log3\in2\pi\Z$. Taking standard parts,
$\st(f)\log2=2\pi k$ and $\st(f)\log3=2\pi l$. If $\st f\ne0$, then
$\log2/\log3=k/l$, which is impossible. So $k=l=0$ and
$f=-\nu(P\log2)/\log2=-\nu(P\log3)/\log3$. Comparing coefficients at each
$\omega^{-1/y}$ gives $\mathfrak s(r_y\log2)/\log2=\mathfrak s(r_y\log3)/\log3$, that
is, $\mathfrak s'(r_y)=\mathfrak s'(cr_y)/c$. By the displayed property of $\mathfrak s'$
this forces $r_y=0$ for every $y$. So $P=0$, then $f=0$ and $T=0$, a
contradiction.
\end{proof}
Part (a) is the pigeonhole argument of \rl{trigonometry:per:thm:cofinal},
applied to the prime values of a character. Part (b) was observed in the
review of 06, and it covers every family of \cref{gz:prop:families}; for
$E_\lambda$ it gives $T=\omega-\lambda$. Part (c) needs characters whose
values are $\cis\delta$ with nonzero infinitesimal $\delta$, values in
$1+\mm$ off $S^1(\C)$. By (a) no ordinary-circle phase can serve.
Consequence: 06's resonance theorem applies verbatim to every phase with a
resonance (\cref{gz:thm:resonance}), as 06's Remark 7.4 anticipates. This
includes every explicitly specified phase in the sources and the
collection's explicit ordinary-circle families ($\Exp$, $E_\lambda$,
$E_{\lambda,M}$, $E_\theta$). 03's statements for arbitrary phases, 04's
prescribed prime phases (\cref{gz:prop:primephases}, partial characters on
a set-sized subspace), and the collection's exponentials whose characters
take values in $1+\mm$ (\rl{trigonometry:per:thm:robustness}) are covered
only when the phase in question has a resonance. The theorem does not apply
to the phase in (c), and whether a coherent galaxy zeta exists for that
phase is open (\cref{gz:q:nonresonant}).

\part{The finite plane}\label{gz:part:finite}

\section{Canonical finite lifts and divisor rigidity}\label{gz:sec:finite}

\subsection{The lift}
\begin{definition}[Canonical finite lift \src{01--06}]\label{gz:def:lift}
Let $f$ be an ordinary meromorphic function near $c\in\C$, with Laurent
germ $f(c+u)=\sum_{n\ge-m}a_nu^n$. For $\eps\in\mm$, $\eps\ne0$, put
\[
f^\#(c+\eps)=\sum_{n=-m}^{-1}a_n\eps^n+\Sh_{n\ge0}a_n\eps^n .
\]
At $\eps=0$ use the ordinary value if $f$ is regular at $c$. At a pole no
field value is assigned and the meromorphic germ is retained. Removable
singularities are removed in the ordinary germ before lifting.
\end{definition}
\Cref{gz:lem:substitution} makes this well defined, and the centre
$c=\st z$ is unique. A pole does not delete its monad: every nonzero
infinitesimal displacement has a value, always of infinite modulus, since
it is $g(0)\eps^{m}(1+\eta)$ with $m\le-1$ (\cref{gz:thm:divisor}; 05 and
06 say ``usually'', error 53). This is the analysis report's canonical
lift (\rl{c:p4:def-lift}), extended here to meromorphic germs. 03 writes
$f^{\mathrm H}$, and 01 and 05 write $f^\sharp$.

\begin{proposition}[Functoriality \src{01--06}]\label{gz:prop:functorial}
Lifting respects sums, products, reciprocals where nonzero, analytic
composition where the ordinary germs compose (expanded around the ordinary
value of the inner germ), and coordinate differentiation. Every identity of
ordinary meromorphic germs becomes an identity of lifted germs on every
common finite domain. For conjugation, put $\bar f(z)=\overline{f(\bar z)}$.
Then $(\bar f)^\#(\bar s)=\overline{f^\#(s)}$, and so
$f^\#(\bar s)=\overline{f^\#(s)}$ when $f$ is real-symmetric
($\bar f=f$), as $\Gamma$, $\zeta$ and $\xi$ are.
\end{proposition}
\begin{proof}
The identities are identities of Taylor--Laurent series. After finitely many
negative powers are factored out, all operations are covered by
\cref{gz:lem:substitution}. A unit has nonzero constant term, so its inverse
is a geometric series in an infinitesimal. Composition is expanded after
subtracting the ordinary value of the inner germ. Differentiation is
coefficientwise, and it is the fine derivative by \rl{b:germalgebra}. For
conjugation, the Laurent coefficients of $\bar f$ at $\bar c$ are
$\overline{a_n}$, and conjugation acts coefficientwise on normal forms.
\end{proof}
The conjugation clause corrects 05's Proposition 3.1, which states it
without the bar on $f$ (error 35).

\begin{theorem}[Divisor rigidity \src{01--06}]\label{gz:thm:divisor}
Let $f$ be a nonzero ordinary meromorphic function on a connected ordinary
domain, and write $f(c+u)=u^mg(u)$ with $m\in\Z$ and $g(0)\ne0$. For every
nonzero $\eps\in\mm$,
\[
f^\#(c+\eps)=g(0)\,\eps^m(1+\eta),\qquad\eta\in\mm,\qquad
v\bigl(f^\#(c+\eps)\bigr)=m\,v(\eps).
\]
In particular $f^\#(c+\eps)\ne0$. The zeros and poles of $f^\#$ on the
finite plane are exactly the embedded ordinary zeros and poles of $f$, with
the same multiplicities and orders.
\end{theorem}
\begin{proof}
$g(\eps)/g(0)=1+\eta$ with $\eta\in\mm$ by \cref{gz:lem:substitution}, and
this is a unit of valuation zero. Multiply by $g(0)\eps^m$. At $\eps=0$ the
assertion is the ordinary one. Every finite point has a unique standard
part, so these local statements exhaust $\OO$.
\end{proof}
All six sources prove this (01 Thm 3.2; 02 Thm 3.1, Cor 3.2; 03 Thm 3.2;
04 Thm 3.1, Cor 3.2; 05 Thm 4.1; 06 Thm 3.3). In the holomorphic case it
is also the analysis report's \rl{c:p4:liftzeros} with
\rl{c:p4:rk-liftzeros}. The theorem is
stronger than preservation of standard parts. A fixed classical function
acquires no cloud of displaced zeros by adjoining infinitesimals, not even
at a multiple zero. It concerns the \emph{unchanged} classical germ: a
Hahn-coefficient deformation such as $f^\#+\eta h^\#$ can move zeros within
their monads (\cref{gz:sec:robust}; analysis \rl{d:thm:zerocluster}).

\begin{remark}[The word canonical \src{01, 02, 04, 06}]\label{gz:rem:canonical}
The lift is canonical as the evaluation of a specified germ. It is not a
uniqueness theorem among all finely differentiable functions that agree
with $f$ on $\C$, and agreement at ordinary points does not imply the germ
law.
\end{remark}

The classical facts used below are in \cite{DLMFgamma,DLMFzeta}.

\begin{corollary}[Finite Gamma, zeta and $\xi$ \src{01--06}]\label{gz:cor:finite-divisors}
$\Gp$ has no zeros. Its poles are exactly the ordinary integers
$0,-1,-2,\dots$, simple, with residues $(-1)^m/m!$, and $(1/\Gamma)^\#$ has
zeros exactly there. $\zp$ has one simple pole, at $1$, with residue $1$.
Its zeros are exactly the ordinary zeros of $\zeta$ with their
multiplicities, and the trivial zeros $-2,-4,\dots$ are simple. $\xip$,
the lift of the entire function $\xi$, has exactly the nontrivial zeros
with their multiplicities.
\end{corollary}

\subsection{Explicit germs}\label{gz:sec:germs}
For nonzero $\eps\in\mm$, with Euler's constant $\gamma$, the Stieltjes
constants $\gamma_n$ ($\gamma_0=\gamma$) and $H_m=\sum_{j\le m}1/j$
($H_0=0$) \src{01--06}:
\begin{align}
\Gp(\eps)&=\eps^{-1}-\gamma+\Bigl(\frac{\gamma^2}2+\frac{\pi^2}{12}\Bigr)\eps+\bigO(\eps^2),\label{gz:eq:gammazero}\\
\Gp(-m+\eps)&=\frac{(-1)^m}{m!\,\eps}\bigl(1+(H_m-\gamma)\eps+\bigO(\eps^2)\bigr),\label{gz:eq:gammapole}\\
\zp(1+\eps)&=\eps^{-1}+\Sh_{n\ge0}\frac{(-1)^n\gamma_n}{n!}\eps^n,\label{gz:eq:zetapole}\\
\Lgm\Gp(1+\eps)&=-\gamma\eps+\Sh_{n\ge2}\frac{(-1)^n\zeta(n)}n\eps^n .\label{gz:eq:loggamma}
\end{align}
Here $\bigO(\eps^k)$ means $\eps^k$ times a finite strong evaluation of the
indicated germ. It is not a limit along a sequence. The next coefficient in
\eqref{gz:eq:gammazero} is $-(\gamma^3/6+\gamma\pi^2/12+\zeta(3)/3)\approx-0.9075$.
At a zero $\rho$ of $\zeta$ of order $m$,
$\zp(\rho+\eps)=\frac{\zeta^{(m)}(\rho)}{m!}\eps^m(1+\eta)$. At the trivial
zeros \src{01, 02, 04},
\begin{equation}\label{gz:eq:trivialslope}
\zeta'(-2r)=\frac{(-1)^r(2r)!\,\zeta(2r+1)}{2(2\pi)^{2r}}\ne0,\qquad
\zp(-2r+\eps)=\zeta'(-2r)\,\eps\,(1+\eta),
\end{equation}
so $\zeta'(-2)=-\zeta(3)/(4\pi^2)$. For a simple zero $\rho$ the formal
Newton map satisfies
$N(\rho+\eps)-\rho=\frac{\zeta''(\rho)}{2\zeta'(\rho)}\eps^2+\bigO(\eps^3)$
\src{05}: the valuation at least doubles. At $\eps=\omega^{-1}$
\src{02, 04, 05},
\[
\Gp(\omega^{-1})=\omega-\gamma+\Bigl(\tfrac{\gamma^2}2+\tfrac{\pi^2}{12}\Bigr)\omega^{-1}+\bigO(\omega^{-2}),\qquad
\zp(1+\omega^{-1})=\omega+\gamma_0-\gamma_1\omega^{-1}+\tfrac{\gamma_2}2\omega^{-2}-\cdots,
\]
and at a simple zero, $\zp(\rho+\omega^{-1})\sim\zeta'(\rho)\omega^{-1}$.
At a hypothetical double zero the scale would be $\omega^{-2}$. Small
residuals measure multiplicity. They are not further exact zeros.

\subsection{Identities}\label{gz:sec:finite-identities}
On finite arguments, as identities of lifted meromorphic germs
\src{01--06}:
\begin{align}
\Gp(z+1)&=z\Gp(z),\qquad \Gp(z)\Gp(1-z)=\frac\pi{\sin^\#(\pi z)},\label{gz:eq:finiteref}\\
\prod_{j=0}^{q-1}\Gp\Bigl(z+\frac jq\Bigr)&=(2\pi)^{(q-1)/2}q^{1/2-qz}\Gp(qz)\qquad(q\in\N_{\ge1}),\label{gz:eq:finitegauss}\\
\zp(s)&=2^s\pi^{s-1}\sin^\#(\pi s/2)\Gp(1-s)\zp(1-s),\label{gz:eq:finiteFE}\\
\xip(s)&=\xip(1-s),\qquad\xip(\bar s)=\overline{\xip(s)},\qquad\Xi^\#(-w)=\Xi^\#(w),\label{gz:eq:xisym}
\end{align}
where $\Xi(w)=\xi(\tfrac12+iw)$ is real entire. Powers of ordinary positive
constants use $\Efp$, and the index $q$ is an ordinary integer. At points
where a pole meets a zero, the germs are multiplied first. $0\cdot\Gp(0)$
is not a product of field values. No product with a surreal number of
factors and no recurrence iterated a surreal number of times is used in the
surreal constructions. Internal products occur only inside transported
models (\cref{gz:sec:transported-divisors}).

\subsection{Coefficientwise formulas and the failure of the raw series}\label{gz:sec:coefficientwise}
For $a\in\C$ with $\Re a>1$ and $\eps\in\mm$ \src{02, 04, 06},
\begin{equation}\label{gz:eq:coefdirichlet}
\zp(a+\eps)=\Sh_{k\ge0}\frac{(-\eps)^k}{k!}\Bigl(\sum_{n\ge1}(\log n)^kn^{-a}\Bigr),
\end{equation}
and for $\Re a>0$ \src{02, 04},
\begin{equation}\label{gz:eq:coefmellin}
\Gp(a+\eps)=\Sh_{j\ge0}\frac{\eps^j}{j!}\int_0^\infty e^{-t}t^{a-1}(\log t)^j\,dt .
\end{equation}
The inner sums and integrals are ordinary, absolutely convergent, and
computed in $\C$; only the outer sums are strong. This is summation notion
(c) of \cref{gz:sec:sumnotions}. Likewise $\Gamma(c)\zeta(c)=\int_0^\infty x^{c-1}(e^x-1)^{-1}dx$
gives the coefficients of $(\Gamma\zeta)^\#$ at $\Re c>1$ \src{04}.

\begin{warning}[The raw finite Dirichlet family \src{01, 02, 04, 05, 06}]\label{gz:warn:finite-raw}
At a finite $s=a+\eps$, the family $(n^{-s})_{n\ge1}$ is not strongly
summable, even when $\Re a>1$: every member has the nonzero standard part
$n^{-a}$ at exponent zero. The double family behind \eqref{gz:eq:coefdirichlet}
is not strongly summable either. Raw Euler, Weierstrass and Hadamard
products are equally inadmissible at finite arguments.
\end{warning}

\begin{corollary}[Zero-free halos \src{06}]\label{gz:cor:halofree}
If an ordinary set $U\subseteq\C$ contains no zero of $f$, then $f^\#$ has
no zero on $U^\#$. In particular the zero-free half-plane $\Re s\ge1$, the
localization of nontrivial zeros in $0<\Re s<1$, and every numerically
verified zero-free region hold exactly for $\zp$, with no infinitesimal
boundary layer of new zeros.
\end{corollary}
\begin{proof}
A zero of $f^\#$ is an ordinary zero by \cref{gz:thm:divisor}, and its
standard part is itself.
\end{proof}

\section{The finite Riemann hypothesis and its robust form}\label{gz:sec:robust}

\subsection{Finite RH}
\begin{definition}[\src{04, 05, 06}]\label{gz:def:RHfin}
$\RH_{\mathrm{fin}}$: every $s\in\OO$ with $\xip(s)=0$ has $\Re s=\tfrac12$.
Its \emph{shadow} form \src{04} replaces the conclusion by $\Re\st s=\tfrac12$.
\end{definition}

\begin{theorem}[Finite RH \src{01, 02, 04, 05, 06}]\label{gz:thm:finiteRH}
The following are equivalent: classical RH; $\RH_{\mathrm{fin}}$; its shadow
form; ``every zero of $\zp$ with $0<\Re s<1$ has $\Re s=\tfrac12$'' (01); and
``all zeros of $\xip$ have real part exactly $\tfrac12$'' (02).
\end{theorem}
\begin{proof}
By \cref{gz:thm:divisor} the zeros of $\xip$ and the zeros of $\zp$ in the
finite strip are exactly the classical nontrivial zeros, at the embedded
ordinary points. For an ordinary point, $\Re s=\Re\st s$. Conversely,
restrict each statement to $\C$.
\end{proof}
The equivalence neither strengthens nor weakens RH. What it rules out is
spurious ``infinitesimal counterexamples'': $\xip(\rho+\eps)\sim\xi'(\rho)\eps\ne0$
at a simple zero, and an infinitesimal residual is not a zero. 03's
horizontal form (\cref{gz:cor:horizontalRH}) is the same statement,
because $\xi_\chi$ has no zeros on $\Tp\cup\Tm$. The transferred form is
\cref{gz:thm:RHstar}.

\subsection{Local lemmas}
\begin{lemma}[Formal implicit function \src{05}]\label{gz:lem:implicit}
Let $A(X,u)\in\C[[X,u]]$ satisfy $A(0,0)=0$ and $\partial_XA(0,0)\ne0$.
There is a unique $d(u)\in u\C[[u]]$ with $A(d(u),u)=0$, and
$A=(X-d(u))U(X,u)$ with $U(0,0)=\partial_XA(0,0)$. For every infinitesimal
$\tau$, the only infinitesimal $\eps$ with $A(\eps,\tau)=0$ is $d(\tau)$,
including solutions outside the subfield generated by $\tau$.
\end{lemma}
\begin{proof}
Solve for the coefficients of $d$ recursively; at each stage the new
coefficient is multiplied by $\partial_XA(0,0)\ne0$. Formal division by
$X-d(u)$, using $X^k-d^k=(X-d)\sum_jX^{k-1-j}d^j$, gives $U$. After strong
evaluation $U(\eps,\tau)$ is a unit.
\end{proof}

\begin{lemma}[Constant perturbation of a zero \src{04}]\label{gz:lem:splitting}
Let $f$ be holomorphic near $c$ with $f(c+u)=u^ma(u)$, $m\ge1$,
$a(0)=a_0\ne0$. Let $\eta>0$ be infinitesimal and $\delta\in\{\pm1\}$. Then
$f^\#+\delta\eta$ has exactly $m$ zeros in $c+\mm$, all simple:
$z_j=c+rv_j(r)$ with $r=\eta^{1/m}$, where the $v_j$ are convergent ordinary
germs with $a_0v_j(0)^m=-\delta$ and distinct nonzero $v_j(0)$.
\end{lemma}
\begin{proof}
$F(r,v)=v^ma(rv)+\delta$ has $m$ simple nonzero roots at $r=0$. The
ordinary implicit-function theorem gives the germs $v_j$, and evaluation at
$r$ gives zeros. Conversely, let $z=c+u$ be a zero in the monad. Then
$u\ne0$ and $m\,v(u)=v(\eta)$ by \cref{gz:thm:divisor}, so $v=u/r$ has
valuation zero. Reducing $F(r,v)=0$ modulo $\mm$ shows that $\st v$ is one
of the $v_j(0)$. The local factorization $F=(v-v_j(r))U_j$ with
$U_j(0,v_j(0))\ne0$ becomes a unit after evaluation, so $v=v_j(r)$. Finally
$f'(c+u)=u^{m-1}(ma_0+\text{infinitesimal})\ne0$.
\end{proof}

\begin{proposition}[Local multiplicity conservation \src{04}]\label{gz:prop:conservation}
Let $f\not\equiv0$ and $h$ be ordinary entire functions, $\eta$ a nonzero
infinitesimal, and $F_\eta=f^\#+\eta h^\#$. In the monad of a zero $c$ of
$f$ of order $m$, $F_\eta$ has exactly $m$ zeros counted with multiplicity.
No finite zero lies outside the monads of zeros of $f$. If $f$ and $h$ are
real entire, $\eta$ is real and $c$ is a real simple zero of $f$, the zero
of $F_\eta$ in $c+\mm$ is unique, real and simple, and
$c_\eta=c-\eta\,h(c)/f'(c)+\bigO(\eta^2)$ \src{04}.
\end{proposition}
\begin{proof}
The last claim follows by taking standard parts. If $h\equiv0$ near $c$
there is nothing to prove. Otherwise write $h(c+u)=u^kb(u)$ with $b(0)\ne0$.
If $k\ge m$, then $F_\eta=u^m[a(u)+\eta u^{k-m}b(u)]$ and the bracket is a
unit on the monad. If $k<m$, then $F_\eta=u^k\{u^{m-k}a(u)+\eta b(u)\}$.
Put $r=\eta^{1/(m-k)}$ (a root in $\SC$) and $u=rv$. The argument of
\cref{gz:lem:splitting} gives exactly $m-k$ simple nonzero roots of the
bracket, so the total is $m$. For the last clause, apply
\cref{gz:lem:implicit} to $A(X,u)=f(c+X)+u\,h(c+X)$, with
$\partial_XA(0,0)=f'(c)\ne0$: the unique zero in the monad is $c+d(\eta)$
with $d'(0)=-h(c)/f'(c)$, and the evaluated factorization $A=(X-d(u))U$
with $U$ a unit shows that it is simple. Since $f,h,c,\eta$ are real,
conjugation maps this zero to a zero of $F_\eta$ in the same monad, so
uniqueness makes it real.
\end{proof}

\subsection{Protection of a simple critical-line zero}
\begin{definition}[J-symmetric formal deformations]\label{gz:def:DJ}
Let $J(s)=1-\bar s$. The class $\DJs$ consists of formal deformations
$F(s,u)=\xi(s)+\sum_{n\ge1}f_n(s)u^n$ with ordinary entire $f_n$ satisfying
$f_n(1-\bar s)=\overline{f_n(s)}$. No convergence in $u$ is required. For
real $\tau\in\mm$ and $s=a+\eps\in\OO$ put
\[
F_\tau(s)=\xip(s)+\Sh_{n\ge1,\,k\ge0}\frac{f_n^{(k)}(a)}{k!}\tau^n\eps^k ,
\]
a strongly summable double family by \cref{gz:lem:substitution}.
\end{definition}
The three source classes are subclasses of $\DJs$:
\begin{itemize}
\item 04: $f^\#+\eta h^\#$ with $h$ real entire in the variable $w$ of
 $s=\tfrac12+iw$. In $s$ this is exactly $h(1-\bar s)=\overline{h(s)}$.
\item 05: formal series whose coefficients satisfy both $f_n(1-s)=f_n(s)$
 and $f_n(\bar s)=\overline{f_n(s)}$, which together imply J-symmetry,
 with $\tau>0$.
\item 06: $\xip+\eps h^\#$ with $h\in\Ssym$, the real vector space of entire
 $h$ with both symmetries, with $\eps$ real of either sign.
\end{itemize}
Every $F_\tau$ has standard part $\xi(\st s)$, so its finite zeros lie in
monads of zeros of $\xi$. For real $\tau$,
$F_\tau(J(s))=\overline{F_\tau(s)}$. To see this at a point $\rho$ with
$\Re\rho=\tfrac12$, write $f_n(\rho+X)=\sum_kc_{n,k}X^k$. J-symmetry gives
$\overline{c_{n,k}}=(-1)^kc_{n,k}$, and $J(\rho+\eps)=\rho-\bar\eps$.

\begin{theorem}[Protection of simple critical-line zeros \src{04, 05, 06}]\label{gz:thm:protection}
Let $\rho$ be a simple zero of $\xi$ with $\Re\rho=\tfrac12$, let
$F\in\DJs$, and let $\tau$ be a real infinitesimal. Then $F_\tau$ has exactly
one zero $\rho_\tau$ in $\rho+\mm$, and $\Re\rho_\tau=\tfrac12$ exactly.
Moreover $\rho_\tau=d(\tau)+\rho$ for a formal series $d$ with purely
imaginary coefficients, and
\[
\rho_\tau=\rho-\tau\frac{f_1}{\xi'}+\tau^2\Bigl(\frac{f_1f_1'}{\xi'^2}-\frac{\xi''f_1^2}{2\xi'^3}-\frac{f_2}{\xi'}\Bigr)+\bigO(\tau^3),
\]
all functions evaluated at $\rho$. If $f_1(\rho)\ne0$ and $\tau\ne0$, the
displacement has valuation $v(\tau)$.
\end{theorem}
\begin{proof}
Apply \cref{gz:lem:implicit} to $A(X,u)=F(\rho+X,u)$, for which
$\partial_XA(0,0)=\xi'(\rho)\ne0$. This gives the unique zero
$\rho+d(\tau)$ in the monad. The function $d^*(u)=-\overline{d}(u)$, with
conjugated and negated coefficients, also solves $A(d^*(u),u)=0$, because
$\overline{c_{n,k}}=(-1)^kc_{n,k}$. By uniqueness $d^*=d$, so every
coefficient of $d$ is purely imaginary. Since $\tau$ is real, $d(\tau)$ is
purely imaginary, so $\Re\rho_\tau=\tfrac12$. Equivalently, $J$ maps the
zero to a zero in the same monad, and uniqueness forces $J(\rho_\tau)=\rho_\tau$.
The expansion comes from comparing the coefficients of $u$ and $u^2$ in
$A(d(u),u)=0$. The valuation claim holds because the first coefficient is
a nonzero ordinary number.
\end{proof}
06 gives the second-order formula for linear deformations ($f_2=0$) and
checks it symbolically. 04's Appendix A gives the same coefficients in the
$w$-plane, as $B=h_0h_1/a^2-bh_0^2/a^3$ with $b=f''/2$ (checked for this
merge, \cref{gz:app:checks}). 05 states that every displacement coefficient
is purely imaginary but gives no argument. The involution argument above
supplies the missing step (error 33). The hypotheses that $\tau$ is real and
the deformation J-symmetric are essential. Without them, even a simple
critical-line zero moves off the line at first order.

\subsection{One stability theorem}
\begin{theorem}[Two-sign reality criterion \src{04}]\label{gz:thm:twosign}
Let $f$ be a nonzero real entire function ($f(\bar w)=\overline{f(w)}$) and
fix one positive infinitesimal $\eta$. The following are equivalent: (i)
every ordinary zero of $f$ is real and simple; (ii) every finite zero of
$f^\#+\eta$ and of $f^\#-\eta$ is real.
\end{theorem}
\begin{proof}
Every finite zero of $f^\#\pm\eta$ lies in the monad of a zero of $f$.
Assume (i). In the monad of a real simple zero, \cref{gz:lem:splitting}
gives a unique zero for each sign. Conjugation fixes the equation and the
monad, so the zero is real. Conversely, a nonreal zero of $f$ gives
nonreal perturbed zeros. Let $c$ be a real zero of order $m\ge2$ with real
leading coefficient $a_0$. For $m=2$ choose $\delta$ with the sign of
$a_0$; then $a_0v^2=-\delta$ has nonreal roots. For $m\ge3$ the $m$ distinct
roots of $a_0v^m=-\delta$ are not all real. In both cases
\cref{gz:lem:splitting} gives a zero $c+rv_j(r)$ whose imaginary part has
the nonzero leading coefficient $\Im v_j(0)$ at scale $\eta^{1/m}$.
\end{proof}

\begin{theorem}[Unified stability theorem \src{04, 05, 06}; unified form \merge]\label{gz:thm:stability}
The following are equivalent.
\begin{enumerate}[label=(\roman*)]
\item Classical RH holds and every nontrivial zero of $\zeta$ is simple.
\item For every $F\in\DJs$ and every real infinitesimal $\tau$, every finite
 zero of $F_\tau$ lies on $\Re s=\tfrac12$.
\item For one positive infinitesimal $\eta$, every finite zero of $\xip+\eta$
 and of $\xip-\eta$ lies on $\Re s=\tfrac12$.
\end{enumerate}
If ``lies on $\Re s=\tfrac12$'' is replaced by ``has standard part on
$\Re s=\tfrac12$'' in (ii) or (iii), each becomes equivalent to RH alone.
\end{theorem}
\begin{proof}
(i)$\Rightarrow$(ii): a finite zero has as standard part a zero $\rho$ of
$\xi$, which is simple and on the line by (i). Apply
\cref{gz:thm:protection}. (ii)$\Rightarrow$(iii): take $F=\xi\pm u$, which
lie in $\DJs$, and $\tau=\eta$. (iii)$\Rightarrow$(i): with $s=\tfrac12+iw$,
$\xip(\tfrac12+iw)=\Xi^\#(w)$ by functoriality, and $\Re s=\tfrac12$ means
$w$ is real. Apply \cref{gz:thm:twosign} to the real entire $\Xi$: all
zeros of $\Xi$ are real and simple, which is (i). For the shadow forms,
every finite zero of $F_\tau$ has as standard part a zero of $\xi$. Every
zero of $\xi$ is the standard part of a zero of $\xip\pm\eta$
(\cref{gz:lem:splitting}), and of $F_\tau$ with $F=\xi$.
\end{proof}

\begin{corollary}[The source criteria \src{04, 05, 06}]\label{gz:cor:criteria}
Each of the following is equivalent to RH plus simplicity.
\begin{enumerate}[label=(\alph*)]
\item \src{04} For one fixed positive infinitesimal $\eta$, all finite zeros
 of $\Xi^\#\pm\eta$ are real. Equivalently, for every ordinary real entire
 $h$ (in $w$), all finite zeros of $\Xi^\#+\eta h^\#$ are real; testing
 $h=\pm1$ suffices.
\item \src{05} For every formal deformation with coefficients having both
 symmetries and every positive infinitesimal $\tau$, all finite zeros lie
 on the line; testing $\xi$ and $\xi\pm u$ suffices.
\item \src{06} For every ordinary real $c$ and positive infinitesimal $\eps$,
 all finite zeros of $\xip+c\eps$ lie on the line. Equivalently, the same
 holds for every $h\in\Ssym$ and every real infinitesimal $\eps$, with
 $\xip+\eps h^\#$.
\end{enumerate}
\end{corollary}
\begin{proof}
Each class lies in $\DJs$ and contains the test family of
\cref{gz:thm:stability}(iii), or its $w$-plane form. 06's converse uses the
constant $c=b_m$, where $b_m=i^m\xi^{(m)}(\rho)/m!\in\R\setminus\{0\}$ is
the leading coefficient of $\xi(\rho+iu)$ at an $m$-fold zero. This is the
sign choice of \cref{gz:thm:twosign} (a second route to the converse).
\end{proof}

\begin{remark}[General real entire functions \src{04}]\label{gz:rem:allh}
For a nonzero real entire $f$ and one fixed positive infinitesimal $\eta$,
all zeros of $f$ are real and simple if and only if, for every ordinary
real entire $h$, all finite zeros of $f^\#+\eta h^\#$ are real; testing
$h=\pm1$ suffices (04 Corollary 6.2). The forward direction is the last
clause of \cref{gz:prop:conservation}, and the converse is
\cref{gz:thm:twosign}. \Cref{gz:cor:criteria}(a) is the case $f=\Xi$.
\end{remark}

\begin{remark}[Exact versus shadow \src{04}]\label{gz:rem:shadow}
For constant perturbations every zero in the monad of $\rho$ has standard
part $\rho$, so the shadow statement is RH, while the exact statement for
both signs is RH plus simplicity. 04 calls the two ``inequivalent in
strength'' and the shadow form ``weaker''. The correct wording is
\emph{formally stronger}: the exact statement implies the shadow statement,
and it is equivalent to the shadow statement together with simplicity of
every nontrivial zero. The two can differ in truth value only if RH holds
and some nontrivial zero is multiple, which is not known to happen
(error 22).
\end{remark}

\begin{corollary}[Splitting scale \src{04, 06}]\label{gz:cor:scale}
At an $m$-fold critical-line zero, the off-line zero constructed in the
proof has displacement of valuation $v(\eta)/m$. A symmetric perturbation of
size $\eta$ can produce a much larger transverse displacement. For example,
$u^2+\eps=0$ gives $u=\pm i\sqrt\eps$, so $s=\rho+iu$ moves horizontally by
$\mp\sqrt\eps$, and the two zeros are exchanged by $J$ instead of being
fixed by it. For $u^3$, the zeros of $u^3+\eta$ are $-\eta^{1/3}$ and
$\eta^{1/3}(\tfrac12\pm\tfrac{\sqrt3}2i)$.
\end{corollary}
Both signs are needed in \cref{gz:thm:stability}(iii): a real double zero
with positive quadratic coefficient splits into two real zeros under
$-\eta$ and into two nonreal zeros under $+\eta$ \src{04}. The perturbed
functions keep both symmetries but not the Euler product, and they are not
proposed as new zeta functions. What the theorem shows is that the
functional-equation symmetries alone do not prevent infinitesimal off-line
splitting \src{04}. It proves neither RH nor simplicity \src{04, 05, 06}.
It concerns infinitesimal parameters, not a uniform real neighbourhood of
$0$ valid at all heights \src{05}. The class $\DJs$ is not claimed to be
maximal \merge.

\subsection{Multiplicative deformations and protected counts}\label{gz:sec:protected}
A multiplicative perturbation $f^\#(1+\eta h^\#)$ has exactly the finite
divisor of $f^\#$, because the second factor is a unit at every finite
point \src{04}. Additive and multiplicative deformations must not be
conflated.

Let $N(T)$ count the nontrivial zeros with $0<\Im\rho\le T$, with
multiplicity, and $N_0^{\rm s}(T)$ those that are simple and on the line.
For ordinary $T>0$ that is not a zero ordinate, put
$D_T=\{w:0<\Re w<T,\ |\Im w|<1\}$. Its boundary contains no zero of $\Xi$,
since zeros have $|\Im w|=|\tfrac12-\beta|<\tfrac12$, and it contains every
zero counted by $N(T)$.

\begin{theorem}[Protected counts \src{04}]\label{gz:thm:protected}
Let $h$ be an ordinary real even entire function and $\eta>0$ infinitesimal.
Let $N_\eta(T)$ count the zeros of $\Xi_\eta=\Xi^\#+\eta h^\#$ in $D_T^\#$
with multiplicity, and $N_{\eta,\R}(T)$ the real ones. Then
\[
N_\eta(T)=N(T),\qquad N_{\eta,\R}(T)\ge N_0^{\rm s}(T).
\]
Hence, for every ordinary $c$, $\liminf N_0^{\rm s}(T)/N(T)\ge c$ implies
$\liminf N_{\eta,\R}(T)/N_\eta(T)\ge c$, the $\liminf$ taken over ordinary
$T$ that are not zero ordinates.
\end{theorem}
\begin{proof}
Every finite zero has as standard part a zero of $\Xi$. Apply
\cref{gz:prop:conservation} to the finitely many zero monads in $D_T^\#$
and add. Each simple real zero gives one real deformed zero, by
\cref{gz:thm:protection} in the $w$-plane. The last statement compares
ordinary integer-valued functions of ordinary $T$.
\end{proof}
Evenness preserves $\Xi_\eta(-w)=\Xi_\eta(w)$; the local theorem needs only
realness. If $h$ is not real, a simple critical-line zero can move off the
line. The theorem proves neither RH nor simplicity. Admissible values of
$c$ are lower bounds for $\kappa_{\rm s}$, not for $\kappa_0$; they and
their statuses are in \cref{gz:sec:proportions}.

Call an ordinary critical-line zero a \emph{uniformly stable centre} if the
conclusion of \cref{gz:thm:protection} holds at it for every $h\in\Ssym$ and
every real infinitesimal parameter \src{06}. Every simple critical-line zero
is one, and the proof of \cref{gz:thm:twosign} shows that a multiple
critical-line zero is not. Proportions of stable centres are centre-indexed
counts with the classical total count as denominator. They are not counts of
displaced roots below a surcomplex height.

\part{Positive infinite real part}\label{gz:part:right}

\section{The Dirichlet algebra on \texorpdfstring{$\Tp$}{T+}}\label{gz:sec:dirichlet}

\subsection{Separated monomials}
Every $s\in\Tp$ has a unique decomposition
\begin{equation}\label{gz:eq:sdecomp}
s=X+c+\eps,\qquad X=\pp(\Re s)\in\Pi,\ X>0,\quad c\in\C,\quad\eps\in\mm .
\end{equation}
Since $\Im s$ is finite, $(\Im s)\log n$ is finite for every ordinary $n$,
and $n^{-s}=\Efp(-s\log n)$ needs no phase: it is the same for every
exponential of \cref{gz:sec:expconv}. Explicitly,
\begin{equation}\label{gz:eq:qn}
n^{-s}=q_n\,n^{-c}\,\Ex(-\eps\log n),\qquad q_n=\exp_{\No}(-X\log n).
\end{equation}
Each $q_n$ is a Conway monomial, since $X\log n\in\Pi$. They satisfy
$q_mq_n=q_{mn}$, $1=q_1>q_2>\cdots$, and $q_n/q_m\in\mm$ for $n>m$, so their
valuations form an increasing sequence of order type $\omega$.

\begin{theorem}[Arbitrary-coefficient Dirichlet realization \src{01--06}]\label{gz:thm:dirichlet}
Let $s\in\Tp$, and let $a:\N_{\ge1}\to\C$ be an arbitrary arithmetic function,
with no growth condition. Then
\[
\DD_s(a)=\Sh_{n\ge1}a(n)\,n^{-s}
\]
exists. The map $\DD_s$ is an injective unital $\C$-algebra homomorphism
from arithmetic functions under Dirichlet convolution
$(a*b)(n)=\sum_{d\mid n}a(d)b(n/d)$ to $\SC$. If $a\ne0$ and $n_0$ is its
least nonzero index, then $\DD_s(a)=a(n_0)\,n_0^{-s}(1+\eta)$ with $\eta\in\mm$.
\end{theorem}
\begin{proof}
Expand every term by \eqref{gz:eq:qn} and apply \cref{gz:lem:mixed} with
$\mathfrak t_n=q_n$ and the one infinitesimal $\eps$. This gives strong
summability, including all the expansions inside the terms. Large ordinary
coefficients $a(n)n^{-c}$ do not change supports. For a product, expand both
families. A fixed $q_k$ receives contributions only from the finitely many
factorizations $k=mn$, and $m^{-s}n^{-s}=k^{-s}$ by the exponential laws,
so regrouping gives the convolution. The function $\delta_1$ maps to $1$.
For the leading term, divide by $a(n_0)n_0^{-s}$. Every remaining term with
$n>n_0$ has leading monomial $q_n/q_{n_0}\in\mm$, and all its other
monomials are smaller, so the normalized tail is a strong family of
positive valuation. Its sum is infinitesimal, which gives the leading form
and hence injectivity.
\end{proof}
\begin{remark}[Scope \src{02, 04, 06}]\label{gz:rem:dirichlet-scope}
The theorem embeds the whole algebra of \emph{formal} Dirichlet series, for
example $a(n)=\exp(n^2)$ or $\exp(\exp n)$. It is a strong-normal-form
evaluation, not a statement about classical analytic continuation of such
data. It holds pointwise in $s$ and is not effective: equality of values
reduces to equality of all coefficients, and this gives no decision
procedure \src{02}. The units of the arithmetic algebra are exactly the $a$
with $a(1)\ne0$, although every nonzero image is invertible in $\SC$
\src{04, 06}. 04 states the same theorem in a set-sized Hahn field:
$\mathcal E_q(a)=\sum a(n)q^{\log n}$ with $q=t^\delta$ has valuation
$\delta\log n_0$ and leading coefficient $a(n_0)$. The realization above is
the case $q=e^{-X}$, whose real powers are the monomials $q_n$. 02, 03 and
04 point out that the multiplication itself is the familiar formal
Dirichlet-series algebra (compare \cite[\S2]{Pong}). What is new is the
exact evaluation in $\SC$.
\end{remark}

\begin{proposition}[Local analytic dependence \src{01, 02, 03, 05, 06}]\label{gz:prop:Dtaylor}
For $s\in\Tp$, every arithmetic function $a$ and \emph{every} $h\in\mm$,
\[
\DD_{s+h}(a)=\Sh_{k\ge0}\frac{h^k}{k!}\DD_s(\delta^ka),\qquad(\delta a)(n)=-(\log n)\,a(n),
\]
the double family being strongly summable. Hence
$D^k\DD_s(a)=\Sh_na(n)(-\log n)^kn^{-s}$, the local derivatives are fine
derivatives, and sum, product and quotient rules hold for every
infinitesimal increment.
\end{proposition}
\begin{proof}
Add $\supp h$ to the support set in the proof of \cref{gz:thm:dirichlet}
and expand $\Ex(-h\log n)$. \Cref{gz:lem:mixed} with two infinitesimals
justifies the regrouping. The derivation $\delta$ satisfies the Leibniz
rule because $\log(mn)=\log m+\log n$. For a quotient,
$\DD_s(\delta^kb)/\DD_s(b)=(-\log n_0)^k(1+\eta_k)/(1+\eta)$ is finite for
every $k$. So $\DD_{s+h}(b)/\DD_s(b)-1$ is a strong sum of terms of
positive valuation, hence infinitesimal, for every $h\in\mm$.
\end{proof}
02's Proposition 5.3 states the quotient rule only ``for $h$ sufficiently
small''. As the proof shows, every infinitesimal $h$ works (error 9).

\begin{definition}[Dirichlet zeta \src{01--06}]\label{gz:def:zD}
$\zD(s)=\DD_s(\mathbf1)=\Sh_{n\ge1}n^{-s}$ on $\Tp$, and
$\zD^{(k)}(s)=\DD_s(n\mapsto(-\log n)^k)$.
\end{definition}

\begin{theorem}[Euler identities and zero-freeness \src{01--06}]\label{gz:thm:euler}
For every $s\in\Tp$, with strong sums and the strong expansion of the
product:
\begin{align*}
\zD(s)&=\Ph_p(1-p^{-s})^{-1},&\zD(s)^{-1}&=\Sh_{n\ge1}\mu(n)n^{-s},\\
\Lgm\zD(s)&=\Sh_p\Sh_{k\ge1}\frac{p^{-ks}}k=\Sh_{n\ge2}\frac{\Lambda(n)}{\log n}n^{-s},&
-\frac{\zD'(s)}{\zD(s)}&=\Sh_{n\ge1}\Lambda(n)n^{-s}.
\end{align*}
Moreover $\zD(s)\in1+\mm$, so $\zD$ has no zeros or poles on $\Tp$, and:
\begin{enumerate}[label=(\alph*)]
\item $\zD(s)-\sum_{n\le N}n^{-s}=(N+1)^{-s}(1+\eta_N)$ with $\eta_N\in\mm$,
 for every ordinary $N\ge1$. In particular $\zD(s)-1\sim2^{-s}$ and
 $\Lgm\zD(s)\sim2^{-s}$.
\item $\zD^{(k)}(s)=(-\log2)^k2^{-s}(1+\eta_k)$ for $k\ge1$.
\item For real $x>\R$: $\zD(x)>1$, $(-1)^k\zD^{(k)}(x)>0$ for every
 ordinary $k$, and $x\mapsto\zD(x)$ is strictly decreasing on
 $\{x\in\No:x>\R\}$. Also $\zD(x)-1\prec x^{-m}$ for every ordinary $m$.
\item For completely multiplicative $a$, $\DD_s(a)=\Ph_p(1-a(p)p^{-s})^{-1}$.
 In particular Dirichlet $L$-series with character coefficients are defined
 on $\Tp$ and lie in $1+\mm$.
\end{enumerate}
\end{theorem}
\begin{proof}
The expanded Euler product has exactly one term for each ordinary integer,
by unique factorization, and its family is the family of
\cref{gz:thm:dirichlet}. The identity $\mu*\mathbf1=\delta_1$ gives the
reciprocal. $\zD-1$ is a strong sum of infinitesimals, so its near-one
logarithm exists. In the formal logarithm of $1+\sum_{n\ge2}n^{-s}$, the
coefficient at $n$ receives contributions only from the finitely many
factorizations of $n$ into factors at least $2$, of length at most
$\log_2n$. So the logarithm is computed coefficientwise in the arithmetic
algebra, where it is $1/k$ at $p^k$ and $0$ elsewhere. Differentiation by
\cref{gz:prop:Dtaylor}, or applying $\DD_s$ to $\Lambda*\mathbf1=\log$,
gives the von Mangoldt formula. (a) and (b) are the leading-term clause of
\cref{gz:thm:dirichlet} applied to the tail, to the logarithm coefficients
and to $(-\log n)^k$. (c) For real arguments every term beyond $1$ is
positive, and so is every term of $(-1)^k\zD^{(k)}$. If $x_2>x_1>\R$, the
differences $n^{-x_1}-n^{-x_2}$ ($n\ge2$) are positive and form a strong
family. A nonempty strong family of positive elements has positive sum,
because the leading coefficients at the largest occurring monomial are
positive; no mean value theorem is used. Finally
$2^{-x}x^m=\exp(-x\log2+m\log x)\in\mm$. (d) is the same factorization
argument, since each $a(p)p^{-s}$ is infinitesimal however large $a(p)$
is. For a Dirichlet character $\varpi$, $a(1)=1$.
\end{proof}
All six sources prove these identities: 01 Thms 7.3--7.4; 02 Thm 6.1 and
Cor 6.2; 03 Thm 7.3; 04 Thm 7.3; 05 Cor 8.3; 06 Thm 5.3 and its
complete-monotonicity paragraph. In 05's Corollary 8.3, $\zD'$ is used
before it is defined; here it is defined in \cref{gz:prop:Dtaylor} first
(error 31). 06 calls the zero-freeness theorem ``a genuine new domain
result''. It is a result on a new domain for the specified extension, and
no priority is claimed (error 42).

\begin{proposition}[Summability dichotomy \src{02}]\label{gz:prop:dichotomy}
Let $\Im s$ be finite. The raw family $(n^{-s})_{n\ge1}$ is strongly
summable if and only if $\Re s$ is positive infinite.
\end{proposition}
\begin{proof}
The positive infinite case is \cref{gz:thm:dirichlet}. If $\Re s$ is
finite, every term has a nonzero coefficient at exponent zero. If $\Re s$ is
negative infinite, its purely infinite part is negative, and the leading
monomials of $n^{-s}$ strictly increase with $n$. So the union of supports
is not well ordered in the valuation order.
\end{proof}

\begin{remark}[What the right-hand theory does not determine \src{01, 04, 05, 06}]\label{gz:rem:flat-dirichlet}
At positive infinite $X$, $e^{-X^2}\prec n^{-X}$ for every ordinary $n\ge2$,
because $X^2-X\log n$ is positive infinite \src{06}. So agreement with every
fixed Dirichlet asymptotic order is weaker than equality with the exact
strong sum. This is the zeta analogue of flat Gamma gauges. The finite
prescription $\zp$ and $\zD$ live on disjoint domains, and their union is
an explicitly specified partial function. No path from the finite halo to
$\Tp$ exists, and $\zD$ is not claimed to be the analytic continuation of
$\zp$ \src{01, 04, 05}. The $L$-series chart gives no surcomplex GRH,
since the hard region would again have finite real part and infinite
imaginary part \src{04}. $\Tp$ contains no point of the infinite-height
strip $\Vs$, so zero-freeness on $\Tp$ says nothing about RH at infinite
height \src{05, 06}.
\end{remark}

\begin{example}[At $\omega$ \src{02, 03, 04, 05}]\label{gz:ex:omega}
$\zD(\omega)=1+2^{-\omega}+3^{-\omega}+4^{-\omega}+\cdots$, where the ellipsis
denotes the complete strong family, and
$(\zD(\omega)-1-2^{-\omega})/3^{-\omega}=1+\eps$ with $\eps\in\mm_\R$. Further,
\begin{align*}
\Lgm\zD(\omega)&=2^{-\omega}+3^{-\omega}+\tfrac12\,4^{-\omega}+5^{-\omega}+7^{-\omega}+\tfrac13\,8^{-\omega}+\tfrac12\,9^{-\omega}+\cdots\quad(\text{no }6^{-\omega}),\\
\zD(s)^{-1}&=1-2^{-s}-3^{-s}-5^{-s}+6^{-s}-7^{-s}+10^{-s}-\cdots .
\end{align*}
Each $n^{-\omega}$ is smaller than every $\omega^{-N}$. At $s=\omega+c+\eps$
the first nonconstant scale is $2^{-c}2^{-\omega}\Ex(-\eps\log2)$.
\end{example}

\section{Inverse zeta on \texorpdfstring{$\Tp$}{T+}}\label{gz:sec:inverse}

This section is source 01's. Nothing else in the six sources treats inverse
zeta. Lagrange inversion is classical \cite{Gessel}, and inverse zeta has a
literature \cite{Kawalec}. The claim is the proved package with the
surreal domain and support assertions stated here, not priority for the
idea of inverting zeta.

For $n\ge3$ put $\alpha_n=\log_2n$ and $\beta_n=\alpha_n-1=\log_2(n/2)$, and
write $\beta=\beta_3=\log_2(3/2)\approx0.584963$.

\begin{theorem}[Inverse-zeta fibre \src{01}]\label{gz:thm:headline}
For every positive infinitesimal $u\in\No$, the fibre
$\{s\in\Tp:\zD(s)=1+u\}$ consists of exactly one point $s_k(u)$ for each
ordinary $k\in\Z$. Its imaginary standard part is $-2\pi k/\log2$, and
\[
s_k(u)=-\frac{\log u}{\log2}-\frac{2\pi ik}{\log2}+\frac{e^{2\pi ik\beta}}{\log2}u^\beta+\frac u{\log2}-\frac{(2\beta+1)e^{4\pi ik\beta}}{2\log2}u^{2\beta}+o(u^{2\beta}).
\]
The full expansion is a strongly summable generalized power series, with
coefficients and exponent collisions given in \cref{gz:thm:arithmetic}. The
branch $s_0(u)$ is real, $s_{-k}(u)=\overline{s_k(u)}$, and $s_0$ inverts the
decreasing bijection $x\mapsto\zD(x)-1$ from $\{x>\R\}$ onto the positive
infinitesimals.
\end{theorem}
Here $o(v)$ means an element whose quotient by $v$ is infinitesimal, a
dominance statement at the given argument. The proof occupies the rest of
this section.

\begin{lemma}[The weight monoid \src{01}]\label{gz:lem:weightmonoid}
For finitely supported $\nu=(\nu_n)_{n\ge3}$ with $\nu_n\in\N$, put
$|\nu|=\sum\nu_n$ and $B(\nu)=\sum\nu_n\beta_n$. For every real $R$ only
finitely many $\nu$ satisfy $B(\nu)\le R$.
\end{lemma}
\begin{proof}
Every generator is at least $\beta_3>0$, so $|\nu|\le R/\beta_3$, and an
occurring generator satisfies $\beta_n\le R$, that is, $n\le2^{R+1}$.
\end{proof}
The exponents are not distinct: $\beta_6=\beta_3+\beta_4$. It is finite
multiplicity, not absence of collisions, that allows evaluation and
regrouping.

Fix a positive infinitesimal $u$ and $k\in\Z$, and put
$c_{n,k}=e^{2\pi ik\alpha_n}=e^{2\pi ik\beta_n}$. Consider
\begin{equation}\label{gz:eq:Veq}
V+\Sh_{n\ge3}c_{n,k}\,u^{\beta_n}V^{\alpha_n}=1,\qquad V\in1+\mm,
\end{equation}
with near-one binomial powers of $V$.

\begin{lemma}[Unique near-one solution \src{01}; existence step completed \merge]\label{gz:lem:Vunique}
Equation \eqref{gz:eq:Veq} has exactly one solution $V_k(u)$ in $1+\mm$. It
is the strong evaluation of a generalized power series in $u$ with
exponents in the monoid generated by the $\beta_n$.
\end{lemma}
\begin{proof}
\emph{Existence.} Take formal variables $x_n$ of weight $\beta_n$. These
weights satisfy the hypothesis of \cref{gz:lem:evalhom}, by
\cref{gz:lem:weightmonoid}. The equation
\[
W=-\sum_{n\ge3}x_n(1+W)^{\alpha_n}
\]
has a unique formal solution with zero constant term. The right side is a
formally summable family, each of whose terms multiplies every coefficient
of $W$ by a positive-weight variable, so the coefficients are determined
recursively in increasing weight. Its explicit coefficients are given in
\cref{gz:lem:multilagrange}. Apply the evaluation $x_n\mapsto c_{n,k}u^{\beta_n}$
of \cref{gz:lem:evalhom}. Being a strongly additive ring homomorphism with
$\ev((1+W)^a)=(1+\ev W)^a$, it turns the formal equation into
\eqref{gz:eq:Veq} for $V=1+\ev W$, and $\ev W\in\mm$. (01 asserted this
step without a lemma; error 3.)

\emph{Uniqueness among all near-one elements.} Let $V,V'\in1+\mm$ both
solve \eqref{gz:eq:Veq}. For $a\in\C$ the formal divided-difference identity
$V^a-V'^a=(V-V')D_a(V-1,V'-1)$ holds with $D_a\in\C[[X,Y]]$ of constant
term $a$, and all evaluated supports lie in the Neumann monoid of
$\supp(V-1)\cup\supp(V'-1)$, independently of $a$. Subtracting the two
equations gives
$(V-V')\bigl(1+\Sh_nc_{n,k}u^{\beta_n}D_{\alpha_n}(V-1,V'-1)\bigr)=0$. The
bracket is $1$ plus an infinitesimal strong sum, hence a unit, so $V=V'$.
\end{proof}

Define
\[
s_k(u)=-\frac{\log u}{\log2}-\frac{2\pi ik}{\log2}-\frac{\Lgm V_k(u)}{\log2}.
\]
Its real part is positive infinite and its imaginary part finite, so
$s_k(u)\in\Tp$.

\begin{theorem}[Complete fibre \src{01}]\label{gz:thm:fibre}
For every $0<u\in\mm_\R$, $\{s\in\Tp:\zD(s)=1+u\}=\{s_k(u):k\in\Z\}$,
without repetition, and branch $k$ is characterized by
$\st(\Im s)=-2\pi k/\log2$.
\end{theorem}
\begin{proof}
By construction $2^{-s_k}=uV_k$, and $n^{-s_k}=c_{n,k}u^{\alpha_n}V_k^{\alpha_n}$
for $n\ge3$. After division by $u$, $\zD(s_k)-1=u$ is exactly
\eqref{gz:eq:Veq}. Conversely, let $s=x+iy\in\Tp$ with
$\zD(s)-1=u$. The first-tail formula gives $u=2^{-s}(1+\eta)$, so
$V=2^{-s}/u\in1+\mm$. Its modulus gives $\exp_{\No}(-x\log2)/u\in1+\mm_\R$,
and its phase gives $e^{-i\st(y)\log2}=1$, so $\st y=-2\pi k/\log2$
for a unique $k\in\Z$. The kernel of $\cis$ and the near-one logarithm give
$-s\log2=\log u+\Lgm V+2\pi ik$. Substituting into every $n^{-s}$ and
dividing by $u$ gives \eqref{gz:eq:Veq}, and \cref{gz:lem:Vunique} forces
$V=V_k(u)$, so $s=s_k(u)$. Distinct $k$ have distinct $\st\Im s$.
\end{proof}
The exhaustion step matters. A formal inverse in one power-series ring
would not exclude solutions whose supports lie outside it. The
divided-difference argument and the first-tail formula exclude them
throughout $\Tp$.

\begin{corollary}[Conjugation and the real inverse \src{01}]\label{gz:cor:realinverse}
$V_{-k}(u)=\overline{V_k(u)}$ and $s_{-k}(u)=\overline{s_k(u)}$. The point
$s_0(u)$ is the unique real solution, and $x\mapsto\zD(x)-1$ is a bijection
from $\{x\in\No:x>\R\}$ onto $\mm_{\R,>0}$.
\end{corollary}
\begin{proof}
Conjugation sends $c_{n,k}$ to $c_{n,-k}$ and preserves strong sums;
uniqueness gives the first formula. A real solution has $k=0$.
Monotonicity and positivity are \cref{gz:thm:euler}(c).
\end{proof}

\begin{corollary}[All nonzero infinitesimal values \src{01}]\label{gz:cor:complexfibre}
The image of $s\mapsto\zD(s)-1$ on $\Tp$ is exactly $\mm\setminus\{0\}$. For
each such $w$ the fibre has exactly one point for every $k\in\Z$, obtained
by replacing $u^a$ with $w^a=\Efp(a\Log w)$ and $\log u$ with $\Log w$.
Replacing $\Log w$ by $\Log w+2\pi ij$ relabels branch $k$ as $k+j$.
\end{corollary}
\begin{proof}
\Cref{gz:lem:realpowers,gz:lem:evalhom} cover the complex powers, and the
existence and divided-difference arguments are unchanged. The condition
$2^{-s}/w\in1+\mm$ again determines $k$. Conversely every value is a nonzero
infinitesimal by \cref{gz:thm:euler}. A change of $\Log w$ by $2\pi ij$
multiplies $w^a$ by $e^{2\pi ija}$ and shifts the base point by
$-2\pi ij/\log2$.
\end{proof}
The statement concerns fibres and values. No covering-map claim in the
fine topology is made \src{01}. The fibres over the right wing are in
\cref{gz:cor:wingfibre}.

\subsection{Coefficients}
For a multi-index $\nu$ put $\nu!=\prod\nu_n!$. Rising factorials have
ordinary length, and $(a)_0=1$.

\begin{lemma}[Weighted Lagrange formula \src{01}]\label{gz:lem:multilagrange}
For the formal solution $V=1+W$,
\[
V^r=1+\sum_{\nu\ne0}\frac{(-1)^{|\nu|}r(r+B+1)_{|\nu|-1}}{\nu!}x^\nu,\qquad
\log V=\sum_{\nu\ne0}\frac{(-1)^{|\nu|}(B+1)_{|\nu|-1}}{\nu!}x^\nu,
\]
with $B=B(\nu)$.
\end{lemma}
\begin{proof}
Below a weight cutoff only finitely many variables occur. Write
$W=t\Theta(W)$ with $\Theta(T)=-\sum_nx_n(1+T)^{\alpha_n}$ and apply Lagrange's
formula $[t^d]F(W)=\frac1d[T^{d-1}]F'(T)\Theta(T)^d$ to $F=(1+T)^r$
\cite{Gessel}. Since $\Theta$ is linear in the $x_n$, a monomial $x^\nu$
occurs only in $\Theta^d$ with $d=|\nu|$, with coefficient
$\frac{(-1)^dd!}{\nu!}(1+T)^{d+B}$, since $\sum\nu_n\alpha_n=d+B$. Extracting
$T^{d-1}$ from $r(1+T)^{r-1+d+B}$ gives
$\frac{(-1)^dr(d-1)!}{\nu!}\binom{r+d+B-1}{d-1}$, which is the stated
coefficient. Differentiate at $r=0$ for $\log V$. The univariate identity is
a universal polynomial identity in the coefficients, so it needs no analytic
hypothesis.
\end{proof}
Evaluating \cref{gz:lem:multilagrange} at $x_n=c_{n,k}u^{\beta_n}$ gives
\[
s_k(u)=-\frac{\log u}{\log2}-\frac{2\pi ik}{\log2}+\frac1{\log2}\Sh_{\nu\ne0}\frac{(-1)^{|\nu|+1}(B+1)_{|\nu|-1}}{\nu!}\,e^{2\pi ikB}u^B,
\]
an exact strong value at every positive infinitesimal $u$, whatever its
normal form, by \cref{gz:lem:weightmonoid,gz:lem:evalhom} \src{01}.

Let $\mathcal M=\{\prod_{j=1}^d(n_j/2):d\ge0,\ n_j\ge3\}\subseteq\Q_{\ge1}$,
let $A_d(r)$ count the ordered $d$-tuples with product $r$, and for $r>1$ put
\[
C_r=\sum_{d\ge1}\frac{(-1)^{d+1}A_d(r)}{d!}(1+\log_2r)_{d-1},
\]
a finite sum, since $d\le\log r/\log(3/2)$ and each $n_j\le2r$.

\begin{theorem}[Arithmetic coefficient formula \src{01}]\label{gz:thm:arithmetic}
For every positive infinitesimal $u$ and $k\in\Z$,
\[
s_k(u)=-\frac{\log u}{\log2}-\frac{2\pi ik}{\log2}+\frac1{\log2}\Sh_{r\in\mathcal M,\,r>1}C_r\,e^{2\pi ik\log_2r}\,u^{\log_2r}.
\]
The exponents are locally finite in $\R_{>0}$, and $C_r$ includes every
cancellation between distinct factorizations. The same formula holds for
nonzero complex infinitesimal $w$.
\end{theorem}
\begin{proof}
Since $2^{B(\nu)}=\prod_n(n/2)^{\nu_n}$, two multi-indices have the same
exponent exactly when they have the same product $r$. For fixed $d=|\nu|$
and $r$, $\sum_{\nu}1/\nu!=A_d(r)/d!$, and every term has the phase
$e^{2\pi ikB}=e^{2\pi ik\log_2r}$. Regroup the evaluated
\cref{gz:lem:multilagrange}, which is legitimate by strong summability.
\end{proof}
\begin{center}\small
\begin{tabular}{@{}lllllllll@{}}
\toprule
$r$ & $3/2$ & $2$ & $9/4$ & $5/2$ & $3$ & $27/8$ & $7/2$ & $15/4$\\
$C_r$ & $1$ & $1$ & $-\tfrac12(1+\log_2\tfrac94)$ & $1$ & $-\log_23$ & $\tfrac16(1+L)(2+L)$ & $1$ & $-(1+\log_2\tfrac{15}4)$\\
\bottomrule
\end{tabular}
\end{center}
Here $L=\log_2(27/8)$, and $C_4=-\tfrac12$. The coefficient at $r=3$ shows
why the collision bookkeeping matters. The length-one factorization
$6/2$ contributes $1$, and the two ordered length-two factorizations
$(3/2)(4/2)$ contribute $-(1+\log_23)$, so $C_3=-\log_23$. The $23$
coefficients with $r<8$ are recorded in
\texttt{data/01-inverse-zeta-inverse-coefficients.csv}.
The smallest elements of $\mathcal M$ above $1$ are $3/2<2<9/4<5/2$, so the
first exponents are $\beta<1<2\beta<\log_2(5/2)$. With $C_{3/2}=C_2=1$,
$C_{9/4}=-\tfrac12(1+2\beta)$ and the phase $e^{2\pi ik\log_2r}$, this gives
the expansion displayed in \cref{gz:thm:headline} \src{01}.

\begin{theorem}[The integer-power sector \src{01}]\label{gz:thm:integersector}
$C_{2^m}=(-1)^{m+1}/m$ for every ordinary $m\ge1$, independently of the
branch. The only rational exponents occurring are the positive integers, so
the rational-exponent sector of the correction is exactly
$\Lgm(1+u)/\log2$.
\end{theorem}
\begin{proof}
Every member of $\mathcal M$ has denominator a power of $2$. So discarding
the indices with an odd prime factor is a multiplicative projection of the
generalized series ring. Applied to \eqref{gz:eq:Veq}, only $n=2^a$
($a\ge2$) survive, with phase $1$. The projected equation
$V_2+\sum_{a\ge2}u^{a-1}V_2^a=V_2/(1-uV_2)=1$ gives $V_2=(1+u)^{-1}$ and
$-\log V_2=\Lgm(1+u)$. If $\log_2r=p/q$ in lowest terms, then $r^q=2^p$
forces $r=2^m$ and $q=1$, and every $2^m$ is a product of $m$ factors
$4/2$.
\end{proof}
Equivalently $\sum_{d=1}^m\frac{(-1)^{d+1}}{d!}\binom{m-1}{d-1}(m+1)_{d-1}=\frac{(-1)^{m+1}}m$.
The projection acts on the formal $u$-expansion. For a nonmonomial $u$ it
must not be confused with discarding monomials of the evaluated normal form
\src{01}.

\begin{corollary}[The real branch is rightmost \src{01}]\label{gz:cor:rightmost}
For $k\ne0$,
$s_k(u)-s_0(u)+\frac{2\pi ik}{\log2}\sim\frac{e^{2\pi ik\beta}-1}{\log2}u^\beta$ and
$\Re s_k(u)-s_0(u)\sim\frac{\cos(2\pi k\beta)-1}{\log2}u^\beta<0$. The fibre is
not an exact vertical arithmetic progression.
\end{corollary}
\begin{proof}
$\beta$ is irrational, since otherwise a power of $3$ would equal a power of
$2$. So $e^{2\pi ik\beta}\ne1$. Subtract the two expansions: the first
exponents are $\beta<1<2\beta<\log_2(5/2)$.
\end{proof}
\begin{example}[\src{01}]
$s_0(\omega^{-1})=\frac{\log\omega}{\log2}+\frac{\omega^{-\beta}}{\log2}+\frac{\omega^{-1}}{\log2}-\frac{2\beta+1}{2\log2}\omega^{-2\beta}+o(\omega^{-2\beta})$,
which exhibits the intermediate scale $\omega^{-\log_2(3/2)}$.
\end{example}
The ordinary-complex illustration of 01's suite takes $u=10^{-8}$ and the
terms with $r<8$, and compares them with roots of $\zeta(s)=1+u$ computed
at $100$ digits:
\begin{gather*}
s_0\approx26.575454935879472,\qquad
s_{\pm1}\approx26.575398807000944\mp9.064735632317692\,i,\\
s_{\pm2}\approx26.575439317457936\mp18.129414142729476\,i,
\end{gather*}
with errors of about $1.0\times10^{-23}$ and zeta residuals below
$7\times10^{-32}$. These numbers illustrate signs and coefficients. They are
not error bounds or surreal proofs \src{01}.

\section{Algebraic independence at infinite arguments}\label{gz:sec:independence}

Two sources prove independence theorems for zeta values at a fixed
infinite argument, for different families over different base fields. 02
treats dilations $\zD^{(j)}(kx)$ at a positive infinite real $x$, together
with $x$, $\LSt(x)$ and $\GSt(x)$. 03 treats real shifts
$\zD^{(k)}(X+a)$ at a positive infinite $X\in\Pi$, over fields containing a
slow Hahn field and Gamma values. Both keep the standard literature in
view. Arithmetic-function independence has a classical theory, with Pong
\cite{Pong} an explicit antecedent. No novelty is claimed for the
functional-independence principle \src{02}. What is new is the exact
evaluation at one surreal argument, over the stated base fields. Standard
parts destroy the information, since $\st\zD(kx)=1$, and nothing here
concerns the ordinary values $\zeta(2),\zeta(3),\dots$ \src{02}.

\subsection{Dilations}\label{gz:sec:dilation}
For distinct primes $p_1,\dots,p_M$, the prime projection
$\operatorname{pr}(a)=\sum_{\alpha\in\N^M}a(p_1^{\alpha_1}\cdots p_M^{\alpha_M})Z^\alpha\in\C[[Z_1,\dots,Z_M]]$
preserves convolution, since every divisor of an integer supported on these
primes is supported on them too. It turns $\delta$ into
$\delta_Z=-\sum_i(\log p_i)Z_i\partial_{Z_i}$. Let $a_{k,0}$ be the
indicator of perfect $k$th powers and $a_{k,j}=\delta^ja_{k,0}$. Then
\[
A_{k,0}=\operatorname{pr}(a_{k,0})=\prod_{i=1}^M(1-Z_i^k)^{-1},\quad A_{k,j}=\delta_Z^jA_{k,0},\quad
\DD_s(a_{k,j})=k^j\zD^{(j)}(ks).
\]
The factor $k^j$ appears because $\delta$ differentiates in $s$, whereas
$\zD^{(j)}(ks)$ differentiates before dilating.

\begin{lemma}[A nonzero Jacobian \src{02}]\label{gz:lem:jacobian}
Let $\mathcal K\subset\N_{\ge1}$ be finite, $J_k\ge0$ for $k\in\mathcal K$, and
$M=\sum_k(J_k+1)$. For $M$ distinct primes, the Jacobian of
$\{A_{k,j}:k\in\mathcal K,\ 0\le j\le J_k\}$ with respect to $Z_1,\dots,Z_M$ is
nonzero.
\end{lemma}
\begin{proof}
The lowest nonconstant homogeneous part of $A_{k,j}$ is
$\sum_i(-k\log p_i)^jZ_i^k$. The lowest-degree part of the Jacobian
is therefore $J_0=\det[k(-k\log p_i)^jZ_i^{k-1}]$, and higher parts
cannot cancel it. Partition the columns into groups $I_k$ of sizes $J_k+1$
and look at the coefficient of $\prod_k\prod_{i\in I_k}Z_i^{k-1}$.
Distinct $k$ give distinct exponents, including $k=1$ with exponent $0$, so
the monomial forces the rows of $k$ onto the columns of $I_k$. Up to a sign
and powers of $k$, the coefficient is $\prod_k\det[(\log p_i)^j]_{j\le J_k,\,i\in I_k}$,
a product of Vandermonde determinants in distinct reals.
\end{proof}

\begin{lemma}[Formal Jacobian criterion \src{02}]\label{gz:lem:jaccriterion}
Power series $F_1,\dots,F_M\in\C[[Z_1,\dots,Z_M]]$ with nonzero Jacobian are
algebraically independent over $\C$.
\end{lemma}
\begin{proof}
Let $P\ne0$ be a relation of least degree. Differentiating $P(F)=0$ gives a
linear system with the Jacobian as matrix, invertible over the fraction
field. So every $\partial_jP(F)=0$, and by minimality every $\partial_jP$ is
zero. In characteristic zero $P$ is then constant, a contradiction.
\end{proof}

\begin{theorem}[Dilation jets \src{02}]\label{gz:thm:zetajets}
For every $s\in\Tp$, complex or real, $\{\zD^{(j)}(ks):k\ge1,\ j\ge0\}$ is
algebraically independent over $\C$. So is
$\{\zD^{(j)}(qs):q\in\Q_{>0},\ j\ge0\}$.
\end{theorem}
\begin{proof}
A relation uses finitely many pairs; enlarge them to initial segments
$j\le J_k$. By the two lemmas the projected $A_{k,j}$ are independent, hence
so are the $a_{k,j}$, and injectivity of $\DD_s$ with the nonzero scalars
$k^j$ transfers this to the values. For rational dilations use a common
denominator $M$ and work at $s/M\in\Tp$.
\end{proof}
The only arithmetic input is that distinct primes have distinct logarithms.
No conjecture on logarithms of primes (such as Schanuel's) is used.

\subsection{Leading forms and the Gamma scale}\label{gz:sec:leading}
Fix a positive infinite real $x$. Let $\mathcal B_x=\DD_x(\Arith)$ and
$\mathcal F_x^{\rm Dir}=\Frac(\mathcal B_x[x])$. The needed scale
inequalities, for ordinary $c>0$ and $d,h\ge0$, are
$(\log x)^h\prec x$ ($h\ge1$) and $x^d(\log x)^h\prec e^{cx}$. They follow
from the elementary theory of $(\No,\exp)$ \cite{vdDE}.

\begin{lemma}[Polynomial coefficients \src{02}]\label{gz:lem:polyleading}
Let $a_0,\dots,a_D\in\Arith$, not all zero, and $P_n(T)=\sum_da_d(n)T^d$. Let
$n_0$ be the least index with $P_{n_0}\ne0$, with highest term $cT^{d_0}$.
Then $\sum_dx^d\DD_x(a_d)=c\,x^{d_0}n_0^{-x}(1+\eta)$ with $\eta\in\mm$. So
$\Arith[T]\to\SC$, $T\mapsto x$, is injective.
\end{lemma}
\begin{proof}
The left side is $\Sh_nP_n(x)n^{-x}$, and $P_{n_0}(x)=cx^{d_0}(1+\eta')$.
After division by the $n_0$ term, the term at $n>n_0$ has monomial part
$(q_n/q_{n_0})x^{e}$ with $|e|\le D$. Since $q_n/q_{n_0}=\exp(-x\log(n/n_0))$
and $x^D\prec\exp(x\log(n/n_0))$, we get $q_n/q_{n_0}\prec x^{-D}$, and every
such monomial is infinitesimal. The normalized tail is a finite sum of strong
arithmetic families times fixed factors, with support of positive
valuation, so its sum is infinitesimal.
\end{proof}
The step $q_n/q_{n_0}\prec x^{-D}$ is the one line that 02's proof leaves
implicit (error 10).

\begin{theorem}[Leading forms \src{02}]\label{gz:thm:fieldleading}
Every nonzero $f\in\mathcal F_x^{\rm Dir}$ is uniquely
$c\,x^dq^{-x}(1+\eta)$ with $c\in\C^\times$, $d\in\Z$, $q\in\Q_{>0}$ and
$\eta\in\mm$. Hence the family $\{\zD^{(j)}(kx)\}$ is algebraically
independent over $\C(x)$.
\end{theorem}
\begin{proof}
Apply \cref{gz:lem:polyleading} to a numerator and a denominator. For
uniqueness, $c'x^{d'}q'^{-x}\sim1$ forces $q'=1$ by exponential separation,
then $d'=0$, then $c'=1$. For the second claim, clear denominators; the
injectivity in \cref{gz:lem:polyleading} makes every coefficient vanish as an
arithmetic function, and \cref{gz:thm:zetajets} finishes.
\end{proof}

\begin{lemma}[A logarithmic scale \src{02}]\label{gz:lem:logscale}
$\LSt(x)=x\log x\,(1+\eps)$ with $\eps\in\mm_\R$. $\LSt(x)$ is transcendental
over $\mathcal F_x^{\rm Dir}$, and every nonzero element of
$\mathcal F_x^{\rm Dir}(\LSt(x))$ is uniquely
$c\,x^d(\log x)^hq^{-x}(1+\eta)$ with $c\in\C^\times$, $d,h\in\Z$,
$q\in\Q_{>0}$ and $\eta\in\mm$.
\end{lemma}
\begin{proof}
In $P(x,\LSt(x))$ the monomial $U^dV^h$ has leading form $x^{d+h}(\log x)^h$.
The lexicographically largest pair $(d+h,h)$ determines a unique original
$(d,h)$, whose term dominates the others. With arithmetic coefficients,
write the expression as $\Sh_nP_n(x,\LSt(x))n^{-x}$ with uniformly bounded
degrees. Choose the least $n_0$, and then the largest pair. Later indices
have an exponential disadvantage that dominates fixed powers of $x$ and
$\log x$. So $\Arith[U,V]\to\SC$, $(U,V)\mapsto(x,\LSt(x))$, is injective,
and clearing denominators gives transcendence over $\mathcal F_x^{\rm Dir}$.
Applying the leading term to a numerator and a denominator gives the
leading form. Uniqueness follows by comparing, in turn, the exponential
scale, the power of $x$, the power of $\log x$ and the coefficient.
\end{proof}

\begin{theorem}[Joint scale separation \src{02}]\label{gz:thm:joint}
$\LSt(x)$ and $\GSt(x)$ are algebraically independent over
$\mathcal F_x^{\rm Dir}$. In fact $\GSt(x)$ is transcendental over
$\mathcal F_x^{\rm Dir}(\LSt(x))$. More generally, every positive real $G$
with $\log G/x>\R$ is transcendental over $\mathcal F_x^{\rm Dir}(\LSt(x))$.
\end{theorem}
\begin{proof}
A nonzero coefficient $b$ has $\log|b|/x=-\log q+\text{infinitesimal}$, a
finite slope, whereas $\log\GSt(x)/x=\LSt(x)/x\sim\log x$ is infinite. In
$\sum_{r\le R}b_r\GSt(x)^r=0$ with $b_R\ne0$, each lower term relative to
the top term is $\exp(\log|b_r/b_R|-(R-r)\LSt(x))\in\mm$. The same argument
applies to any $G$ with infinite slope.
\end{proof}

\begin{theorem}[Principal independence \src{02}]\label{gz:thm:indep-main}
For every positive infinite real $x$, the family
$\{x,\LSt(x),\GSt(x)\}\cup\{\zD^{(j)}(kx):k\ge1,\ j\ge0\}$ is algebraically
independent over $\C$. The derivatives are coordinate derivatives.
\end{theorem}
\begin{proof}
\Cref{gz:thm:zetajets} gives the jets. \Cref{gz:lem:polyleading} adjoins
$x$, \cref{gz:lem:logscale} adjoins $\LSt(x)$ even over $\mathcal F_x^{\rm Dir}$,
and \cref{gz:thm:joint} adjoins $\GSt(x)$.
\end{proof}
The Gamma step is robust: an infinitesimal relative perturbation of the
baseline changes nothing. The log-Gamma step uses the scale $x\log x$ of
$\LSt$, not a uniqueness property of surreal Gamma. The theorem is
compatible with the gauge nonuniqueness of the gamma-functions report, whose
members all have the same growth scales.

\subsection{Shifts}\label{gz:sec:shift}
\begin{lemma}[Chebyshev determinant \src{03}]\label{gz:lem:chebyshev}
For distinct reals $a_1,\dots,a_d$, integers $m_i\ge1$, $r=\sum m_i$, and
distinct reals $t_1,\dots,t_r$, the matrix
$[e^{-a_it_j}(-t_j)^k]_{(i,k),j}$ ($0\le k<m_i$) is nonsingular over $\C$.
\end{lemma}
\begin{proof}
A nonzero real exponential polynomial $\sum_ie^{-a_it}P_i(t)$ with
$\deg P_i<m_i$ has at most $r-1$ real zeros. For one exponent this is the
polynomial bound. Otherwise multiply by $e^{a_1t}$ and differentiate $m_1$
times. The first polynomial disappears, and each other term becomes
$e^{(a_1-a_i)t}(D+a_1-a_i)^{m_1}P_i$, which is nonzero of the same degree.
Rolle's theorem and induction finish. A singular real matrix would give such
a function vanishing at $r$ points, and a complex null vector gives a real
one.
\end{proof}
This is the only place where an ordinary mean-value argument is used. No
surreal Rolle theorem is invoked.

\begin{lemma}[Finite-prime independence \src{03}; proof completed \merge]\label{gz:lem:primeindependence}
Let $F\supseteq\C$ be a field, $S$ a finite set of primes, and $x_p$ the
arithmetic function supported at $p$ with value $1$. For distinct reals $a_i$
and integers $m_i\ge1$, the functions $A_{a_i,k}(n)=n^{-a_i}(-\log n)^k$
($0\le k<m_i$, value at $n=1$ as in \cref{gz:def:zD}) satisfy no nonzero
polynomial relation with coefficients in $F[x_p:p\in S]$, or, after clearing
denominators, in its fraction field.
\end{lemma}
\begin{proof}
Let $r=\sum m_i$ and choose $r$ distinct primes $q_j\notin S$. Projecting onto
integers supported on $S\cup\{q_j\}$ is an algebra map onto $F[[x_p,t_j]]$.
With $E_i=\prod_{p\in S}(1-p^{-a_i}x_p)^{-1}$ and
$D=-\sum_p(\log p)x_p\partial_{x_p}-\sum_j(\log q_j)t_j\partial_{t_j}$, the
projection of $A_{a_i,k}$ is
$Y_{i,k}=D^k\bigl(E_i\prod_j(1-q_j^{-a_i}t_j)^{-1}\bigr)$. Every
$t$-coefficient of $Y_{i,k}$ is a rational function in the $x_p$ that is
regular at $x=0$: it is obtained from products of the factors
$(1-p^{-a_i}x_p)^{-1}$ by $D_S=-\sum_p(\log p)x_p\partial_{x_p}$. Hence
$Y_{i,k}\in R[[t]]$, where $R=F(x)\cap F[[x]]$ is the local ring at $0$, and
$R[[t]]$ embeds both in $F(x)[[t]]$ and in $F[[x,t]]$. At $t=0$
the Jacobian has entries
$q_j^{-a_i}\sum_{b\le k}\binom kb D_S^{k-b}E_i\,(-\log q_j)^b$. For
each $i$ its row block is a lower-triangular transform, with diagonal
$E_i\ne0$, of the rows $q_j^{-a_i}(-\log q_j)^b$. So the Jacobian is
invertible over $F(x)$ by \cref{gz:lem:chebyshev} with $t_j=\log q_j$. The
formal inverse function theorem (solve degree by degree, one finite linear
system over $F(x)$ at each degree) shows that the $Y_{i,k}$ are algebraically
independent over $F(x)$. A relation $P$ with coefficients in $F[x_p]\subseteq R$
satisfies $P(Y)=0$ in $F[[x,t]]$, hence in $R[[t]]$ and in $F(x)[[t]]$, so $P=0$.
\end{proof}
03's proof works in $F(x)[[t]]$ although the projected objects live in
$F[[x,t]]$, and $F[[x]]$ is not a subring of $F(x)$. The passage through
$R[[t]]$ supplies the missing step (error 17). \emph{Antecedent.} For $F=\C$
the lemma is a special case of Pong's Theorem 5.5 \cite{Pong}: take the
function $f=\mathbf1$ and $\alpha_i=-a_i$, independence over the algebra
$\mathcal S\supseteq\C[x_p]$ of arithmetic functions with finitely generated
support. Pong's proof also chooses fresh primes and shows that a determinant
is nonzero. The $F$-general form follows from the case $F=\C$ by linear
disjointness. The 2020 Addendum \cite{PongAdd} was not checked. 03 cites Pong as an antecedent but presents its proof as
self-contained. The new content of this subsection is the Hahn-support
separation in \cref{gz:lem:sloweval} and \cref{gz:thm:indep-slow,gz:thm:gamma-zeta-independence}
(error 18). Pong's Corollary 5.6 and Example 5.5 (Ostrowski) are the
classical-analytic shadow of \cref{gz:thm:indep-monomial}.

Fix $X\in\Pi$, $X>0$, and write $q_n=e^{-X\log n}$ as in \cref{gz:sec:dirichlet}.
Then $\ev_X(a)=\Sh_na(n)q_n$ is the $c=\eps=0$ case of
\cref{gz:thm:dirichlet}, and it is injective \src{03}.

\begin{theorem}[Shift jets over a monomial field \src{03}]\label{gz:thm:indep-monomial}
$\{\zD^{(k)}(X+a):a\in\R,\ k\in\N\}$ is algebraically independent over
$M_X=\C(e^{-\lambda X}:\lambda\in\R)$.
\end{theorem}
\begin{proof}
Take finitely many shifts $a_i$, all derivatives $k<m_i$, and a relation
cleared of denominators, $\sum_{j\le N}e^{-\lambda_jX}P_j(Y)=0$, with distinct
$\lambda_j$ and complex polynomials $P_j$, not all zero. Each $P_j(Y)$ is the
evaluation of an arithmetic function with exponents in $\log\N_{\ge1}$. So
terms from different summands can share a monomial only if
$\lambda_j-\lambda_\ell\in\log\Q_{>0}$. Distinct cosets modulo $\log\Q_{>0}$
have disjoint supports and vanish separately. In one coset, factor out one
$e^{-\lambda_0X}$; the other factors are $r^{-X}$ with $r\in\Q_{>0}$,
involving a finite set $S$ of primes. Clearing negative prime exponents
gives a nonzero relation $Q((p^{-X})_{p\in S},Y)=0$. By injectivity of
$\ev_X$ this is a relation among arithmetic functions, contradicting
\cref{gz:lem:primeindependence} with $F=\C$.
\end{proof}
The base field must be a \emph{generated} field. Over a Hahn completion
containing the zeta series, transcendence would fail trivially \src{03}.

Assume now that $X$ is a positive infinite Conway monomial, and let
$F^{\rm Hahn}_X=\{\Sh_{r\in S}c_rX^{-r}:S\subseteq\R\text{ well ordered}\}\cong\C((t^\R))$
with $t=X^{-1}$. Since $X/\log X$ is infinite, the monomials
$e^{-\lambda X}X^{-r}=\exp(-\lambda X-r\log X)$ are exact Conway monomials,
ordered lexicographically in $(\lambda,r)$.

\begin{lemma}[Slow coefficients \src{03}]\label{gz:lem:sloweval}
For arbitrary $f_n\in F^{\rm Hahn}_X$, $\Sh_nf_ne^{-X\log n}$ is strongly
summable, and it realizes an injective homomorphism from the arithmetic
algebra over $F^{\rm Hahn}_X$. Its block at first-scale exponent $\log n$
is exactly $f_n$.
\end{lemma}
\begin{proof}
The support is $\bigcup_n\{\log n\}\times\supp f_n$ in the lexicographic
order. The first coordinates increase and each fibre is well ordered, so the
union is well ordered, and different fibres are disjoint. Products use the
finitely many factorizations of $n$.
\end{proof}

\begin{theorem}[Over the slow Hahn field \src{03}]\label{gz:thm:indep-slow}
For a positive infinite monomial $X$, the shift family of
\cref{gz:thm:indep-monomial} is independent over
$E_X=F^{\rm Hahn}_X(e^{-\lambda X}:\lambda\in\R)$, the generated field (not a
further Hahn completion).
\end{theorem}
\begin{proof}
Repeat the previous proof with $F^{\rm Hahn}_X$ in place of $\C$. Distinct
$\lambda$ give monomials that are linearly independent over $F^{\rm Hahn}_X$.
\Cref{gz:lem:sloweval} gives injectivity and separates cosets by their first
coordinate. Finish with \cref{gz:lem:primeindependence} for $F=F^{\rm Hahn}_X$.
\end{proof}

\begin{theorem}[Over a field containing the Gamma ray \src{03}]\label{gz:thm:gamma-zeta-independence}
Let $T_X=e^{X\log X}$ and $E_X^\Gamma=E_X(T_X)$. The shift family is
algebraically independent over $E_X^\Gamma$. This field contains
$\GSt(X+b)=\Gamma_\chi(X+b)$ for every ordinary real $b$ and every phase
$\chi$.
\end{theorem}
\begin{proof}
The field generated by $E_X$ and the jets lies in the Hahn field on the
two-scale group, and $X\log X\gg X\gg\log X$. So integer powers of $T_X$ lie
in distinct first-coordinate fibres of a three-scale workspace. Hence $T_X$
is transcendental over that field, and a relation $\sum_jT_X^jP_j(Y)=0$
splits into $P_j(Y)=0$, which contradicts \cref{gz:thm:indep-slow}. By
\cref{gz:thm:translation},
$\GSt(X+b)=\sqrt{2\pi}\,T_Xe^{-X}X^{b-1/2}\Ex(\SSt(X)+Q_b(X))$. The last factor
lies in $\C[[X^{-1}]]\subset F^{\rm Hahn}_X$, and the value is real and
phase-free.
\end{proof}
The Gamma values are not claimed to be mutually independent: the recurrence
relates them. Coefficients from all of $\SC$ are not allowed, since every
element is algebraic over a field that contains it \src{03}. Example: no
algebraic equation for $\zD(\omega)$ has nonzero coefficients in the field
generated by $\C((\omega^{-\R}))$, all $e^{-\lambda\omega}$, and all
$\GSt(\omega+b)$ with $b\in\R$.

\subsection{How the two independence theorems compare}\label{gz:sec:indep-compare}
The following table compares the two theorems.
\begin{center}\small
\begin{tabular}{@{}P{0.13\textwidth}P{0.40\textwidth}P{0.40\textwidth}@{}}
\toprule
 & 02 (\cref{gz:sec:dilation,gz:sec:leading}) & 03 (\cref{gz:sec:shift})\\
\midrule
Point & any positive infinite real $x$; complex $s\in\Tp$ for the jets alone & positive $X\in\Pi$; a monomial for the Hahn fields\\
Family & dilations $\zD^{(j)}(kx)$, $\zD^{(j)}(qx)$, with $x,\LSt(x),\GSt(x)$ & real shifts $\zD^{(k)}(X+a)$\\
Base field & $\C$; $\C(x)$; $\mathcal F_x^{\rm Dir}$ for $\LSt,\GSt$ & $M_X$; $E_X$; $E_X^\Gamma\ni\GSt(X+b)$\\
Gamma & $\LSt(x),\GSt(x)$ independent over $\mathcal F_x^{\rm Dir}$ (all Dirichlet values and $x$) & Gamma values in the base field; no mutual independence\\
\bottomrule
\end{tabular}
\end{center}
Neither theorem implies the other. For a positive infinite monomial $X$,
the overlap $a=0$ (03, over $E_X\supseteq\C(X)$) and $k=1$ (02) gives
independence of $\{\zD^{(j)}(X)\}$ over $\C(X)$; 03's field $M_X$ does not
contain $X$. 03's base field $E_X^\Gamma$ contains $X$, all real powers of
$X$ and $\GSt(X)$, but not $\log X$. Indeed $E_X^\Gamma$ lies in the Hahn
field whose monomials are $\exp(jX\log X-\lambda X-r\log X)$ ($j\in\Z$,
$\lambda,r\in\R$), and the leading monomial $\exp(\pp(\log\log X))$ of
$\log X$ is not of this form, since $0<\pp(\log\log X)\prec\log X$. So
$E_X^\Gamma$ does not contain $\LSt(X)$ either: that would give
$\log X=(\LSt(X)+X-\tfrac12\log2\pi-\SSt(X))/(X-\tfrac12)\in E_X^\Gamma$.
02 covers dilations and nonmonomial $x$, but has no Hahn
coefficients and no shifts. 02's dilation family, indicators of $k$th powers
times log powers, is not of Pong's form $T^\alpha\partial^jf$, and whether
another classical result covers dilations was not searched. A combined
dilation-and-shift statement over 03's field is proved by neither source and
is open (\cref{gz:q:combined}). 01's real inverse $s_0(u)$ involves
$\log u$ and $u^{\log_2(3/2)}$, which is consistent with the independence
of $x$ and $\zD(x)$.

\section{The Stirling logarithm at infinity}\label{gz:sec:stirling}

\subsection{Definition}
On the Stirling half-plane $\Sinf=\{\Re z>0,\ |z|\text{ infinite}\}$ put
\begin{equation}\label{gz:eq:LSt}
\LSt(z)=\Bigl(z-\frac12\Bigr)\Log z-z+\frac12\log2\pi+\SSt(z),\qquad
\SSt(z)=\Sh_{k\ge1}\frac{B_{2k}}{2k(2k-1)}z^{1-2k},
\end{equation}
with the principal finite angle. The series is strongly summable because
$z^{-1}\in\mm$, whatever the factorial growth of the Bernoulli numbers
\src{01--06}. At real $x>\R$ put $\GSt(x)=\exp_{\No}\LSt(x)>0$. With $\Gp$ on
positive finite arguments this is the positive-real baseline on $\No_{>0}$.
It satisfies $\Gamma(x+1)=x\Gamma(x)$ at every $x\in\No_{>0}$, by
\eqref{gz:eq:finiteref} at finite $x$ and by applying $\exp_{\No}$ to
\cref{gz:thm:stirling} at infinite $x$ \src{05}.
On real $x>\R$, $\LSt$ is the gamma-functions report's $L_0$
(\rl{eq:L0-infinite}); on $\{x>\R,\ y\in\OO_\R\}$ it is its $\widetilde L_0$
(\rl{eq:complex-infinite}, \rl{eq:log-infinite-complex}). The strong
interpretation of Stirling's series at infinite arguments is part of the
Costin--Ehrlich extension framework \cite{CE}, which treats log-Gamma
explicitly, and Berarducci and Mantova treat transseries as surreal functions \cite{BM}. The classical series is an asymptotic expansion, valid in
sectors $|\arg z|\le\pi-\delta$ \cite{DLMF511}. On $\Sinf$, \eqref{gz:eq:LSt}
is a \emph{prescription}. It also covers points such as
$\eps+iT$ with $\eps\in\mm_{>0}$ and $T>\R$ (a doubt raised in the review
of 06). Nothing ties
its values there to the analytic behaviour of a Gamma function; 05 calls
it \emph{bare Stirling} for this reason.

\begin{proposition}[Recurrence, Gauss and formal uniqueness \src{01, 02, 03, 05, 06}]\label{gz:thm:stirling}
On $\Sinf$, for ordinary $m\ge1$,
\[
\LSt(z+1)-\LSt(z)=\Log z,\qquad
\sum_{r=0}^{m-1}\LSt\Bigl(z+\frac rm\Bigr)=\LSt(mz)+\frac{m-1}2\log2\pi+\Bigl(\frac12-mz\Bigr)\log m .
\]
Among expressions $(z-\tfrac12)\Log z-z+c+\sum_{k\ge1}a_kz^{-k}$ with
$a_k\in\C$, the recurrence determines every $a_k$ once $c$ is fixed.
\end{proposition}
\begin{proof}
Work formally with $D=d/dZ$ and $u/(e^u-1)=\sum B_nu^n/n!$. Then
$\frac D{e^D-1}\log Z=\log Z-\frac1{2Z}-\sum_{k\ge1}\frac{B_{2k}}{2k}Z^{-2k}=\LSt'(Z)$,
and multiplying by $e^D-1$, the formal shift, gives
$\LSt'(Z+1)-\LSt'(Z)=1/Z$. So $\LSt(Z+1)-\LSt(Z)-\log Z$ has zero derivative.
Its expansion has zero constant term, so it vanishes. For Gauss, the
difference $R_m$ of the two sides satisfies $R_m(Z+1/m)=R_m(Z)$ by the
recurrence. Expanding $\Log(Z+\frac rm)=\Log Z+\Lgm(1+\frac r{mZ})$, the
$\Log Z$ terms agree because $\sum_{r<m}(Z+\frac rm-\frac12)=mZ-\frac12$,
the $\log m$ terms cancel, and the constant terms are $\frac m2\log2\pi$ on
both sides, since the constant $\frac rm$ of
$(Z+\frac rm-\frac12)\Lgm(1+\frac r{mZ})$ cancels the $-\frac rm$ of
$-(Z+\frac rm)$. So $R_m$ is a series in negative powers of $Z$
\merge. A nonzero such series
with leading term $cZ^{-d}$ changes by $-(cd/m)Z^{-d-1}\ne0$ under the shift,
so $R_m=0$. Uniqueness is the same argument: a nonzero difference
$\sum a_kz^{-k}$ with first term $a_kz^{-k}$ has shift difference
$-ka_kz^{-k-1}+\cdots\ne0$. Substituting $Z=z$ is legitimate by
\cref{gz:lem:substitution}, using $\Log(z+a)=\Log z+\Lgm(1+a/z)$ on $\Sinf$
for finite $a$ with $z+a\in\Sinf$, and $\Log(mz)=\Log z+\log m$.
\end{proof}
These identities duplicate the gamma-functions report's \rl{lem:formal-shift},
\rl{eq:baseline-recurrence}, \rl{lem:periodic} and \rl{prop:gauss0}, which
prove them on real arguments by coefficient matching and a periodicity
lemma. The proof above is the Bernoulli-operator route of 02, 03 and 05.
01, 04 and 06 instead derive the recurrence (01 also Gauss) from the
Hurwitz shift and distribution identities; that route is
\cref{gz:rem:secondroute} \src{01, 04, 06}. The uniqueness holds only among normalized formal series.
It does not exclude corrections smaller than every inverse power
(\cref{gz:sec:nonunique}) \src{01, 05}.

The formal logarithmic derivative is
\begin{equation}\label{gz:eq:psiSt}
\psiSt(z)=\Log z-\frac1{2z}-\Sh_{k\ge1}\frac{B_{2k}}{2k}z^{-2k}.
\end{equation}
At positive infinite real $A$, the raw classical formula
$\psi'(A)=\sum_{n\ge0}(A+n)^{-2}$ must not be used. Every summand has the
leading monomial of $A^{-2}$ with the same coefficient, so the family is not
strongly summable. The derivative of \eqref{gz:eq:psiSt} is the correct
object \src{01, 02, 03, 04}. The gamma-functions report records the same
guardrail in its unlabelled warning ``An illegitimate substitution into an
ordinary sum''.

\subsection{Translation and the ratio cocycle}
\begin{theorem}[Bernoulli translation formula \src{02, 03, 05}]\label{gz:thm:translation}
Let $Z\in\SC$ with $Z^{-1}\in\mm$, let $a\in\OO$, and let $\Log Z$ be a
logarithm of $Z$ with finite imaginary part. Put
$\mathcal L(Z+a)=(Z+a-\tfrac12)(\Log Z+\Lgm(1+a/Z))-(Z+a)+\tfrac12\log2\pi+\SSt(Z+a)$.
Then
\begin{equation}\label{gz:eq:shiftedStirling}
\mathcal L(Z+a)=\Bigl(Z+a-\frac12\Bigr)\Log Z-Z+\frac12\log2\pi+\Sh_{n\ge1}\frac{(-1)^{n+1}B_{n+1}(a)}{n(n+1)}Z^{-n}.
\end{equation}
If $Z,Z+a\in\Sinf$ with the principal branches, then $\mathcal L=\LSt$, and
\begin{align*}
\LSt(Z+a)-\LSt(Z)-a\Log Z=Q_a(Z)&:=\Sh_{n\ge1}\frac{(-1)^{n+1}(B_{n+1}(a)-B_{n+1})}{n(n+1)}Z^{-n}\\
&=\frac{a(a-1)}{2Z}-\frac{a(a-1)(2a-1)}{12Z^2}+\cdots .
\end{align*}
Hence, for every phase, $\Gamma_\chi(Z+a)/(\Gamma_\chi(Z)E_\chi(a\Log Z))=\Ex(Q_a(Z))$,
independent of $\chi$. At $a=1$, $Q_1=0$ gives the recurrence again.
\end{theorem}
\begin{proof}
Work formally in $u=Z^{-1}$, with $a$ an indeterminate and $\Log Z$ a
formal symbol of derivative $1/Z$. Using $B_1(a)=a-\tfrac12$, the
$Z$-derivative of the right side is $\Log Z+\sum_{n\ge1}(-1)^{n+1}B_n(a)Z^{-n}/n$.
The $Z$-derivative of the left side is $\psiSt(Z+a)$ expanded in $a/Z$. By
the binomial theorem and $B_n(a)=\sum_j\binom njB_ja^{n-j}$ its coefficient
of $Z^{-n}$ is the same, which is 03's computation of $L'(z+a)$. So the two
sides differ by a constant in $Z$. Both have constant term
$\tfrac12\log2\pi$: on the left, the term $Z\cdot a/Z$ from
$(Z+a-\tfrac12)\Lgm(1+a/Z)$ cancels $-a$. So \eqref{gz:eq:shiftedStirling}
holds in $\Q[a]((u))$ together with the log terms. To evaluate, write $a=c+\delta$
with $c\in\C$ and $\delta\in\mm$. Every coefficient is a polynomial in
$\delta$, so the family is jointly strongly summable in $u$ and $\delta$
(\cref{gz:lem:substitution}). This joint evaluation is the step that 02's
(8.6) asserts without detail (error 11). On $\Sinf$ the principal branches
satisfy $\Log(Z+a)=\Log Z+\Lgm(1+a/Z)$, since both sides are logarithms of
$Z+a$ with imaginary parts differing by an infinitesimal. Apply $E_\chi$ and
its local law on the infinitesimal $Q_a(Z)$. Finally $B_{n+1}(1)=B_{n+1}$
for $n\ge1$.
\end{proof}
03 proves this on $\Tp$, 02 for real $x>\R$ and $h\in\OO$, and 05 as a
formal identity at infinite $|Z|$ (its eq.\ 6.11, the surreal form of DLMF
5.11.8). The evaluation at $Z=\pm iT$, used in \cref{gz:thm:strip}, needs
only the first statement.

\begin{theorem}[Ratio cocycle \src{05}]\label{gz:thm:cocycle}
For $|Z|$ infinite and finite $a,b,c$, put
$Q_{a,b}(Z)=\Sh_{n\ge1}\frac{(-1)^{n+1}(B_{n+1}(a)-B_{n+1}(b))}{n(n+1)}Z^{-n}$
and $R_{a,b}=\Ex(Q_{a,b})$. Then $R_{a,b}$ is well defined and nonzero, and
\begin{gather*}
R_{a,b}R_{b,c}=R_{a,c},\quad R_{a,a}=1,\quad R_{1,0}=1,\quad
\overline{R_{a,b}(Z)}=R_{\bar a,\bar b}(\bar Z),\\
\frac{R_{a,b}(Z+1)}{R_{a,b}(Z)}=\frac{1+a/Z}{1+b/Z}(1+1/Z)^{-(a-b)},\qquad
R_{a,b}(Z)=1+\frac{(a-b)(a+b-1)}{2Z}+\bigO(Z^{-2}).
\end{gather*}
The powers are near-one powers. For example
$R_{1/2,0}(Z)=1-\frac1{8Z}+\frac1{128Z^2}+\frac5{1024Z^3}-\frac{21}{32768Z^4}+\bigO(Z^{-5})$.
\end{theorem}
\begin{proof}
Expand the Bernoulli polynomials at the standard parts of $a,b$ and evaluate
jointly in $Z^{-1}$, $a-\st a$, $b-\st b$. The cocycle follows by cancelling
the middle polynomial, $R_{1,0}=1$ from $B_k(1)=B_k$ ($k\ge2$), and
conjugation from the real coefficients. For the shift, subtract the two
instances of \eqref{gz:eq:shiftedStirling} and use the recurrence:
$Q_{a,b}(Z+1)-Q_{a,b}(Z)=\Lgm(1+a/Z)-\Lgm(1+b/Z)-(a-b)\Lgm(1+1/Z)$. The leading
term uses $B_2(a)=a^2-a+\tfrac16$.
\end{proof}
$Q_{a,b}=Q_a-Q_b$, where $Q_a$ is 03's $R_a$, and at $b=0$ with $a=h$ it is
the correction in 02's (8.6). The leading term $(a-b)(a+b-1)/(2Z)$ matches
03's $a(a-1)/(2z)$. The coefficient $-21/32768$ is an addition to 05's
expansion, confirmed numerically in the review of 05.
\begin{remark}[Reading $R_{a,b}$ as a Gamma ratio \src{05}]\label{gz:rem:cocycle-gloss}
05 glosses $R_{1/2,0}$ as the formal normalization of
$\Gamma(Z+\tfrac12)/(Z^{1/2}\Gamma(Z))$, and the same reading makes
$R_{a,b}$ the formal normalization of $\Gamma(Z+a)/(Z^{a-b}\Gamma(Z+b))$. The identities hold for every infinite
$Z$, including negative real $Z$, but the gloss is meaningful only away from
the negative real direction, that is, in sectors $|\arg Z|\le\pi-\delta$.
On the negative real axis the classical ratio carries a non-flat quotient
of sines coming from reflection (error 34). $R_{a,b}$ is a formal, phase-free
object, and no claim is made that exponentially small corrections vanish.
\end{remark}

\subsection{The phase estimate and the finite-phase locus}\label{gz:sec:phasedomain}
\begin{proposition}[Phase estimate \src{01, 02, 03, 04, 05, 06}]\label{gz:prop:phase}
Let $x>\R$ and let $y$ be finite. Then
\[
\Im\LSt(x+iy)=y\log x-\frac y{2x}+\Bigl(\frac{y^2}6-\frac1{12}\Bigr)\frac y{x^2}+\bigO\Bigl(\frac y{x^3}\Bigr),
\]
so $\Im\LSt(x+iy)-y\log x\in\mm_\R$. Consequently $\Efp(\LSt(z))$ is defined
exactly on $\Om$. There $\Gamma_{\rm fp}=\Efp\circ\LSt$ is zero-free and
satisfies the recurrence and every finite Gauss formula. A nonzero ordinary
imaginary displacement $x+iy$ ($y\in\R^\times$) always leaves $\Om$.
\end{proposition}
\begin{proof}
Expand $\Log(x+iy)=\log x+\Lgm(1+iy/x)$ and $(x+iy)^{-1}$ in $y/x$. The term
$y$ from $(x-\tfrac12)\Im\Lgm(1+iy/x)$ cancels $-\Im z=-y$, leaving
$y\log x-y/(2x)+y^3/(6x^2)+\cdots$, and $\Im\SSt(z)=-y/(12x^2)+\cdots$.
Adding an infinitesimal does not change finiteness. $\Om$ is preserved by
$z\mapsto z+1$, $z\mapsto z+j/q$ and $z\mapsto qz$, since
$\log(x+j/q)-\log x\in\mm$ and $\log(qx)=\log x+\log q$. Apply $\Efp$ to
\cref{gz:thm:stirling}. For ordinary $y\ne0$, $y\log x$ is infinite.
\end{proof}
This is the gamma-functions report's \rl{lem:phase} and \rl{thm:phase-domain}.
Those results are stronger: they cover every gauge $a$, the finite-real-part
part of the tube, and fine holomorphy and fine-openness. They are cited
here, not claimed. The expansion is printed with the corrected remainder:
01's proof of its Theorem 6.2 omits the terms of order $y^3/x^2$ and
$y/x^2$ (error 4). 05 writes $\Im\LSt=y\log x(1+\eta)$, a weaker consistent
form, and ``usually'' where ``always'' is correct (error 30). Examples
\src{01, 02}, as in the gamma-functions report's example after
\rl{thm:phase-domain}: $x+i/\log x\in\Om$ with
$\Gamma_{\rm fp}(x+i/\log x)/\GSt(x)=e^i(1+\eta)$, while $x+i\notin\Om$.
Purely real positive infinite arguments always lie in $\Om$. The locus is
not a maximal-domain or impossibility statement \src{01, 02}. The next
proposition makes that precise.

\begin{proposition}[$\Om$ is the locus of phase independence \merge]\label{gz:prop:phase-locus}
Let $z=x+iy$ with $x>\R$ and $y$ arbitrary. All phase exponentials $E_\chi$
give the same value $E_\chi(\LSt(z))$ if and only if $\Im\LSt(z)$ is finite,
if and only if $z\in\Om$. If $p=\pp(\Im\LSt(z))\ne0$ and $M$ is the leading
monomial of $p$, the phases $\chi_{\lambda,M}$ give
$E_{\chi_{\lambda,M}}(\LSt(z))=e^{i\lambda c_M(p)}\,\Gamma_{\Exp}(z)$ for every
$\lambda\in\R$, so $z$ carries a whole circle of values.
\end{proposition}
\begin{proof}
If $\Im\LSt(z)$ is finite, every $E_\chi$ agrees with $\Efp$ there. If $y$ is
finite, $\Im\LSt(z)$ is finite exactly on $\Om$ by \cref{gz:prop:phase}.
Suppose $y$ is infinite, say positive. Then
$\Im\LSt(z)=(x-\tfrac12)\Arg z+y(\log|z|-1)+\Im\SSt(z)$. Here $\Arg z>0$ is
finite, $\log|z|-1\ge\log y-1$ is positive infinite, and $\Im\SSt(z)$ is
infinitesimal. So $\Im\LSt(z)$ is positive infinite, and similarly for
negative $y$. Hence $\Im\LSt(z)$ is finite exactly on $\Om$. Finally, if
$p\ne0$ then $c_M(p)\ne0$, and
$\Psi_{\chi_{\lambda,M}}(\Im\LSt)=e^{i\lambda c_M(p)}\cis(\fin\Im\LSt)$.
\end{proof}
This turns ``$\Efp\circ\LSt$ is defined exactly on $\Om$'', which is 01's
Theorem 6.2 and the gamma-functions report's \rl{thm:phase-domain}, into a
statement that stays true when $\Exp$ is available everywhere. 01's ``defined
if and only if'' is correct for $\Efp$. It does not mean that Gamma is
undefined elsewhere (error 5).

\subsection{Gamma with a named phase}\label{gz:sec:gammachi}
\begin{proposition}[Phase Gamma on $\Sinf$ \src{01, 02, 03, 06}]\label{gz:prop:gammachi}
Let $E:\SC\to\SC^\times$ be any additive-to-multiplicative homomorphism
extending $\Efp$. Then $\Gamma_{E,\rm St}=E\circ\LSt$ is nonzero on $\Sinf$ and
satisfies $\Gamma_{E,\rm St}(z+1)=z\,\Gamma_{E,\rm St}(z)$ \src{06}. For a
phase $\chi$, $\Gamma_\chi=E_\chi\circ\LSt$ moreover commutes with
conjugation and satisfies every finite Gauss formula on $\Sinf$, with
$m^{1/2-mz}:=E_\chi((\tfrac12-mz)\log m)$. It agrees with $\GSt$ on real
$x>\R$ and with $\Gamma_{\rm fp}$ on $\Om$. It has the local law
$\Gamma_\chi(z+h)=\Gamma_\chi(z)\Ex(\LSt(z+h)-\LSt(z))$ for increments with
$h/z$ and $h(1+|\Log z|)$ infinitesimal.
\end{proposition}
\begin{proof}
$E(\Log z)=z$ because $\Arg z$ is finite. Apply $E$ to
\cref{gz:thm:stirling}. A homomorphism into $\SC^\times$ does not vanish. The
local law expands $\LSt(z+h)$ jointly in $z^{-1}$ and $h/z$; the linear
coefficient contains $\Log z$.
\end{proof}
06 does not assume the modulus law $|E|=\exp\Re$; the other clauses use it
through \cref{gz:prop:Echi}. 02 states the theorem on $\Tp$ with $E_\theta$,
01 on the wing with $E_\lambda$, and 03 on $\Tp$ with every phase.
$\Gamma_\chi$, and $\Gamma_{E,\rm St}$ for every $E$, satisfies the
recurrence on $\Sinf$, but \emph{not} reflection at any point of the
vertical strip (\cref{gz:thm:strip}(e)). 06 does not claim reflection
there, so this is a missing remark in 06, not an error (error 48).

Each $\Gamma_\chi$ can be joined with $\Gp$ on $\OO$, which is disjoint from
$\Sinf$. The result is a union of prescriptions on disjoint domains: it
agrees with the ordinary $\Gamma$ on $\C$ because $\Gp$ does, and no
identity linking the two pieces is asserted \src{01}.

\begin{example}[Phase dependence \src{01, 02, 03}]\label{gz:ex:phasedep}
\leavevmode
\begin{enumerate}[label=(\alph*)]
\item \src{01} At $x=\exp_{\No}(\omega)$, $\Im\LSt(x+i)=\omega+\delta$ with
 $\delta\in\mm$. So $\Gamma_{\chi_\lambda}(x+i)/\Gamma_{\Exp}(x+i)=e^{i\lambda}$,
 and $\lambda=0,\pi$ give opposite values at one point while agreeing on the
 real axis.
\item \src{02} For every positive infinite real $x$ and every ordinary
 $\vartheta$, there is an explicit $\theta=\theta_x$ (a multiple of the
 coefficient of one monomial of $\pp(\log x)$) with
 $\Gamma_{\chi_\theta}(x+i)/\Gamma_{\Exp}(x+i)=e^{i\vartheta}$. The functional
 depends on $x$ (error 12).
\item \src{03} For a positive infinite monomial $X$, $\log X\in\Pi$, so
 $\Exp(i\log X)=1$ and
 $\Gamma_{\Exp}(X+i)/\Gamma_{\Exp}(X)=\Ex(Q_i(X))=1-\frac{1+i}{2X}+\bigO(X^{-2})$.
 At the smaller displacement $i\vartheta/\log X$ the phase is finite, and
 $\Gamma_{\Exp}(X+i\vartheta/\log X)/\GSt(X)=e^{i\vartheta}(1+\eps)$.
\end{enumerate}
\end{example}
For $z\in\Sinf$ the phase-free factor is
\[
\Ex(\SSt(z))=1+\frac1{12z}+\frac1{288z^2}-\frac{139}{51840z^3}-\frac{571}{2488320z^4}+\frac{163879}{209018880z^5}+\bigO(z^{-6}),
\]
and
$\Gamma_{\Exp}(\omega)=\GSt(\omega)=\sqrt{2\pi}\,\omega^{\omega-1/2}e^{-\omega}\Ex\bigl(\frac1{12\omega}-\frac1{360\omega^3}+\frac1{1260\omega^5}-\frac1{1680\omega^7}+\cdots\bigr)$,
where $\omega^{\omega-1/2}=\exp_{\No}((\omega-\tfrac12)\log\omega)$ involves no
phase \src{01, 02, 03, 04}.

\subsection{Three mechanisms of non-uniqueness}\label{gz:sec:nonunique}
\begin{enumerate}[label=(\arabic*)]
\item \emph{Phase choice off $\Om$} \src{01, 02, 03, 06}: the functions
 $E_\chi\circ\LSt$ for different $\chi$
 (\cref{gz:prop:phase-locus,gz:ex:phasedep}). No real value changes.
\item \emph{Flat periodic gauges on the positive reals} \src{04, 05}: 05's
 $\GSt\cdot e^{h^{05}_\lambda}$ with $h^{05}_\lambda(x)=\lambda e^{-\pp(x)}$
 at infinite $x$ and $0$ at finite $x$, and 04's $\GSt\cdot\exp(\pi e^{-\pp(A)})$
 (\cref{gz:prop:flathurwitz}). Both keep the recurrence and every finite
 value, and the corrections are smaller than every inverse power, since
 $x=\pp(x)+\bigO(1)$ and $e^{-P}P^N\in\mm$. \emph{Both break every Gauss
 formula at infinite arguments} \merge: for $h=\lambda e^{-\pp}$ with
 $\lambda\ne0$ and ordinary $q\ge2$,
 $\sum_{r<q}h(x+r/q)=q\lambda e^{-\pp(x)}\ne\lambda e^{-q\pp(x)}=h(qx)$,
 because $e^{(q-1)\pp(x)}\ne1/q$. Neither source claims Gauss, so this is a
 missing statement, not an error (error 23).
\item \emph{The convex gauge family} of the gamma-functions report
 (\rl{cor:headline}, \rl{thm:main}): $h_a(x)=a(M(x))\pp(x)$, for example
 $a_\lambda(m)=\lambda e^{-m}$, which is 06's gauge. It keeps the recurrence,
 every Gauss formula, strict log-convexity and all Stirling estimates.
 06 checks directly that it keeps the recurrence and every finite Gauss
 formula, since $M(x+1)=M(x)$, $h(qx)=qh(x)$ and $h(x+j/q)=h(x)$ \src{06}.
 Strict log-convexity, the Stirling estimates and the sharp
 classification belong to that report and are not reproved.
\end{enumerate}
\emph{Reconciliation.} Gamma at infinity is non-unique in three independent
ways. The first changes no real value. The last two change no finite value
and no asymptotic expansion, and only the third keeps Gauss multiplication.
A small logarithmic difference need not be a small absolute difference of
Gamma values (\rl{warn:absolute-difference}). The member $a=0$, the
baseline of every source, is the one singled out by the scalar chain rule
for the Berarducci--Mantova derivation (\rl{thm:BM-all}). The sources fix
it by prescription, not by that derivation \src{03}.

\section{Hurwitz zeta at an infinite parameter}\label{gz:sec:hurwitz}

\subsection{Definition and domain}
Let $A>\R$ be real and $s\in\OO$. With $(s)_j=s(s+1)\cdots(s+j-1)$ and a
named phase for $A^w=E_\chi(w\log A)$, put
\begin{equation}\label{gz:eq:zH}
\zH(s,A)=\frac{A^{1-s}}{s-1}+\frac{A^{-s}}2+\Sh_{k\ge1}\frac{B_{2k}}{(2k)!}(s)_{2k-1}A^{1-s-2k},\qquad
\zHh(s,A)=A^{s-1}\zH(s,A).
\end{equation}
This takes the classical large-parameter expansion (DLMF 25.11.43
\cite{DLMF2511}) as an exact definition \src{01, 02, 03, 04, 06}. Since
$A^{1-s}$ factors out, $\zHh$ is phase-free, and $\zH(s,A)$ is phase-free
exactly when $(\Im s)\log A$ is finite, that is, on $\SA$ \src{01, 02}.
For finite $s$ with $(\Im s)\log A$ infinite its value depends on the phase
\src{03}, and every identity below holds for every phase.
\begin{proposition}[Existence \src{01, 02, 03}]\label{gz:prop:hurwitz-exists}
For every finite $s$ the family in \eqref{gz:eq:zH} is strongly summable.
\end{proposition}
\begin{proof}
Write $s=c+\eps$. Each $(s)_{2k-1}$ is a polynomial in $\eps$ with complex
coefficients, and the $k$th term carries $A^{-2k}$ after $A^{1-s}$ is
factored out. Apply \cref{gz:lem:substitution} to $A^{-1}$ and $\eps$
jointly.
\end{proof}
At $s\ne1$ the factor $1/(s-1)$ is an ordinary field inverse, even when $s$
is infinitely close to $1$; no separate interpretation is needed (01's
wording, error 6). \emph{Parameter germs.} At a finite centre $s_0$, write
$A^{-s_0-h}=A^{-s_0}\Ex(-h\log A)$. This is a strongly evaluable germ when
$h\log A\in\mm$, but not for every infinitesimal $h$ \src{02, 03, 04, 06}.
04 uses the germ at $0$ ($s\in\mm$, $s\log A\in\mm$) and separate
prescriptions at $s=-m$. 06 uses ordinary real $s$ with germs. 01 uses the
formal germ $\zH(S,A)\in\SC[[S]]$, evaluable when $S\log A\in\mm$.

\begin{warning}[The raw shifted sum \src{01, 02, 03, 04, 06}]\label{gz:warn:hurwitz-raw}
For real $A>\R$ and every $s\in\SA$, in particular every ordinary real
$\sigma$ including $\sigma=0$ and the negative integers \src{02}, the family
$((A+n)^{-s})_{n\ge0}$ is not strongly summable. Every member is
$A^{-s}(1+n/A)^{-s}=A^{-s}(1+\eps_n)$ with $\eps_n\in\mm$, so all members
have the same leading monomial and coefficient. So \eqref{gz:eq:zH} is not
a rearrangement of a strong sum of $(A+n)^{-s}$. Substituting $A=1$ into an infinite-parameter
definition does not give a global Riemann zeta \src{06}. No relation
$\zeta(s)=\sum_{n<A}n^{-s}+\zH(s,A)$ with omnific $A$ is used anywhere
\src{03}.
\end{warning}

\begin{theorem}[Normalized unit \src{03}]\label{gz:thm:hurwitz-unit}
For every finite $s$, $(s-1)\zHh(s,A)=1+\eps(s,A)$ with $\eps(s,A)\in\mm$. So
$\zH(\cdot,A)$ is locally meromorphic in finite $s$, has a simple pole with
residue $1$ at $s=1$ and no other poles, and has no zeros at finite $s$. The
normalized factor and the zero-freeness are independent of the phase.
\end{theorem}
\begin{proof}
$(s-1)\zHh$ has constant term $1$, and every other term carries a positive
power of $A^{-1}$ with a coefficient polynomial in $s$, so their joint sum
is infinitesimal. For $s\ne1$ both $(s-1)^{-1}$ and $A^{1-s}$ are nonzero.
At $s_0=1$, $A^{1-s}=1+\bigO(h\log A)$ and the principal part of $\zHh$ is
$1/h$.
\end{proof}
Consistently, $\zH(-m,A)=-B_{m+1}(A)/(m+1)\ne0$ (\cref{gz:thm:lerch}).

\subsection{Identities}
\begin{theorem}[Hurwitz identities \src{01, 02, 03, 04, 06}]\label{gz:thm:hurwitz-identities}
For real $A>\R$, finite $s$, ordinary $q,m\ge1$, and every phase,
meromorphically in $s$:
\begin{gather*}
\zH(s,A)-\zH(s,A+1)=A^{-s},\qquad\partial_A\zH(s,A)=-s\,\zH(s+1,A),\\
\sum_{r=0}^{m-1}\zH\Bigl(s,\frac{A+r}m\Bigr)=m^s\zH(s,A),\qquad
\sum_{j=0}^{q-1}\zH\Bigl(s,A+\frac jq\Bigr)=q^s\zH(s,qA),
\end{gather*}
and $\Res_{s=1}\zH(s,A)=1$. The two distribution identities are the same
after $A\mapsto A/q$ \src{01, 03}.
\end{theorem}
\begin{proof}
Work over $\C(s)$ in $A^{1-s}\C(s)[[A^{-1}]]$ with $D=d/dA$ and
$D^{-1}A^{-s}=A^{1-s}/(1-s)$, with no integration constant. The Bernoulli
identity $(1-e^D)^{-1}=-D^{-1}+\tfrac12-\sum_{k\ge1}\frac{B_{2k}}{(2k)!}D^{2k-1}$
applied to $A^{-s}$ gives exactly \eqref{gz:eq:zH}, since
$D^{2k-1}A^{-s}=-(s)_{2k-1}A^{-s-2k+1}$. Multiplying by $1-e^D$ gives the
shift. The operator $e^D$ acts as $A\mapsto A+1$ through the binomial series
$(A+1)^w=A^w\sum_i\binom wiA^{-i}$. So the coefficient of $A^{1-s-n}$ in
$e^D\zH$ is the finite sum $\sum_{j+i=n}c_j\binom{1-s-j}i$, where $c_j$ is
the coefficient of $A^{1-s-j}$ in $\zH$. The regrouping is therefore
coefficientwise finite. This is the detail that 04's proof of its Theorem
8.1 sketches (error 25). Since $\zH(s,A)=(1-e^D)^{-1}A^{-s}$ and the operator
commutes with $D$, $\partial_A\zH(s,A)=(1-e^D)^{-1}(-sA^{-s-1})=-s\,\zH(s+1,A)$;
termwise this is $(s)_{2k-1}(s+2k-1)=s(s+1)_{2k-1}$. For distribution, the
difference
$W(A)=\sum_{r<m}\zH(s,(A+r)/m)-m^s\zH(s,A)$ is $1$-periodic by the shift
identity, and $m^{1-s}W\in A^{1-s}\C(s)[[A^{-1}]]$. A nonzero element of
this space with first term $c_j(s)A^{1-s-j}$ has forward difference
beginning $c_j(s)(1-s-j)A^{-s-j}\ne0$ over $\C(s)$, so $W=0$ \src{03}. To
specialize, write $s=c+\eps$. Affine changes of $A$ use
$\log(A+r)=\log A+\Lgm(1+r/A)$, and the leading factors are exponentiated by
the homomorphism $E_\chi$. Every coefficient identity is polynomial in
$\eps$, so strong evaluation in $A^{-1}$ and $\eps$ proves the displayed
identities. Poles and removable points are handled as germ identities. The
residue comes from the first term alone.
\end{proof}
\emph{Second route} \src{01}: the geometric operator identity
$\sum_{j<q}e^{jD/q}=(e^D-1)/(e^{D/q}-1)$ gives the distribution formula
directly, because the derivative with respect to $qA$ is $D/q$.

\begin{corollary}[Formal uniqueness \src{03}]\label{gz:cor:hurwitz-unique}
Over $\C(s)$, the shift equation $Z(s,A+1)-Z(s,A)=-A^{-s}$ has exactly one
solution in $A^{1-s}\C(s)[[A^{-1}]]$, namely \eqref{gz:eq:zH}.
\end{corollary}
The uniqueness is formal, before specialization. At $s=-m$, adding a
constant in $A$ preserves the specialized difference equation, and it is
coherence in $s$ that fixes these constants. The corollary does not
classify every function satisfying the difference equation, and flat
Hurwitz sectors are not excluded (\cref{gz:prop:flathurwitz}).

\begin{theorem}[Bernoulli values and the Lerch identity \src{01, 02, 03, 04, 06}]\label{gz:thm:lerch}
For real $A>\R$ and ordinary $m\ge0$, and independently of the phase:
\begin{gather*}
\zH(-m,A)=-\frac{B_{m+1}(A)}{m+1},\qquad
\zH(0,A)=\tfrac12-A,\qquad
\zH(-1,A)=-\tfrac{A^2}2+\tfrac A2-\tfrac1{12},\\
\zH(-2,A)=-\tfrac{A^3}3+\tfrac{A^2}2-\tfrac A6,\qquad
\zH(-3,A)=-\tfrac{A^4}4+\tfrac{A^3}2-\tfrac{A^2}4+\tfrac1{120},\\
\partial_s\zH(s,A)\big|_{s=0}=\LSt(A)-\tfrac12\log2\pi,\qquad
\GSt(A)=\sqrt{2\pi}\exp\bigl(\partial_s\zH(s,A)|_{s=0}\bigr),\\
\zH(1+h,A)=h^{-1}-\psiSt(A)+\bigO(h)\quad(h\log A\in\mm),\\
\psiSt^{(m)}(A)=(-1)^{m+1}m!\,\zH(m+1,A)\quad(m\ge1).
\end{gather*}
For example $\psiSt'(A)=A^{-1}+\tfrac12A^{-2}+\Sh_{j\ge1}B_{2j}A^{-2j-1}$.
\end{theorem}
\begin{proof}
At $s=-m$ the rising factorial $(s)_{2k-1}$ vanishes once $2k-1>m$, and
$(-m)_{2k-1}=-m!/(m-2k+1)!$ for $2k-1\le m$. With
$B_{m+1}(A)=\sum_j\binom{m+1}jB_jA^{m+1-j}$ and the vanishing of odd
Bernoulli numbers beyond $B_1$, this gives $-B_{m+1}(A)/(m+1)$. At $s=0$
the derivatives of the first two terms are $A\log A-A$ and $-\tfrac12\log A$.
Since $(s)_{2k-1}$ vanishes at $0$ with derivative $(2k-2)!$, the $k$th term
contributes $\frac{B_{2k}}{2k(2k-1)}A^{1-2k}$. At $s=1+h$,
$A^{-h}=1-h\log A+\bigO(h^2)$ and $(1)_{2k-1}=(2k-1)!$, which give the
constant $-\log A+\frac1{2A}+\Sh\frac{B_{2k}}{2k}A^{-2k}=-\psiSt(A)$.
Differentiating \eqref{gz:eq:psiSt} $m$ times and comparing with
$\zH(m+1,A)$ through $(m+1)_{2k-1}=(2k+m-1)!/m!$ gives the polygamma formula.
All centres are real, so only real exponentiation is used.
\end{proof}
Classically the Lerch identity is $\zeta'(0,A)=\log\Gamma(A)-\tfrac12\log2\pi$
(DLMF 25.11.18). Here it is verified for the normal forms, with no strong sum
over $(A+n)^{-s}$ (01). The Lerch identity is a normalization principle: the
baseline's logarithm comes from a coherent spectral family. It does not
remove flat deformations if unprescribed flat Hurwitz sectors are allowed
\src{02, 03}. The value $-1/12$ in $\zH(-1,A)$ is the coherence-selected
constant of a polynomial tail. It does not make $1+2+3+\cdots$ a strong sum,
since that family is not strongly summable \src{02, 03, 06}.

\begin{remark}[Second route to the Stirling identities \src{01, 04, 06}]\label{gz:rem:secondroute}
Differentiating the shift identity at $s=0$ and using the Lerch identity
gives $\LSt(A+1)-\LSt(A)=\log A$ (01 Theorem 6.1; 04 Theorem 8.1; 06
Proposition 6.1).
Differentiating the distribution identity gives Gauss (01 Theorem 6.1). The
identities are formal in inverse powers, so they extend to complex $z\in\Sinf$
on the principal branch. This route is independent of
\cref{gz:thm:stirling}'s proof.
\end{remark}

\begin{proposition}[A flat Hurwitz modification \src{04}]\label{gz:prop:flathurwitz}
Let $h(A)=\exp(-\pp(A))$ for positive infinite $A$ and $h(A)=0$ for finite
$A$. On the local $s$-germs at the ordinary integers, where values are
prescribed, put $\widetilde\zeta_{\mathrm H}(s,A)=\zH(s,A)+h(A)\sin^\#(\pi s)$.
Then $\widetilde\zeta_{\mathrm H}$ satisfies the same shift identity, and
the correction $h(A)\sin^\#(\pi s)$ vanishes at every ordinary integer $s$
and is smaller than every ordinary inverse power of $A$. But
$\partial_s\widetilde\zeta_{\mathrm H}|_{s=0}=\partial_s\zH|_{s=0}+\pi h(A)$,
so the corresponding Gamma is $\GSt\cdot\exp(\pi h(A))$, with the same
recurrence and the same finite values.
\end{proposition}
\begin{proof}
$\pp(A+1)=\pp(A)$, so $h$ is $1$-periodic. $\pp(A)$ differs from $A$ by a
finite amount and dominates every multiple of $\log A$, so $A^Nh(A)\in\mm$.
The sine vanishes at integers and has derivative $\pi$ at $0$.
\end{proof}
The modification uses only the integer germs and does not extend 04's germ
domain to the whole finite $s$-plane. By \cref{gz:sec:nonunique}(2) it
breaks every Gauss formula at infinite $A$. The Stirling--Hurwitz values are
thus a \emph{selected} normalization, and no uniqueness, convexity or Gauss
claim is made for the correction \src{04}.

\section{The right wing}\label{gz:sec:wing}

$\Tp$ has finite imaginary part. On the wing $\Wing=\{\Re s>\R\}$ the terms
$E_\chi(-s\log n)$ need a phase. \emph{Reconciliation.} $\zD$ is zero-free
on $\Tp$ in all six sources. 02's Warning 6.4, which confines zero-freeness
to $\Tp$, concerns the raw $\zD$ (error 13). Its own Proposition 11.3, 03's
Proposition 13.1, and the next proposition concern wing sums with a named
phase. 04 states the Dirichlet chart for finite imaginary part only and
omits the extension under $\Exp$ (error 24). 05 defines nothing at
infinite imaginary part, so no source asserts zero-freeness there without
naming a phase.

\begin{proposition}[Wing sums for ordinary-circle phases \src{02, 03}; general form \merge]\label{gz:prop:wing}
Let $\chi(\Pi)\subseteq S^1(\C)$, and let $s=x+iy\in\Wing$ with $Y=\pp(y)$ and
$f=\fin(y)$. Then $(E_\chi(-s\log n))_{n\ge1}$ is strongly summable, and
\[
\zeta^\Wing_\chi(s)=\Sh_nE_\chi(-s\log n)=\DD_{x+if}(a),\qquad a(n)=\chi(-Y\log n)\in S^1(\C),
\]
with $a$ completely multiplicative. So $\zeta^\Wing_\chi$ is a zero-free
Euler product, locally analytic, and agrees with $\zD$ on $\Tp$. For
$E_\theta$, $\zeta^\Wing_{\chi_\theta}(s)=\zD(\pi_\theta(s))$ with
$\pi_\theta(x+iy)=x+i(\fin y+\theta(\pp y))$ \src{02}. For $\Exp$,
$\zeta^\Wing_{\Exp}(s)=\zD(x+i\fin y)$, and this is invariant under
$s\mapsto s+iY'$ for every $Y'\in\Pi$ \src{03}.
\end{proposition}
\begin{proof}
$-y\log n=-Y\log n-f\log n$ with $Y\log n\in\Pi$ and $f\log n$ finite. So
$E_\chi(-s\log n)=\Efp(-(x+if)\log n)\,\chi(-Y\log n)$, and
$x+if\in\Tp$. \Cref{gz:thm:dirichlet} allows arbitrary complex coefficients.
Complete multiplicativity comes from the character property. Zero-freeness
and the Euler product are \cref{gz:thm:euler}(d). An infinitesimal increment
changes $f$ and $x$ but not $Y$, so the local germ is that of
\cref{gz:prop:Dtaylor}. For $\chi_\theta$, $a(n)=n^{-i\theta(Y)}$.
\end{proof}

\begin{corollary}[A continuum of wing values \src{02}]\label{gz:cor:continuum}
For real $x>\R$ and nonzero $Y\in\Pi$, the values $\zeta^\Wing_{\chi_\theta}(x+iY)$,
as $\theta$ ranges over explicit coefficient functionals, include the
continuum of pairwise distinct values $\zD(x+it)$, $t\in\R$.
\end{corollary}
\begin{proof}
Scale a coefficient functional so that $\theta(Y)=t$. If $\zD(x+it)=\zD(x+iu)$,
injectivity of $\DD_x$ gives $e^{-i(t-u)\log n}=1$ for all $n$. At $n=2,3$
this makes $\log2/\log3$ rational unless $t=u$.
\end{proof}

\begin{proposition}[Wing sums can fail \merge]\label{gz:prop:wingfail}
Assume choice for a Hamel basis of $\R$. There is a phase $\chi$, whose
character values are $\cis\delta$ with infinitesimal $\delta$ (values in
$1+\mm$, off $S^1(\C)$), for which $(E_\chi(-s\log n))_n$ is not
strongly summable at $s=\omega+i\omega\in\Wing$.
\end{proposition}
\begin{proof}
The $\log p$ are $\Q$-independent, so there is a $\Q$-linear
$\mathfrak s:\R\to\mm_\R$ with $\mathfrak s(\log p)=\omega^{-\omega^2/p}$ for every
prime $p$. Let $\chi(P)=\cis(\mathfrak s(c_\omega(P)))$. Then
$E_\chi(-s\log p)=\exp_{\No}(-\omega\log p)\,\cis(-\omega^{-\omega^2/p})$, and
$\exp_{\No}(-\omega\log p)=\omega^{-\omega\log p}$ is a monomial. So the $p$th
member contains the exponent $\omega\log p+\omega^2/p$ (in the valuation
convention) with coefficient $-i$, and no other part of that member cancels
it. For $p<p'$,
$(\omega\log p+\omega^2/p)-(\omega\log p'+\omega^2/p')=\omega^2\frac{p'-p}{pp'}-\omega\log\frac{p'}p>0$.
So the union of supports contains an infinite strictly decreasing sequence
and is not well ordered.
\end{proof}
This makes precise 03's caution: ``For general admissible phases, the
termwise reduction used in its proof need not hold.'' No source asserts anything on $\Wing$
without naming a phase.

\begin{corollary}[Wing fibres \merge; from \src{01, 03}]\label{gz:cor:wingfibre}
For every nonzero $u\in\mm$,
$\{s\in\Wing:\zeta^\Wing_{\Exp}(s)=1+u\}=\{s_k(u)+iY:k\in\Z,\ Y\in\Pi\}$.
\end{corollary}
\begin{proof}
$\zeta^\Wing_{\Exp}(x+iy)=1+u$ if and only if $x+i\fin y$ lies in the $\Tp$
fibre, $\{s_k(u)\}$ by \cref{gz:thm:fibre,gz:cor:complexfibre}. Since
$\Im s_k(u)$ is finite, $x+iy=s_k(u)+i\pp(y)$, and conversely
$\fin(\Im s_k+Y)=\Im s_k$.
\end{proof}
None of this extends the completed Gamma--zeta package to $\Vs$, and there is
no value at $\tfrac12+i\omega$ \src{02, 03}. A formal Stirling series can be
summable on large complex sectors while its branch, Stokes data and
compatibility across sectors remain unspecified \src{03}.

\part{Global objects on the horizontal tube}\label{gz:part:tube}

This part is source 03's construction, stated for an arbitrary phase
$\Psi_\chi$ and renamed as in \cref{gz:tab:rename}. The domain is the
horizontal tube $\Th=\Tm\sqcup\OO\sqcup\Tp=\{\Im s\in\OO_\R\}$. It contains
every ordinary complex number and points such as $\omega+i$ and
$-e^\omega+3i$, but not $\tfrac12+i\omega$. A function is \emph{locally
Hahn-holomorphic} if near each point it is given by a specified strongly
evaluable Taylor germ, and \emph{locally Hahn-meromorphic} if it is locally
a quotient of such germs. A pole is a Laurent germ, not a field value. Every
point has a fine neighbourhood inside its own region, so the definitions
below give local analytic data without asserting a connected continuation
between $\Tm$, $\OO$ and $\Tp$. On $\Tm$, reflection and the functional
equation are \emph{definitions}. The content is that these definitions are
consistent with the independently constructed pieces, have the stated local
structure, and obey exact divisor and cancellation theorems \src{03}.

\section{Reflection forces infinite Gamma poles}\label{gz:sec:forced}
\begin{theorem}[Forced poles \src{03, 06}]\label{gz:thm:forced}
Let $\Psi_\chi$ be any phase.
\begin{enumerate}[label=(\alph*)]
\item \src{06} Let $G$ be a function with values in $\SC$ on a set of real
 surreals, satisfying $G(x)G(1-x)\Sin_\chi(\pi x)=\pi$ whenever $G(x)$ and
 $G(1-x)$ are both defined. Then for every $x\in\cI(\Psi_\chi)$, $G$ is
 undefined at $x$ or at $1-x$. For $\Exp$, this holds for every omnific
 integer $x$. In particular no $G$ defined at every real surreal except
 possibly the \emph{ordinary} nonpositive integers satisfies the relation.
\item \src{03} If $g$ is locally Hahn-holomorphic on $\Th$, nonzero at every
 positive infinite real point, and satisfies $g(z)g(1-z)=\Sin_\chi(\pi z)/\pi$,
 then $g(x)=0$ at every $x\in\cI(\Psi_\chi)$ with $x<-\R$. So a zero-free
 meromorphic Gamma $1/g$ with reflection has poles there.
\end{enumerate}
\end{theorem}
\begin{proof}
By \cref{gz:thm:periods}, $\Sin_\chi(\pi x)=0$ for every $x\in\cI(\Psi_\chi)$,
so the relation cannot hold at $x$ with both values defined. $\cI$ has
positive infinite elements $x$, and then $x$ and $1-x$ are both infinite.
For (b), $1-x>\R$, so $g(1-x)\ne0$ and $g(x)=0$.
\end{proof}
Both parts are one-line consequences of \rl{trigonometry:thm:infiniteperiods}
and \rl{trigonometry:cor:globalzeros}, so forced poles are immediate once
infinite periods exist. 06 presents the existence of infinite periods as
known, citing the trigonometry report but not this label, and only its
application to Gamma as new (error 41). The pole set of \cref{gz:thm:gamma-tube} satisfies (a): for
$x\in\cI$ with $x>\R$, the element $1-x\in\cI$ is a pole below $-\R$; for
finite $x\in\cI=\Z$, either $x\le0$ or $1-x\le0$. \emph{Reconciliation:} 03
builds Gamma with reflection on $\Th$ while 06 proves that reflection with
only ordinary poles is impossible. These are the same fact. 06's theorem
concerns arbitrary functions and is part (a); 03's is part (b), the form for
the local germ class, which also says at which of $x$, $1-x$ the pole lies.
The conclusion is not that global Gamma
functions are impossible. A global reflection formula must allow extra
infinite poles, restrict its domain, or change another part of the
prescription \src{06}. The ordinary nonpositive integers are the complete
pole set only for the finite lift. Transported models choose the first
alternative, with poles at the internal integers
(\cref{gz:sec:transported-divisors}). The obstruction uses no contour
integral, no Stirling coefficient and no hypothesis on primes \src{03}.

\section{Gamma on the horizontal tube}\label{gz:sec:gammatube}
Define the reciprocal
\[
g_\chi(z)=\begin{cases}(1/\Gamma)^\#(z),&z\in\OO,\\ E_\chi(-\LSt(z)),&z\in\Tp,\\ \dfrac{\Sin_\chi(\pi z)}\pi\,\Gamma_\chi(1-z),&z\in\Tm,\end{cases}
\]
where $\Gamma_\chi(1-z)=E_\chi(\LSt(1-z))$ on $\Tp$. Put $\Gamma_\chi=1/g_\chi$,
read meromorphically at zeros of $g_\chi$. On $\Tp\subset\Sinf$ this is
the $\Gamma_\chi$ of \cref{gz:prop:gammachi}, so $\Gamma_\chi$ is one function
on $\Sinf\cup\Th$.

\begin{theorem}[Gamma on $\Th$ \src{03}]\label{gz:thm:gamma-tube}
$g_\chi$ is locally Hahn-holomorphic on $\Th$. $\Gamma_\chi$ has no zeros,
and its poles are simple and form exactly
\[
\Z_{\le0}\cup\bigl(\cI(\Psi_\chi)\cap\{x<-\R\}\bigr).
\]
As identities of meromorphic germs on $\Th$,
\[
\Gamma_\chi(z+1)=z\Gamma_\chi(z),\qquad g_\chi(z)g_\chi(1-z)=\frac{\Sin_\chi(\pi z)}\pi,\qquad
\Gamma_\chi(z)\Gamma_\chi(1-z)=\frac\pi{\Sin_\chi(\pi z)},
\]
and at a real pole $x$, $\Res_{z=x}\Gamma_\chi(z)=\Cos_\chi(\pi x)/\Gamma_\chi(1-x)$.
For $\Exp$ the pole set is exactly $\Oz_{\le0}$.
\end{theorem}
\begin{proof}
The first two cases are holomorphic by \cref{gz:thm:divisor} and
\cref{gz:prop:gammachi}. On $\Tm$ both factors are holomorphic and the second
is nonzero. So the zeros of $g_\chi$ on $\Tm$ are those of $\Sin_\chi(\pi z)$,
which are real, simple and equal to $\cI\cap\{x<-\R\}$ by
\cref{gz:thm:periods}. The finite zeros are the ordinary nonpositive
integers, and there are none on $\Tp$. For left-infinite $z$, the recurrence
on $\Tp$ gives $\Gamma_\chi(1-z)=(-z)\Gamma_\chi(-z)$, and $\Sin_\chi$ changes
sign under $z\mapsto z+1$ (since $\Psi_\chi(\pi)=-1$), so $g_\chi(z)=zg_\chi(z+1)$.
The other regions use the finite germ identity and \cref{gz:thm:stirling}.
Reflection on $\Tm$ is the definition; on $\Tp$ apply it to $1-z$ and use
$\Sin_\chi(\pi(1-z))=\Sin_\chi(\pi z)$; on $\OO$ it is the lifted classical
identity. At a pole $x$ write $z=x+h$: then
$g_\chi(x+h)=h\Cos_\chi(\pi x)\Gamma_\chi(1-x)+\bigO(h^2)$, and
$\Cos_\chi(\pi x)=\pm1$, which gives the residue. It reduces to $(-1)^n/n!$ at
$x=-n$. Finally $\cI(\Phi)=\Oz$ and $\Oz\cap\OO_\R=\Z$.
\end{proof}

\begin{theorem}[Gauss on $\Th$ \src{03}]\label{gz:thm:gauss-tube}
For every ordinary $m\ge1$ and $z\in\Th$,
$\prod_{r<m}\Gamma_\chi(z+r/m)=(2\pi)^{(m-1)/2}m^{1/2-mz}\Gamma_\chi(mz)$
meromorphically, with $m^{1/2-mz}=E_\chi((\tfrac12-mz)\log m)$.
\end{theorem}
\begin{proof}
The finite case is classical and the case $\Tp$ is \cref{gz:prop:gammachi}.
For $z\in\Tm$ use $\prod_{r<m}\Sin_\chi(\pi(z+r/m))=2^{1-m}\Sin_\chi(\pi mz)$.
To prove it, replace each sine by exponentials and factor $U^m-1$ over the
ordinary $m$th roots of unity; only the group law and finite phases enter.
The complementary arguments $1-z-r/m$ permute the $w+r/m$ with
$w=1/m-z\in\Tp$, so their Gamma product is
$(2\pi)^{(m-1)/2}m^{mz-1/2}\Gamma_\chi(1-mz)$. Multiply the reciprocal
definitions. The computation is an identity of holomorphic reciprocal germs,
so no division by a vanishing sine occurs.
\end{proof}

\begin{proposition}[Uniqueness for fixed data \src{03}]\label{gz:prop:tube-unique}
Once the phase, the finite lift and the exact Stirling branch on $\Tp$ are
fixed, there is exactly one zero-free locally meromorphic Gamma on $\Th$
satisfying reflection.
\end{proposition}
\begin{proof}
Reflection determines the reciprocal, with its germs at the zeros, on $\Tm$
from the already specified values on $\Tp$. Existence is
\cref{gz:thm:gamma-tube}.
\end{proof}
This is uniqueness for fixed data, and it is immediate: reflection is the
definition of $g_\chi$ on $\Tm$. It is not a surreal Bohr--Mollerup
theorem, and the gauges of \cref{gz:sec:nonunique} show that the positive
branch is itself a selection.

\begin{example}[The canonical case \src{03}]\label{gz:ex:canonical-gamma}
$\Res_{z=-\omega}\Gamma_{\Exp}(z)=1/\GSt(1+\omega)$, because
$\Cos(\pi\omega)=\cos(\fin\pi\omega)=1$. By reflection (a direct
consequence of \cref{gz:thm:gamma-tube}, not printed in 03),
$\Gamma_{\Exp}(-\omega-\tfrac12)=-\pi/\GSt(\omega+\tfrac32)$, because
$\Sin(-\pi\omega-\tfrac\pi2)=\sin(-\tfrac\pi2)=-1$. On real $x<-\R$,
$\Gamma_{\Exp}(x)=\pi/(\sin(\pi\fin x)\GSt(1-x))$ away from $\Oz$.
\end{example}

\section{Zeta and \texorpdfstring{$\xi$}{xi} on the horizontal tube}\label{gz:sec:zetatube}
With the same phase everywhere, put
\[
\FE_\chi(s)=2^s\pi^{s-1}\Sin_\chi(\pi s/2)\,\Gamma_\chi(1-s),
\]
where $2^s=E_\chi(s\log2)$ and $\pi^{s-1}=E_\chi((s-1)\log\pi)$. On $\Th$ these
powers are phase-free, because $\Im s$ is finite. (03 writes $\chi_\phi$ for
$\FE_\chi$; the letter $\chi$ is the character here.)

\begin{lemma}[The reflection factor \src{03}]\label{gz:lem:FEfactor}
On $\Th$, as meromorphic germ identities, $\FE_\chi(s)\FE_\chi(1-s)=1$ and
$\FE_\chi(s)=\pi^{s-1/2}\Gamma_\chi(\tfrac{1-s}2)/\Gamma_\chi(\tfrac s2)$.
\end{lemma}
\begin{proof}
In the product, the powers give $2/\pi$, and
$\Sin_\chi(\pi s/2)\Sin_\chi(\pi(1-s)/2)=\tfrac12\Sin_\chi(\pi s)$. Reflection
cancels the rest, in the field of germs. For the second identity use
duplication, $\Gamma_\chi(1-s)=\frac{2^{-s}}{\sqrt\pi}\Gamma_\chi(\tfrac{1-s}2)\Gamma_\chi(1-\tfrac s2)$
(\cref{gz:thm:gauss-tube} with $m=2$), and reflection at $s/2$.
\end{proof}

Define
\[
\zeta_\chi(s)=\begin{cases}\zp(s),&s\in\OO,\\ \zD(s),&s\in\Tp,\\ \FE_\chi(s)\,\zD(1-s),&s\in\Tm.\end{cases}
\]
This is not circular: for left-infinite $s$, the Gamma factor lies on the
independent positive branch.

\begin{theorem}[Zeta on $\Th$ \src{03}]\label{gz:thm:zeta-tube}
$\zeta_\chi$ is locally Hahn-meromorphic on $\Th$, with a single pole at $1$,
of residue $1$. It satisfies $\zeta_\chi(s)=\FE_\chi(s)\zeta_\chi(1-s)$
throughout $\Th$ and commutes with conjugation. Its zeros are the ordinary
zeros of $\zeta$ with their multiplicities, together with the simple real
zeros $\{x<-\R:x\in2\cI(\Psi_\chi)\}$. At such a zero,
\[
\zeta_\chi'(x)=2^{x-1}\pi^x\Cos_\chi(\pi x/2)\,\Gamma_\chi(1-x)\,\zD(1-x)\ne0 .
\]
\end{theorem}
\begin{proof}
On $\OO$ the claims are the germ claims, with no displaced zeros by
\cref{gz:thm:divisor}. On $\Tp$, $\zD$ is holomorphic and zero-free
(\cref{gz:thm:euler}). On $\Tm$ every factor except $\Sin_\chi(\pi s/2)$ is
holomorphic and nonzero, so the zeros are those of $\Sin_\chi(\pi s/2)$,
real and simple by \cref{gz:thm:periods}. Differentiating the product at
such a zero leaves only the sine derivative. The functional equation is the
definition on $\Tm$. On $\Tp$ it follows from the definition at $1-s$ and
\cref{gz:lem:FEfactor}, and on $\OO$ it is the lifted classical equation.
Conjugation is preserved by each ingredient, including the strong Dirichlet
sum with real coefficients.
\end{proof}

\begin{corollary}[Canonical trivial zeros \src{03}]\label{gz:cor:canonical-zeros}
For $\Exp$ the zeros of $\zeta_{\Exp}$ on $\Th\cap\R$ below $-\R$ are exactly
$2\Oz_{<0}$ restricted to $x<-\R$. Together with $-2,-4,\dots$ they are
$2\Oz_{<0}$. For example $\zeta_{\Exp}(-\omega)=0$, since $-\omega/2\in\Oz$.
But $-\omega-1\notin2\Oz$, since $(-\omega-1)/2$ has finite part
$-\tfrac12$. Moreover $\zeta_{\Exp}'(-2\omega)=2^{-2\omega-1}\pi^{-2\omega}\GSt(1+2\omega)\zD(1+2\omega)$.
\end{corollary}
``Even'' means membership in the additive group $2\Oz$. No recurrence is
iterated a surreal number of times, and no product with an omnific number of
factors occurs \src{03}.

\begin{theorem}[Completed-divisor rigidity \src{03}]\label{gz:thm:xi-tube}
$\xi_\chi(s)=\tfrac12s(s-1)\pi^{-s/2}\Gamma_\chi(s/2)\zeta_\chi(s)$ extends to a
locally Hahn-holomorphic function on $\Th$ with $\xi_\chi(s)=\xi_\chi(1-s)$.
It has no zeros on $\Tp\cup\Tm$ and equals $\xip$ on $\OO$. So its zero
divisor is exactly the classical nontrivial divisor, independent of $\chi$,
with no zeros at nonstandard finite points and none at infinite real part.
\end{theorem}
\begin{proof}
With \cref{gz:lem:FEfactor},
$\pi^{-s/2}\Gamma_\chi(s/2)\FE_\chi(s)=\pi^{-(1-s)/2}\Gamma_\chi(\tfrac{1-s}2)$,
and $s(s-1)$ is symmetric, which proves the functional equation for germs.
On $\Tp$ every factor is holomorphic and nonzero. On $\Tm$ the functional
equation expresses $\xi_\chi$ through the nonzero function on $\Tp$, so every
apparent singularity is removable with a nonzero value. On $\OO$ the
expression is the classical identity for the entire $\xi$, so it equals
$\xip$, whose zeros are given by \cref{gz:thm:divisor}.
\end{proof}

\begin{corollary}[Horizontal RH \src{03}]\label{gz:cor:horizontalRH}
For every phase, ``$\xi_\chi(s)=0$ with $s\in\Th$ implies $\Re s=\tfrac12$''
is equivalent to classical RH.
\end{corollary}
The two statements quantify over the same zeros. The content is the
cancellation theorem, and nothing is said at $\tfrac12+iT$ with infinite
$T$, which lies outside $\Th$ \src{03}. At a nontrivial zero $\rho$ of order
$m$, $\xi_\chi(\rho+\eps)=\frac{\xi^{(m)}(\rho)}{m!}\eps^m(1+\eta)\ne0$,
unconditionally and for every $m$. Simplicity is not smuggled in.

\begin{proposition}[Explicit cancellation \src{03}]\label{gz:prop:cancel}
Let $A>\R$ with $A\in\cI(\Psi_\chi)$. For $h$ small enough for the local
germs (possibly much smaller than $1/\log A$),
\begin{gather*}
\Gamma_\chi(-A+h/2)=\frac{2\Cos_\chi(\pi A)}{\Gamma_\chi(1+A)}h^{-1}+\bigO(1),\\
\zeta_\chi(-2A+h)=2^{-2A-1}\pi^{-2A}\Cos_\chi(\pi A)\Gamma_\chi(1+2A)\zD(1+2A)\,h+\bigO(h^2),
\end{gather*}
and, since $\Cos_\chi(\pi A)^2=1$,
$\xi_\chi(-2A)=A(2A+1)2^{-2A}\pi^{-A}\frac{\Gamma_\chi(1+2A)}{\Gamma_\chi(1+A)}\zD(1+2A)=\xi_\chi(1+2A)$.
The last equality follows from duplication,
$2^{-2A}\Gamma_\chi(1+2A)/\Gamma_\chi(1+A)=\Gamma_\chi(A+\tfrac12)/\sqrt\pi$.
\end{proposition}
The computation keeps the local multiplicities and gives the removable
value; no symbol $0\cdot\infty$ is used. At $A=\omega$ with $\Exp$, the
residue $1/\GSt(1+\omega)$ and the zero derivative at $-2\omega$ live on very
different scales, but their product is fixed by duplication.

\subsection{Phase changes move the infinite divisors}
\begin{proposition}[A movable infinite zero \src{01, 03}]\label{gz:prop:moving}
Fix a positive infinite monomial $M$ and consider the phases $\chi_{\lambda,M}$
of \cref{gz:prop:families}(b). They all give the same finite Gamma and zeta,
the same Gamma on the positive real axis, and the same zeta on $\Tp$. But the
left-infinite real zeros of $\zeta_{\chi_{\lambda,M}}$ in $-M+\OO_\R$ are
exactly $-M+\lambda+2k$ ($k\in\Z$), so $\zeta_{\chi_{\lambda,M}}(-M)=0$ if and
only if $\lambda\in2\Z$. The completed $\xi$ divisors are identical.

For $M=\omega$ this is 01's one-sided prescription:
$\zeta_{\chi_\lambda}$ on $\Tm$ has the zero set
$\{\sigma<-\R:P_\lambda(\sigma)\in2\Z\}$, all real and simple, and
\[
\zeta_{\chi_\lambda}(-\omega)=-A_\omega\sin(\pi\lambda/2),\qquad
A_\omega=2^{-\omega}\pi^{-\omega-1}\GSt(1+\omega)\zD(1+\omega)>0 .
\]
In particular $\zeta_{\Exp}(-\omega)=0$ with $\zeta_{\Exp}'(-\omega)=\tfrac\pi2A_\omega$,
while $\zeta_{\chi_1}(-\omega)=-A_\omega\ne0$.
\end{proposition}
\begin{proof}
$-M+r\in2\cI(\Psi_{\chi_{\lambda,M}})$ if and only if $r-\lambda\in2\Z$, by
\cref{gz:prop:families}(b); apply \cref{gz:thm:zeta-tube}, and
\cref{gz:thm:xi-tube} for $\xi$. For $M=\omega$, 01's
$Z_{-,\lambda}(s)=2^s\pi^{s-1}\sin_\lambda(\pi s/2)\Gamma_{\mathrm{St},\lambda}(1-s)Z_\infty(1-s)$
is literally $\FE_{\chi_\lambda}(s)\zD(1-s)$. At $-\omega$,
$\Sin_{\chi_\lambda}(-\pi\omega/2)=\sin(\lambda c_\omega(-\pi\omega/2))=-\sin(\pi\lambda/2)$.
The derivative is \cref{gz:thm:zeta-tube} with $\Cos(-\pi\omega/2)=1$.
\end{proof}

\begin{remark}[01's remark after its left-zero theorem, corrected]\label{gz:rem:stale01}
01 says its result ``does not say that a canonical surreal zeta must have a
zero at $-\omega$, or that it must not'', and that a simultaneous functional
equation on a larger Gamma domain ``would require an additional
construction and proof''. That construction is 03's
(\cref{gz:thm:gamma-tube,gz:thm:zeta-tube}), and it contains 01's
prescription. With Ehrlich and Kaplan's canonical $\Exp$ ($\lambda=0$) the
zero at $-\omega$ is present. No reflection or functional-equation
requirement forces it, since every $\lambda$ gives an FE-compatible package
on $\Th$ (error 2). The canonical initial-integer-part convention is a
meaningful way to choose, and it is not a consequence of reflection
\src{03}. 01's non-claim, that its one-sided prescription is not a proved
global functional equation, stays true of that prescription taken alone.
01's third research direction, a coherent continuation deciding which left
zeros survive, is thereby answered on $\Th$ relative to a named phase:
every phase gives such a package, with left zeros
$2\cI(\Psi_\chi)\cap\{x<-\R\}$. At infinite imaginary height it remains
open (\cref{gz:q:nonresonant}).
\end{remark}
Agreement on all ordinary arguments, on the positive-real Gamma and on the
right-infinite zeta therefore does not determine the negative-infinite
divisor. Every phase has infinite poles \src{03}.

\part{Infinite imaginary parts: obstructions and underdetermination}\label{gz:part:infinite}

The genuinely new points for RH are $\sigma+iT$ with $0<\sigma<1$ and $T$
infinite, the strip $\Vs$. None of the constructions of Parts II--IV assigns
zeta values there. The finite lift covers $\OO$. The Dirichlet algebra and
the wing have positive infinite real part. The tube has finite imaginary
part. This part shows what goes wrong in the obvious attempts \src{02, 03,
04, 05, 06}.

\section{Reflection versus the Stirling modulus on the vertical strip}\label{gz:sec:strip}

\begin{theorem}[The reflection defect \src{05} for $\sigma=\tfrac12$; all $\sigma$ \merge]\label{gz:thm:strip}
Let $T>\R$ and let $\sigma$ be finite with $0<\sigma<1$, and put
$z=\sigma+iT$. Then $z,1-z\in\Sinf$, and:
\begin{enumerate}[label=(\alph*)]
\item $\LSt(z)+\LSt(1-z)=\log2\pi-\pi T+i\pi(\sigma-\tfrac12)$ exactly: every
 inverse-power term cancels. In particular $2\Re\LSt(\tfrac12+iT)=\log2\pi-\pi T$.
\item For every homomorphism $E$ extending $\Efp$, and so for every phase,
 $\Sin_E(\pi z)=(E(i\pi z)-E(-i\pi z))/(2i)$ is phase-free, and
 $|\pi/\Sin_E(\pi z)|=2\pi e^{-\pi T}/|1-e^{2\pi iz}|$ with
 $e^{2\pi iz}=e^{-2\pi T}\cis(2\pi\sigma)$.
\item Let $G$ be defined at $z$ and $1-z$ with Stirling moduli,
 $|G(w)|=\exp_{\No}\Re\LSt(w)$ for $w=z,1-z$. Then
 $G(z)G(1-z)=\pi/\Sin_E(\pi z)$ fails unless $|1-e^{2\pi iz}|=1$. For each $T$
 this happens for exactly two abscissae, $\sigma_-(T)=\arccos(e^{-2\pi T}/2)/2\pi$
 (the arccosine taken as the unique finite angle in $[0,\pi]$ with the given
 cosine, which exists by \cref{gz:prop:polar})
 and $\sigma_+(T)=1-\sigma_-(T)$, which are infinitesimally close to
 $\tfrac14$ and $\tfrac34$.
\item Any $G$ satisfying reflection at $z$ has the total log-modulus defect
 \[
 \log|G(z)|+\log|G(1-z)|-\Re\bigl(\LSt(z)+\LSt(1-z)\bigr)=-\tfrac12\log\bigl(1-2e^{-2\pi T}\cos2\pi\sigma+e^{-4\pi T}\bigr),
 \]
 which is smaller than every $T^{-N}$ and nonzero unless $\sigma=\sigma_\pm(T)$.
 At $\sigma=\tfrac12$ with $G(\bar z)=\overline{G(z)}$ \src{05},
 $\log|G(\tfrac12+iT)|=\Re\LSt(\tfrac12+iT)-\tfrac12\log(1+e^{-2\pi T})$,
 and the correction is $-\tfrac12e^{-2\pi T}+\tfrac14e^{-4\pi T}-\tfrac16e^{-6\pi T}+\cdots$.
\item \merge{} For every additive-to-multiplicative $E:\SC\to\SC^\times$
 extending $\Efp$ (no modulus law is assumed), and so for every phase and
 for 06's $\Gamma_{E,\rm St}=E\circ\LSt$,
 \[
 \Gamma_{E,\rm St}(z)\,\Gamma_{E,\rm St}(1-z)=\bigl(1-e^{2\pi iz}\bigr)\,\frac\pi{\Sin_E(\pi z)} .
 \]
 Since $e^{2\pi iz}\ne0$, $\Gamma_{E,\rm St}$ violates reflection at $z$ for
 every finite $\sigma\in(0,1)$, including $\sigma_\pm(T)$: there the factor
 $1-e^{2\pi iz}$ has modulus $1$ but argument $\pm e^{-2\pi T}(1+\eps)$
 at $\sigma_\pm(T)$, with $\eps\in\mm$. The same holds at $\sigma-iT$, so $\Gamma_{E,\rm St}$
 violates reflection at every point of $\Vs$.
\end{enumerate}
\end{theorem}
\begin{proof}
(a) Apply \cref{gz:thm:translation} twice. First take $Z=iT$ with
$\Log Z=\log T+i\pi/2$ and $a=\sigma$. The principal $\Log(\sigma+iT)$ equals
$\log T+i\pi/2+\Lgm(1+\sigma/(iT))$: both are logarithms of $\sigma+iT$, and
both imaginary parts lie in $(-\pi/2,\pi/2)$, since
$\Im\Lgm(1-i\sigma/T)=-\sigma/T+\cdots<0$. Then take $Z'=-iT$ with
$\Log Z'=\log T-i\pi/2$ and $a'=1-\sigma$. The sum of the two elementary
parts is
$(iT+\sigma-\tfrac12)(\log T+\tfrac{i\pi}2)-iT+(-iT+\tfrac12-\sigma)(\log T-\tfrac{i\pi}2)+iT+\log2\pi=\log2\pi-\pi T+i\pi(\sigma-\tfrac12)$.
The $n$th inverse-power terms are
$\frac{(-1)^{n+1}}{n(n+1)}\bigl(B_{n+1}(\sigma)(iT)^{-n}+B_{n+1}(1-\sigma)(-iT)^{-n}\bigr)=0$,
because $B_{n+1}(1-\sigma)=(-1)^{n+1}B_{n+1}(\sigma)$ and
$(-iT)^{-n}=(-1)^n(iT)^{-n}$. (b) $i\pi z=-\pi T+i\pi\sigma$ has finite
imaginary part, so $E(\pm i\pi z)=e^{\mp\pi T}\cis(\pm\pi\sigma)$ and
$\Sin_E(\pi z)=-\frac{e^{\pi T}\cis(-\pi\sigma)}{2i}(1-e^{2\pi iz})$. (c) By
(a) the Stirling moduli give $|G(z)G(1-z)|=2\pi e^{-\pi T}$. With
$q=e^{-2\pi T}$, $|1-q\cis2\pi\sigma|^2=1-2q\cos2\pi\sigma+q^2$ equals $1$
exactly when $\cos2\pi\sigma=q/2$. For finite $\sigma\in(0,1)$ this has
exactly the two solutions stated. (d) Take logarithms of the moduli in (b).
At $\sigma=\tfrac12$, $1-z=\bar z$ and $|1-e^{2\pi iz}|=1+e^{-2\pi T}$.
Conjugation splits the defect equally. (e) By (a),
$w_0=\LSt(z)+\LSt(1-z)=\log2\pi-\pi T+i\pi(\sigma-\tfrac12)$ has finite
imaginary part, so $E(w_0)=\Efp(w_0)=-2\pi i\,e^{-\pi T}\cis(\pi\sigma)$.
Since $E$ is a homomorphism,
$\Gamma_{E,\rm St}(z)\Gamma_{E,\rm St}(1-z)=E(\LSt(z))E(\LSt(1-z))=E(w_0)$.
By (b), $\pi/\Sin_E(\pi z)=-2\pi i\,e^{-\pi T}\cis(\pi\sigma)/(1-e^{2\pi iz})$,
and comparing the two gives the identity. Here $e^{2\pi iz}=\Efp(2\pi iz)$
has modulus $e^{-2\pi T}\ne0$, so $1-e^{2\pi iz}\ne1$. With $q=e^{-2\pi T}$
and $\cos2\pi\sigma_\pm=q/2$, $1-e^{2\pi iz}=1-\tfrac{q^2}2-iq\sin2\pi\sigma_\pm$,
where $\sin2\pi\sigma_\pm=\mp(1-q^2/4)^{1/2}$; this has modulus $1$ and
argument $\pm q(1+\eps)$. Finally, for $z'=\sigma-iT$ put $w=1-z'=(1-\sigma)+iT$.
Then $\Gamma_{E,\rm St}(z')\Gamma_{E,\rm St}(1-z')=\Gamma_{E,\rm St}(w)\Gamma_{E,\rm St}(1-w)$,
and $\Sin_E(\pi z')=\Sin_E(\pi-\pi w)=\Sin_E(\pi w)$ because $E(\pm i\pi)=-1$,
so the identity at $w$ gives the failure at $z'$.
\end{proof}
05 proves the case $\sigma=\tfrac12$ (its Theorem 7.1 and Corollary 7.2),
including $\sin(\pi(\tfrac12+iT))=\cosh\pi T$. It uses only a finite real
phase and the real exponential. Classically the correction is encoded in
$|\Gamma(\tfrac12+it)|^2=\pi/\cosh\pi t$. The extension to every finite
$\sigma\in(0,1)$, and the two exceptional abscissae where the modulus test
is silent, were established for this merge. The reconciliation plan had
claimed the obstruction for every $\sigma$. The exceptions $\sigma_\pm(T)$
were found while the proof was rechecked. At them the classical identity
$\Re(\log\Gamma(z)+\log\Gamma(1-z))=\log2\pi-\pi T$ holds as well (checked
numerically at ordinary $T$, \cref{gz:app:checks}), so the modulus test alone
claims nothing there for a general $G$ with Stirling moduli; for
$\Gamma_{E,\rm St}$ itself, (e) still gives failure. Clause (e) was added
after a post-merge audit (\cref{gz:app:audit}). The symbolic check of (a)
through $T^{-14}$ and the numerical check of (d) are in the reconciliation
suite; the check of (e) is described in \cref{gz:app:checks}.

\emph{Reconciliation.} $\tfrac12\pm iT$ lies in $\Sinf$ but not in $\Th$, so
the theorem does not touch 03's reflection on the tube
(\cref{gz:thm:gamma-tube}). It does apply to 06's $\Gamma_{E,\rm St}=E\circ\LSt$
on $\Sinf$ for every $E$ extending $\Efp$ (by (e) no modulus law is
needed), in particular to $\Gamma_{\Exp}$ and every $\Gamma_\chi$. These
satisfy the recurrence on $\Sinf$ (\cref{gz:prop:gammachi}) but violate
reflection at every point of $\Vs$, by (e) (the remark missing from 06,
error 48). The exceptional abscissae concern only an arbitrary $G$ with
Stirling moduli. A conjugation-
and reflection-compatible Gamma on the line $\tfrac12+i\R_{>\R}$ must add
the flat correction, and even then it has fixed only a modulus: no phase and
no global Gamma function \src{05}. The theorem is not an objection to
asymptotic Stirling expansions, to positive-real surreal Gamma, or to
sectorial resummation. It is an exact incompatibility test for a proposed
surreal prescription \src{05}. It matters for zeta directly, since
$\Gamma(s/2)$ in $\xi$ lies in exactly this regime when $s\in\Vs$. A global
$\xi$ cannot use bare Stirling values, keep every classical Gamma identity
and expect compatibility. Phase data and beyond-all-orders corrections must
both be accounted for \src{05}.

\section{Resonant galaxies}\label{gz:sec:galaxies}

\subsection{A phase does not give strong summability}
At $s=\sigma+iT$ with ordinary $\sigma$ and infinite $T$, and for any unit
phase, the Dirichlet terms $n^{-\sigma}\Psi_\chi(-T\log n)$ are finite units
with nonzero standard parts. The family is not strongly summable, even for
$\sigma>1$ \src{04}. This obstruction is separate from the phase question:
prescribing every phase does not turn the Dirichlet series at infinite
height into a Hahn sum. Similarly, the Riemann--Siegel cutoff
$\sqrt{T/2\pi}$ at infinite $T$ does not give an ordinary finite sum. In a
bare surreal field, ``$n\le\sqrt T$'' means all ordinary $n$, and in an
elementary model it means a hyperfinite internal sum, which is additional
structure \src{04}. With $\Exp$, $n^{-iT}=\cis(-\fin(T)\log n)$, so the terms
at infinite height reduce to terms at finite height. They are ordinarily
summable in $\C$ but not strongly summable.

\subsection{Coherent continuation on a resonant galaxy}
Let $T>0$ be a resonance of $\Psi_\chi$ (\cref{gz:def:resonance}). Then $-T$ is
one too. On the galaxy $\pm iT+\OO$ write $s=\pm iT+w$ with $w$ finite. For
every ordinary $n$,
\begin{equation}\label{gz:eq:resonantpowers}
E_\chi(-s\log n)=\Efp(-w\log n)\,\Psi_\chi(\mp T\log n)=n^{-w}.
\end{equation}
So coefficientwise ordinary Dirichlet summation on $\Re\st w>1$ reproduces
$\zp(w)$, not a new high-frequency function.

\begin{definition}[Coherent meromorphic continuation \src{06}]\label{gz:def:coherent}
A \emph{coherent meromorphic continuation in the galaxy coordinate} is a
section with one set-sized well-ordered Hahn support and ordinary
meromorphic coefficient functions on the connected ordinary $w$-plane,
evaluated at finite $w$ by Taylor--Laurent lifting. It is stronger than
fine analyticity at each point separately.
\end{definition}
06 calls this ``the common-domain notion employed in the repository''. It
\emph{extends} the analysis report's coherent sections (\rl{c:p4:def-HU}),
which have holomorphic coefficients on a connected open set, to meromorphic
coefficients on the whole plane. The identity theorem still applies off the
discrete pole sets (error 45).

\begin{theorem}[Resonant-galaxy obstruction \src{06}; general form written out \merge]\label{gz:thm:resonance}
Let $T>0$ be a resonance of $\Psi_\chi$. Suppose $Z$ uses $E_\chi$ for its
Dirichlet powers, agrees with coefficientwise ordinary Dirichlet summation
on $\Re\st w>1$ in both galaxies $\pm iT+\OO$, and has coherent meromorphic
continuations on their ordinary coordinate planes. Then
$Z(iT+w)=\zp(w)$ and $Z(-iT+w)=\zp(w)$ as meromorphic functions of finite $w$.
In particular:
\begin{enumerate}[label=(\roman*)]
\item $Z$ has poles at $1\pm iT$, so it cannot have the pole at $1$ alone;
\item $Z(3+iT)=\zeta(3)\ne0$, whereas $Z(-2-iT)=0$;
\item the pointwise functional equation
 $Z(s)=2^s\pi^{s-1}\Sin_\chi(\pi s/2)G(1-s)Z(1-s)$, with $E_\chi$ powers,
 fails at $s=3+iT$ for every Gamma factor $G$ that is field-valued at
 $-2-iT$.
\end{enumerate}
\end{theorem}
\begin{proof}
By \eqref{gz:eq:resonantpowers} and the specified summation,
$Z(iT+w)=\zp(w)$ on $\Re\st w>1$. Subtract the canonical section of $\zeta$
from the coherent continuation. On that half-plane every Hahn coefficient of
the difference vanishes at every ordinary $w$, so by the ordinary identity
theorem each coefficient vanishes on the connected plane. The same holds on
$-iT+\OO$. (i) and (ii) are ordinary zeta facts. For (iii), at $s=3+iT$ the
right side is a product of field elements whose last factor is
$Z(-2-iT)=\zeta(-2)=0$, while the left side is $\zeta(3)$.
\end{proof}
The Gamma qualification in (iii) is not vacuous. The shifted Stirling value
\[
G(-2-iT)=\frac{\Gamma_\chi(1-iT)}{(-2-iT)(-1-iT)(-iT)}
\]
is field-valued and nonzero, since $1-iT\in\Sinf$ \src{06}. Introducing an
actual Gamma pole at $-2-iT$ would change the pole prescription and require
a meromorphic cancellation, not a pointwise product of field values.

For $\Exp$ every nonzero $T\in\Pi$ is a resonance; this is the case 06
states as its theorem. 06's Remark 7.4 already notes that every phase with
a common resonance ``encounters the same argument''; the proof above writes
it out. By \cref{gz:prop:resonance}(a), every phase with ordinary-circle
character values has resonances of arbitrarily large modulus. This includes
$E_\lambda$ (with $T=\omega-\lambda$ by part (b)), $E_{\lambda,M}$,
$E_\theta$ and every explicitly specified phase in the sources (error 40).
03's statements for arbitrary phases, 04's prescribed prime phases and
exponentials with character values in $1+\mm$ are covered only when the
phase in question has a resonance. A ``chosen phase'' of this kind does
not escape the obstruction; the galaxy replica is then $\zp(w)$ on the
resonant galaxy. The phase of \cref{gz:prop:resonance}(c) has no resonance,
so the theorem does not apply to it.

\emph{Reconciliation with 03.} The galaxies $\pm iT+\OO$ lie outside $\Th$.
03's ``single pole at $1$'' (\cref{gz:thm:zeta-tube}) is a statement on
$\Th$ and is untouched. Both statements express the $i\Pi$-invariance of
$\Exp$, and the wing periodicity of \cref{gz:prop:wing} is its shadow on the
right.

\begin{remark}[Scope \src{06}]\label{gz:rem:resonance-scope}
The theorem rules out a conjunction of explicit axioms, not every global
surcomplex zeta function. A phase without resonance, a different summation
rule, or a failure of coherence avoids its hypotheses. The phrase
``trivial zeros'' also needs care in a global theory. Extra zeros of a chosen
global sine are not automatically trivial zeros of a completed zeta. Gamma
poles and zeta zeros must cancel in a specified meromorphic construction
before a completed function and its RH can be formulated. On $\Th$,
\cref{gz:thm:xi-tube} is such a construction.
\end{remark}

\section{What weak global analyticity does not decide}\label{gz:sec:weak}

\begin{theorem}[Symmetric completions with different infinite zeros \src{05, 06}]\label{gz:thm:weak}
Let $T>\R$ and let $0<a<\tfrac12$ be ordinary, and put
$c_\pm=\tfrac12\pm iT$. In (a) assume moreover that $T$ is purely infinite
\src{06}; (b) holds for every positive infinite $T$ \src{05}.
\begin{enumerate}[label=(\alph*)]
\item \src{06} The class functions $F_0=\xip$ on $\OO$, $F_0=1$ off $\OO$, and
 $F_1=F_0$ except $F_1(s)=(s-c_\pm)^2-a^2$ on $\pm iT+\OO$, are locally
 Hahn-analytic at every point of $\SC$, agree with $\xip$ on $\OO$, and satisfy
 $F(1-s)=F(s)$ and $F(\bar s)=\overline{F(s)}$. $F_0$ has no infinite zeros.
 $F_1$ has exactly four, $c_\pm\pm a$, all simple and off the line.
\item \src{05} Let $F:\SC\to\SC$ be monadwise analytic (on every monad it has
 a strongly evaluable local power series) with the same two symmetries. Let
 $\mathcal Q=\{\tfrac12\pm a\pm iT\}$, let $U$ be the union of the four
 monads, and let $P_U=\prod_{b\in\mathcal Q}(s-b)$ on $U$ and $P_U=1$ off $U$. Then
 $\widetilde F=P_UF$ has the same symmetries and local property, agrees with
 $F$ wherever $\Im s$ is finite, and vanishes on $\mathcal Q$. It differs
 from $F$ unless $F$ vanishes at every point of $U$ where $P_U\ne1$.
\end{enumerate}
\end{theorem}
\begin{proof}
(a) Finite translation and conjugation preserve $\OO$, and every point has a
fine neighbourhood inside its galaxy, so local analyticity holds. Under
$s\mapsto1-s$ the two galaxies are exchanged and the centred coordinate
changes sign; conjugation also exchanges them. The zeros are $c_\pm\pm a$,
with real parts $\tfrac12\pm a$. (b) The four monads are disjoint: their
centres differ by a nonzero ordinary real part or an infinite imaginary part.
$\mathcal Q$ and $U$ are stable under $s\mapsto1-s$ and conjugation, so the
polynomial $\prod_{b\in\mathcal Q}(s-b)$ has real coefficients and is
invariant under $s\mapsto1-s$, as its degree is even and its root set is
reflection-stable. Hence $P_U$ has both symmetries. On each monad $P_U$ is
a polynomial or $1$, and $U$ contains no point of finite imaginary part.
\end{proof}
05's statement uses $a=\tfrac16$, with zeros at $\tfrac13\pm iT$ and
$\tfrac23\pm iT$, and says ``there is another such function''. That is false
when $F$ vanishes on $U$ (for example $F\equiv0$), hence the qualifier in
(b) (error 32). The functions are not proposed as zeta completions; no Euler
product or Gamma factorization is asserted \src{06}. They are not
counterexamples to RH, and they are not an obstruction to a more rigid global
category \src{05}. They show that finite agreement, local analyticity and
the two completed symmetries do not settle RH at infinite height. A global RH
must name its function and the extra axioms that select it.

\subsection{Shortcuts that do not work, and several infinite RH statements}
The Taylor series of $\xi$ at $0$ cannot be evaluated at an infinite $s$:
infinitely many nonzero standard coefficients multiply powers of decreasing
valuation, so the support condition fails \src{05}. A Hadamard product over
the ordinary zeros is not automatically strongly summable at infinite $s$.
Transferring an asymptotic formula while dropping its remainder leaves room
for a flat term, and \cref{gz:thm:strip} shows that an exact identity can
force such a term \src{04, 05}.

For a specified global or partial function one can ask for the exact-line
statement $\Re\rho=\tfrac12$ at all nontrivial zeros, for the standard-part
statement $\st\Re\rho=\tfrac12$ (which allows $\tfrac12+\eps+iT$), or for an
infinite-height-only statement. These are different \src{05}. Density
theorems need not exclude even rare zeros at a fixed distance from the line,
and excluding infinite-height exceptions alone does not exclude a bounded set
of ordinary exceptions. RH equivalences in this report use the full
universal statement. ``Trivial zeros'' depend on the construction as well.
The finite lift has exactly $-2,-4,\dots$. A tube phase has $2\cI(\Psi_\chi)_{<0}$,
which is $2\Oz_{<0}$ for $\Exp$ (\cref{gz:cor:canonical-zeros}). An
elementary model has the transferred negative even \emph{internal} integers,
which are not automatically $2\Oz_{<0}$ (\cref{gz:sec:transported-divisors}).
A finite-step recurrence cannot reach a negative infinite integer
\src{03, 05}.

\part{Elementary extensions and transfer}\label{gz:part:transfer}

\section{Transported structures}\label{gz:sec:transfer}

\subsection{The structure that is transferred}\label{gz:sec:structure}
Start from a set-sized first-order expansion $\mathfrak A$ of the real field.
It names the real and imaginary coordinate functions of the entire functions
$\xi(s)$, $(s-1)\zeta(s)$ and $1/\Gamma(s)$, with complex inputs as pairs of
reals, and meromorphic functions are recovered by division off the zero set
of a denominator or by domain predicates. It names the real exponential,
$\sin$ and $\cos$, \emph{a constant for every real number}, and, as needed,
counting functions ($N$, $N_0$, $N_0^{\rm s}$, $N(\sigma,\cdot)$), an integer
predicate with finite sums, Chebyshev's $\psi$, and the heat family $H$
\src{04, 05, 06}. The real constants are used as parameters in the
finite-zero and continuity arguments below. 06 names only the functions and
leaves this implicit; Ehrlich and Kaplan's $T_{\rm trig,exp}$ does the same
(error 43). A proper elementary extension $\mathfrak M^*\succ\mathfrak A$ exists,
as a nonprincipal ultrapower \src{04, 06} or by compactness \src{05}. Its
real field ${}^*\R$ is real closed and contains infinite elements, and it
carries transferred functions ${}^*\xi,{}^*\zeta,{}^*\Gamma,{}^*\exp,\dots$.
Transfer (Łoś) applies to first-order formulas in the named symbols. No
transfer principle for all surreals is assumed.

\begin{proposition}[Realization inside the surreals \src{04, 05, 06}]\label{gz:prop:embedding}
Every set-sized real closed ordered field $F\supseteq\R$ has an ordered-field
embedding into $\No$ fixing $\R$. Hence $F[i]$ embeds in $\SC$, and every
named function of $\mathfrak M^*$ can be transported to the image.
\end{proposition}
\begin{proof}
Well-order a transcendence basis of $F$ over $\R$. At a stage with real
closed domain $E$ and embedding $j$, the next basis element $\alpha$
determines the set-sized sets $L=\{j(e):e<\alpha\}$ and $R=\{j(e):e>\alpha\}$.
A surreal $\{L\mid R\}$ lies strictly between them, and it is not in $j(E)$,
since every element of $j(E)$ lies on one side of its own cut. As $j(E)$ is
real closed, it is transcendental over $j(E)$. The signs of polynomials over
a real closed field are determined by the cut, so this extends $j$ to an
ordered embedding of $E(\alpha)$ and then of its real closure inside $\No$.
At limit stages take unions, which are real closed. At the end the
embedding is defined on the real closure of the rational-function field of
the whole basis. $F$ is algebraic over that real closed field, hence equal
to it, so the embedding is defined on all of $F$.
\end{proof}
This is the standard universality of $\No$ \cite{Gonshor,vdDE}, in a relative
form. It supplies surcomplex-valued models at infinite height, but no unique
function on all of $\SC$.

\begin{warning}[Transported operations are extra structure \src{04, 05, 06}]\label{gz:warn:transported}
An ordered-field embedding need not preserve the Gonshor exponential, strong
Hahn sums, restricted analytic functions, a scalar derivation, or $\Exp$. The
transported ${}^*\exp$ is not automatically $\exp_{\No}$. ${}^*\zeta$ must not be
identified with $\zp$ or $\zD$ where their domains overlap, and transported
values are not automatically the strong Taylor--Stirling values. Stars must
not be erased at infinite arguments. The two-frequency test
(\cref{gz:prop:twofreq}) is a first-order sentence true in every elementary
extension and false for $\Exp$, so the transported full phase is not $\Exp$
on any domain containing a resonance. Transferred finite-order Taylor
estimates allow, in a non-Archimedean field, discrepancies smaller than every
power of an infinitesimal. Equality with a specified strong Hahn sum is an
additional compatibility requirement. It would need an embedding preserving
a restricted-analytic structure, and the relevant results
\cite{vdDE,CE,EK} have hypotheses that would have to be checked for the chosen
language.
\end{warning}

\subsection{Divisors in transported models}\label{gz:sec:transported-divisors}
Include $G=1/\Gamma$ and the integer sort. The classical facts
$G(z)=0\iff z\in-\N$ and $G'(-n)=(-1)^nn!$ transfer. So ${}^*\Gamma$ has no
zeros, and its poles are exactly $-{}^*\N_{\ge0}$, simple, with
$\Res({}^*\Gamma,-n)=(-1)^n/{}^*n!$, which is infinitesimal at infinite $n$.
Likewise ${}^*\zeta(-2n)=0$ for every positive internal $n$, with the
transferred derivative formula \eqref{gz:eq:trivialslope}, and ${}^*\xi$ cancels
these zeros and poles. Recurrence, reflection and the functional equation
hold wherever the transported expressions are defined. A Gauss product with
an internal number of factors needs the internal product operation, not a
Hahn product \src{04}.

\emph{Reconciliation.} These poles are at the internal integers, which depend
on the ultrafilter and the embedding. They are not every negative omnific
integer, and ${}^*\N$ is not $\Oz$ \src{04, 05, 06}. The tube construction
(\cref{gz:thm:gamma-tube}) puts the poles on the period group of a phase,
$\Oz_{\le0}$ for $\Exp$. The two are different constructions. Both choose the
alternative that \cref{gz:thm:forced} allows, extra infinite poles \src{06};
inside the model the transported sine has zero set $\pi\,{}^*\Z$. The
positive-infinite Stirling chart and the finite lift do not by themselves
specify an internal lattice of negative poles \src{04}.

\subsection{The transferred Riemann hypothesis}
\begin{theorem}[Transferred RH \src{04, 05, 06}]\label{gz:thm:RHstar}
Let $F={}^*\R$ with the transported structure.
\begin{enumerate}[label=(\alph*)]
\item $\forall s\in F[i]\ ({}^*\xi(s)=0\Rightarrow\Re s=\tfrac12)$ is equivalent
 to classical RH.
\item The same with the added conclusion ${}^*\xi'(s)\ne0$ is equivalent to RH
 together with simplicity of every nontrivial zero \src{04}.
\item The model has zeros of ${}^*\xi$ at infinite heights.
\item Every finite zero of ${}^*\xi$, of ${}^*\zeta$, or of any transported
 nonzero ordinary entire function is an exact ordinary zero.
\end{enumerate}
\end{theorem}
\begin{proof}
(a) and (b) are transfers of first-order sentences about two real
coordinates, in both directions. (c) ``For every real $T$ there is a zero with
imaginary part greater than $T$'' is first order and transfers to infinite
$T\in F$. (d) Let $z$ be finite with standard part $a$. If $f(a)\ne0$, an
ordinary disk about $a$ is zero-free, and that statement transfers. If
$f(a)=0$, an ordinary disk contains only the zero $a$, and ``every zero in
this disk equals $a$'' transfers, with $a$ and the radius named by
constants. Alternatively, an ordinary disk of radius $R$ contains a finite
list of zeros $\rho_1,\dots,\rho_m$, and ``every zero in the disk is one of
them'' transfers. For zeta use the pole predicate as well.
\end{proof}
This is a global RH at finite and infinite heights, but not a new conjecture.
Its extra instances follow from transfer, and a proof of (a) independent of
classical RH would prove classical RH \src{04, 05}. The zeros of the model
are cofinal in the set-sized field $F$, not in $\No$ \src{05}. By contrast,
an independently constructed global surcomplex zeta could have a genuinely
stronger RH unless its connection with the classical function were proved.
Finite-zero agreement with $\xip$ does not give equality of finite values
\src{04, 06}.

\subsection{Internal counting and quantifiers}\label{gz:sec:quantifiers}
${}^*N(T)$ is a named internal function. At infinite $T$ it is an infinite
internal integer, not the external cardinality of the set of ordinary zeros
below $T$, which is countable \src{04, 05, 06}. It is not a Hahn sum and not
a proper-class contour integral. An estimate $A(T)=\bigO(B(T))$ transfers
only in quantified form, with ordinary witnesses $C,T_0$ such that
$|A(T)|\le CB(T)$ for $T\ge T_0$. A bound $T^{\alpha+o(1)}$ transfers as ``for
each ordinary $\eta>0$ there are ordinary $C,T_0$ with
$A(T)\le CT^{\alpha+\eta}$''. The exponent $o(1)$ cannot be suppressed at
infinite $T$, since an infinitesimal exponent times $\log T$ can be large
\src{04, 05}. A theorem for each fixed ordinary $\sigma$ or $\eta$ is not
uniform in an infinitesimal perturbation of the parameter unless its
constants are tracked. A bounded statement transfers with its ordinary
endpoints, and $N(T)$ must not be replaced by $N(\st T)$ at an endpoint
without a separate argument \src{05}.

\section{Classical results at the new scales}\label{gz:sec:partial}

This section gives selected benchmarks with their exact scopes and statuses.
It does not claim that any constant is the best currently known, and it
improves no classical bound. Each classical result has one meaning for the
finite lift (the zero set is unchanged, by \cref{gz:thm:divisor}) and another
in an elementary model (quantified transfer), and the proofs are not rerun
over Hahn coefficients \src{04, 05, 06}.

\subsection{Zero-free regions}
By \cref{gz:cor:halofree}, every ordinary zero-free region lifts exactly to its
halo, boundary included. Mossinghoff, Trudgian and Yang prove zero-freeness
in
\[
\sigma\ge1-\frac1{55.241(\log|t|)^{2/3}(\log\log|t|)^{1/3}}\ (|t|\ge3),\qquad
\sigma\ge1-\frac1{5.558691\log|t|}\ (|t|\ge2),
\]
and with $48.1588$ in place of $55.241$ for all sufficiently large $|t|$
\cite{MTY} \src{04, 05, 06}.

\begin{proposition}[Zero-free layers at infinite height \src{04, 05, 06}]\label{gz:prop:mty}
In an elementary model containing zeta and the logarithm, both explicit
regions exclude zeros of ${}^*\zeta$ for internal $\sigma,t$ with the
transported ${}^*\log$. At every infinite $|T|$ the asymptotic region does
too. The width $w(T)=1/(55.241({}^*\log|T|)^{2/3}({}^*\log{}^*\log|T|)^{1/3})$
is a positive infinitesimal, and the layer $\sigma\ge1-w(T)$ is exactly
zero-free.
\end{proposition}
\begin{proof}
Each region is a first-order implication with ordinary constants. The
asymptotic one has an ordinary threshold, exceeded by every infinite $|T|$.
\end{proof}
This is stronger than $\st\sigma\le1$ and far from $\sigma=\tfrac12$. It holds
for the transported structure. Substituting the Gonshor logarithm after an
arbitrary embedding would need a compatibility that the embedding does not
supply. At positive infinite real part, $\zD$ is zero-free on $\Tp$ without
transfer (\cref{gz:thm:euler}), a different sector.

\subsection{Verified heights}\label{gz:sec:verified}
Platt and Trudgian verified RH rigorously for all zeros with
$0<\Im\rho\le T_{\rm PT}=3\cdot10^{12}$, using interval arithmetic \cite{PT}. The
published abstract adds that all these zeros are simple; arXiv v1 lacks that
sentence \src{05}. 06 declines to infer simplicity from the verification
statement it quotes. 06's argument does not need simplicity.
\Cref{gz:cor:verified} (from 05) does use it, through the published
abstract's statement, whose computation this merge did not recheck.
$T_{\rm PT}$ is a height, not a number of
zeros, and it is not claimed to be the latest record \src{04, 05, 06}. The
canonical finite statement is immediate: every finite zero of $\zp$ with
ordinate in $(0,T_{\rm PT}]$ lies on the line.

\begin{corollary}[Protection in verified monads \src{05}]\label{gz:cor:verified}
Assume the published simplicity statement of \cite{PT}. For every
$F\in\DJs$ and every real infinitesimal $\tau$, every finite zero
$s$ of $F_\tau$ with $0<\Im\st s\le3\cdot10^{12}$ lies on $\Re s=\tfrac12$.
Each verified zero has exactly one deformed zero in its monad.
\end{corollary}
\begin{proof}
$\st s$ is a zero of $\xi$ in the verified range, simple and on the line by
the published statement. Apply \cref{gz:thm:protection}.
\end{proof}
The range refers to standard parts: a zero at a boundary ordinate may move
infinitesimally above or below it. The corollary says nothing at infinite
height. Platt and Trudgian also give
$|\psi(x)-x|\le\frac{\sqrt x}{8\pi}\log^2x$ for $59<x\le2.169\cdot10^{25}$
\src{04}. This transfers with the same ordinary range. It does not become an
all-infinite-$x$ theorem, and the raw sum over ordinary primes below an
infinite surreal is not the transported ${}^*\psi$.

\subsection{Counting and Hardy's theorem}\label{gz:sec:counting}
The Riemann--von Mangoldt formula
$N(T)=\frac T{2\pi}\log\frac T{2\pi}-\frac T{2\pi}+\bigO(\log T)$
\cite{Titch,DLMFzeros} holds exactly for the finite lift at every ordinary
height. In an elementary model there is an ordinary $C$ with
$|{}^*N(T)-\frac T{2\pi}{}^*\log\frac T{2\pi}+\frac T{2\pi}|\le C\,{}^*\log T$ at all
large internal $T$, including every infinite $T$. Hence
${}^*N(T)/(\frac T{2\pi}{}^*\log\frac T{2\pi})\in1+\mm$, equivalently
$\st({}^*N(T)/(T\,{}^*\log T))=1/(2\pi)$ \src{04, 05, 06}. (05 cites MTY for the
formula, but MTY only quotes an explicit version; Titchmarsh is the primary
source.) 04 reports that Rodgers and Tao record the version in the doubled
heat variable; this was not rechecked. Hardy's theorem, in the form ``for
every real bound there is a critical-line zero above it'', transfers to zeros
above every internal bound. For the finite lift it gives only the original
countable set \src{04, 05}.

\subsection{Proportions}\label{gz:sec:proportions}
Let $\kappa_0=\liminf N_0(T)/N(T)$, where $N_0$ counts critical-line zeros with
multiplicity, and $\kappa_{\rm s}=\liminf N_0^{\rm s}(T)/N(T)$, which counts
simple critical-line zeros. These are different quantities.
\begin{itemize}
\item \emph{Published} \cite{PRZZ}: $\kappa_0\ge0.417293$ and
 $\kappa_{\rm s}\ge0.407511$. ``More than five-twelfths'' concerns zeros on the
 line, not simple zeros on the line. 05 uses only $\kappa_0>5/12$, which is
 consistent, since $0.417293>5/12>0.407511$. Whether PRZZ's $N_0$ counts
 multiplicities was not checked by 05.
\item \emph{Preprints, August and September 2026.} Alp\"oge and Furman
 \cite{AF}: at least $2/3$ of the zeros are simple and on the line, and at
 least $5/6$ are distinct. With a Montgomery--Taylor window the constants
 become $0.6725$ and $0.8362$. Lamzouri \cite{Lamzouri}, Theorem 1.1:
 $\kappa_{\rm s}\ge C_0=\tfrac32-\tfrac1{\sqrt2}\cot\tfrac1{\sqrt2}=0.6725007\ldots$,
 distinct-zero proportion at least $(1+C_0)/2=0.8362503\ldots$, and $88.76\%$
 simple or on the line. The closed form $C_0$, the Montgomery--Taylor
 constant, is stated in Lamzouri's Theorem 1.1; AF's $0.6725$ is their
 Montgomery--Taylor-window variant (a precision added to 04, error 28; 06
 was correct in substance). These are preprints. Their proofs were not verified by any
 source or by this merge, and AF's reported formal certificates were not
 run. Every conclusion below is parametric in the accepted bound.
\end{itemize}
For the finite lift both sets of estimates apply unchanged.
\Cref{gz:thm:protected} gives an unconditional protected proportion at least
$0.407511$ for symmetric deformations with the published input, and at least
$C_0$ with the preprint input. Using $0.417293$ there would be wrong: a
multiple critical-line zero may split off the line \src{04, 06}.

\begin{corollary}[Stable proportions \src{04, 05, 06}]\label{gz:cor:stableproportion}
If $\liminf N_0^{\rm s}(T)/N(T)\ge c_*$ for an ordinary $c_*$, the
centre-indexed proportion of uniformly stable centres
(\cref{gz:sec:protected}) is at least $c_*$. In any elementary model naming
the counts, for every ordinary $c<c_*$ and every infinite $T$,
${}^*N_0^{\rm s}(T)\ge c\,{}^*N(T)$; hence
$\st({}^*N_0^{\rm s}(T)/{}^*N(T))\ge c_*$. The same holds for $N_0$ under
$\liminf N_0(T)/N(T)\ge c_*$.
\end{corollary}
\begin{proof}
The halo statement is \cref{gz:thm:protection} at every simple
critical-line centre. The $\liminf$ gives an ordinary $T_c$ with
$N_0^{\rm s}(T)\ge cN(T)$ for $T\ge T_c$, and this transfers. Taking every
$c<c_*$ gives the standard-part form.
\end{proof}
The conclusion is not the exact endpoint inequality with $c=c_*$. A
positive proportion does not show that every zero is on the line or simple,
and an infinite exceptional count of smaller order remains possible. With
the unbounded total count it does give critical-line zeros at infinite
heights in the model \src{04, 05, 06}.

\subsection{Zero density}\label{gz:sec:density}
Ingham's exponent is $A_{\rm I}(\sigma)=3(1-\sigma)/(2-\sigma)$. Guth and
Maynard's Theorem 1.2 gives $A_{\rm GM}(\sigma)=15(1-\sigma)/(3+5\sigma)$,
and with Ingham in the complementary range their (1.4) gives
$N(\sigma,T)\le T^{30(1-\sigma)/13+o(1)}$ \cite{GM}. The two exponents cross
at $\sigma=7/10$, and $A_{\rm GM}(3/4)=5/9$. The paper is published in
\emph{Ann.\ of Math.}\ 203(2) (2026). 04 cites only arXiv v2 and 06 gives no
status (errors 29, 46). Guth--Maynard count zeros with $|\Im\rho|\le T$, as
04's $N^\pm_{\ge\sigma}$ does. The one-sided count $0<\Im\rho\le T$ of 05 and
06 is smaller, so the same bounds hold for it (05's remark, added to 06's
text, error 47).

\begin{proposition}[Fixed-parameter density transfer \src{04, 05, 06}]\label{gz:prop:density}
Fix ordinary $\tfrac12<\sigma<1$ and $\eta>0$. For each applicable exponent
there are ordinary constants with ${}^*N(\sigma,T)\le C_{\sigma,\eta}T^{A(\sigma)+\eta}$
at every infinite internal $T$, with the transported power. Unconditionally,
${}^*N(\sigma,T)/{}^*N(T)$ is infinitesimal for every fixed ordinary
$\sigma>\tfrac12$ and every infinite $T$.
\end{proposition}
\begin{proof}
The first claim is the quantifier rule of \cref{gz:sec:quantifiers}. For the
second, $A_{\rm I}(\sigma)<1$ for $\sigma>\tfrac12$. Choose $\eta$ with
$A_{\rm I}(\sigma)+\eta<1$ and divide by the transferred counting formula:
the ratio is at most $CT^{A_{\rm I}+\eta-1}/{}^*\log T\in\mm$. For
$\sigma\ge1$ the count is zero.
\end{proof}
By reflection and conjugation the same holds to the left of the line, so the
fraction of zeros outside every fixed ordinary-width neighbourhood of the line
is infinitesimal. This is not a proof that all zeros lie on the line. An
infinite population can have infinitesimal relative density, and
fixed-$\sigma$ bounds do not reach $\sigma=\tfrac12+\eps$ with infinitesimal
$\eps$. For the finite lift the count is exactly the ordinary one, and no
density statement for an arbitrary additive Hahn perturbation follows from
its standard-part coefficient \src{04}.

\subsection{Primes and finite positivity tests}\label{gz:sec:positivity}
RH is equivalent to $\psi(x)=x+\bigO(x^{1/2}\log^2x)$ \cite{Titch}. Its
quantified consequences transfer once $\psi$ and the arithmetic are named.
In a bare surreal field there is no specified function ``number of primes
below an infinite $x$''. Ordinary integers, omnific integers and transferred
integers are different data, and a sum over unspecified ``surreal primes''
defines nothing \src{06}. Classical RH equivalences involving
prime-counting errors, derivatives of zeta or positivity conditions
transfer only after every function, sum and index set they use has been
named in the structure. They do not become statements about a separately
chosen Hahn-valued prime-counting function. A pair-correlation statement
likewise needs its counting measure, normalization, test functions and
limit quantifiers; it cannot be replaced by an unspecified ``surreal
distribution of zeros'' \src{05}.

\begin{proposition}[Finite Hermitian tests \src{04, 06}]\label{gz:prop:hermitian}
An ordinary Hermitian matrix $A$ is positive semidefinite (definite) over $\C$
if and only if $u^*Au\ge0$ ($>0$ for $u\ne0$) for all $u\in\SC^n$. A positive
semidefinite ordinary Hermitian form stays so after algebraic scalar
extension to $\SC$.
\end{proposition}
\begin{proof}
Diagonalize $A=UDU^*$ over $\C$. Then $u^*Au=\sum d_j|(U^*u)_j|^2$ in the
ordered real field. A negative $d_j$ already fails at an ordinary vector.
Every vector of an algebraic scalar extension is a finite combination of
ordinary vectors, whose Gram matrix is an ordinary positive semidefinite
matrix.
\end{proof}

\begin{proposition}[A real margin protects positivity \src{04}]\label{gz:prop:margin}
If $A$ is positive definite and $\Delta$ is Hermitian with infinitesimal
entries, then $A+\Delta$ is positive definite over $\SC$. This fails for
positive semidefinite $A$: take $A=\operatorname{diag}(0,1)$ and
$\Delta=\operatorname{diag}(-\eta,0)$.
\end{proposition}
\begin{proof}
With $u^*Au\ge a\sum|u_j|^2$ for an ordinary $a>0$ and $e=\max|\Delta_{jk}|$,
Cauchy--Schwarz gives $|u^*\Delta u|\le de\sum|u_j|^2$, and $de<a$.
\end{proof}
Surreal coefficients in a fixed finite family of Weil-type tests therefore
add no condition. A test over all Hahn-supported infinite families needs a
summation theorem, and a test over all infinitesimal deformations needs
control of null directions \src{04, 06}. RH criteria built from fixed-order
derivatives of $\xi$ at ordinary points are preserved exactly by lifting. A
hyperinteger index requires transfer, and a new Hahn test family is a third
operation. Keeping the three apart prevents false ``stronger RH''
conclusions \src{04}.

\section{Heat flow at infinitesimal time}\label{gz:sec:heat}

Use the normalization of Rodgers and Tao \cite{RT} \src{04, 05, 06}:
\begin{gather*}
H_t(z)=\int_0^\infty e^{tu^2}\Phi_{\rm RT}(u)\cos(zu)\,du,\qquad
\Phi_{\rm RT}(u)=\sum_{n\ge1}(2\pi^2n^4e^{9u}-3\pi n^2e^{5u})e^{-\pi n^2e^{4u}},\\
H_0(z)=\tfrac18\,\xi\bigl(\tfrac12+\tfrac{iz}2\bigr)=\tfrac18\,\Xi(z/2).
\end{gather*}
The last form reconciles 04's stability variable $\Xi(w)=\xi(\tfrac12+iw)$ with
the doubled heat variable; there is no other factor of $2$. The de
Bruijn--Newman constant $\Lambda_{\rm DN}$ is characterized by: all zeros of
$H_t$ are real if and only if $t\ge\Lambda_{\rm DN}$. RH is equivalent to
$\Lambda_{\rm DN}\le0$, and Rodgers and Tao proved $\Lambda_{\rm DN}\ge0$. So
every negative ordinary $t$ has a nonreal zero. (The upper bound
$\Lambda_{\rm DN}\le0.2$ is not used.) The
super-exponential decay of $\Phi_{\rm RT}$ makes $H$ jointly entire in $(t,z)\in\C^2$,
since each derivative of $e^{tu^2}\cos(zu)$ is bounded on compacta by a
polynomial times $e^{Cu^2+Cu}$. Differentiation under the ordinary integral
gives $\partial_tH=-\partial_z^2H$ and
$H_t=\sum_n\frac{(-t)^n}{n!}H_0^{(2n)}$. No surreal integral is used: these are
ordinary germs before lifting or transfer. In the variable
$s=\tfrac12+iz/2$ the formal family $8H_t$ is a deformation in $\DJs$ with both
symmetries, since each $H_0^{(2n)}$ is real and even, and $t$ is real
\src{05}.

\begin{lemma}[Local real-root persistence \src{04, 05, 06}]\label{gz:lem:heatpersist}
Let $a$ be a simple real zero of $H_0$. There are ordinary $r,\delta>0$ such
that for every real $|t|<\delta$, $H_t$ has exactly one zero in $|z-a|<r$, and
it is real, with $z(t)=a+t\frac{H_0''(a)}{H_0'(a)}+\bigO(t^2)$. The same holds
after canonical infinitesimal evaluation and after transfer.
\end{lemma}
\begin{proof}
The analytic implicit-function theorem at $(0,a)$ gives one root with a unit
factorization. Shrink $r,\delta$ until it is the only root in the disk. For
real $t$, conjugation preserves the disk and the function, so uniqueness
makes the root real. Differentiating $H_t(z(t))=0$ with the heat equation
gives $z'(0)=-\partial_tH/\partial_zH=H_0''(a)/H_0'(a)$. Hahn evaluation
preserves the factorization. The existence, uniqueness and reality
statement, with $a,r,\delta$ named by constants, is first order and
transfers.
\end{proof}

\begin{theorem}[Escape at negative infinitesimal time \src{04, 05, 06}]\label{gz:thm:escape}
Let $\mathfrak M^*$ be an elementary extension naming $H$ as in
\cref{gz:sec:structure}, with real field $F$. The hypotheses are stated
clause by clause.
\begin{enumerate}[label=(\roman*)]
\item \emph{Assume RH and simplicity.} For every real infinitesimal $\tau$, of
 either sign, every finite zero of the canonical joint Taylor lift
 $H^\#(\tau,\cdot)$ is real and simple.
\item \emph{Unconditionally,} for every negative infinitesimal $\tau\in F$, the
 transferred ${}^*H_\tau$ has a nonreal zero.
\item \emph{Assume RH and simplicity.} For every negative infinitesimal
 $\tau\in F$, every nonreal zero of ${}^*H_\tau$ has infinite modulus.
\item \emph{Unconditionally,} a finite nonreal zero of ${}^*H_\tau$ at an
 infinitesimal $\tau$ has as standard part either a nonreal zero of $H_0$ or
 a multiple real zero of $H_0$.
\end{enumerate}
At a simple real zero $c$, $c(-\eta)=c-\eta H_0''(c)/H_0'(c)+\bigO(\eta^2)$.
\end{theorem}
\begin{proof}
(i) Under the assumptions every zero of $H_0$ is real and simple. A finite
zero of the joint lift has as standard part a zero $c$ of $H_0$. In the monad
of $c$ the evaluated unit factorization of \cref{gz:lem:heatpersist} leaves
exactly one zero, which is real and simple.
(ii) The ordinary sentence $\forall t<0\ \exists z\,(H_t(z)=0\wedge\Im z\ne0)$
follows from $\Lambda_{\rm DN}\ge0$ and the definition of $\Lambda_{\rm DN}$, with no
hypothesis, and it transfers.
(iv) If $z$ is finite with $a=\st z$, then ordinary continuity of $H$ near
$(0,a)$, with each ordinary tolerance transferred, gives $H_0(a)=0$. If $a$
were a simple real zero, the transferred \cref{gz:lem:heatpersist} would
make $z$ real.
(iii) Under RH and simplicity every zero of $H_0$ is simple and real, so (iv)
leaves no possible standard part for a finite nonreal zero.
\end{proof}
04 states (ii) under the RH hypothesis of its theorem, although (ii) needs
no hypothesis (error 27). 06 places the canonical-lift clause (i) after its
conditional statement in a way that reads as unconditional; it holds under
RH and simplicity (error 44). 05 notes that only its part (b), which is
(iii) here, uses the hypotheses, and it isolates (i) as a consequence of the
deformation theorem. 06 states (iv).

The theorem explains why infinitesimal stability at every ordinary zero is
compatible with $\Lambda_{\rm DN}\ge0$. Each fixed simple zero has a real time
interval protecting it in a fixed disk, but no common interval protects all
zeros at all heights. An infinitesimal time lies inside every individual
interval, so the nonreal zeros forced by the global theorem must escape every
finite disk. It proves infinite modulus only: no direction, rate, leading
monomial or asymptotic scale of the escaping zeros is given \src{04, 05, 06}.
Simplicity matters, because a multiple zero can split nonreally inside its own
finite monad (\cref{gz:cor:scale}). Clauses (i) and (iii) are conditional,
and the theorem proves neither assumption.

\part{Synthesis}\label{gz:part:synthesis}

\section{Synthesis and open questions}\label{gz:sec:synthesis}

There is no single unqualified ``RH over the surcomplex numbers''. There are
precise statements attached to precise extensions.
\begin{itemize}
\item The finite RH (\cref{gz:thm:finiteRH}), the horizontal RH on $\Th$ for
 every phase (\cref{gz:cor:horizontalRH}) and the transferred RH
 (\cref{gz:thm:RHstar}(a)) are each equivalent to classical RH. The first two
 hold for the reason that the divisors are rigid, and the third because of
 elementarity.
\item The robust finite RH (\cref{gz:thm:stability}) and the transferred RH
 with nonvanishing derivative (\cref{gz:thm:RHstar}(b)) are each equivalent
 to RH together with simplicity.
\item For a separately chosen global completion on $\Vs$, RH is an additional
 assertion whose content depends on the construction. Weak axioms do not
 determine it (\cref{gz:thm:weak}). Natural combinations of a phase,
 coefficientwise summation and coherent continuation are inconsistent for
 every resonant phase (\cref{gz:thm:resonance}), and reflection conflicts
 with the Stirling modulus there except at two exceptional abscissae
 $\sigma_\pm(T)$ for each height, while $E\circ\LSt$ fails reflection at
 every point of $\Vs$ (\cref{gz:thm:strip}).
\end{itemize}
The finite theory is fully specified on its domain. The theory at positive
infinite real part is a second exact construction, with an injective
arithmetic algebra, exact Euler products, complete inverse fibres,
independence theorems and exact Stirling and Hurwitz identities. A named phase
extends Gamma, zeta and $\xi$ to the horizontal tube, with divisors that move
with the phase while the $\xi$ divisor stays classical. At infinite height
the global problem needs additional choices, and several natural ones are
inconsistent.

\begin{question}[Nonresonant global completion \src{02, 06}; re-scoped \merge]\label{gz:q:nonresonant}
Specify a phase with no resonance, a summation rule in a finite-real-part
region at infinite height, and a meromorphic continuation class. Can one
keep a single zeta pole, the completed functional equation and a prescribed
Gamma normalization? \Cref{gz:thm:resonance} excludes every phase with a
resonance, and in particular every phase with ordinary-circle character
values. 02 asks for a chosen continuation or summation mechanism near
infinite imaginary height compatible with the Gamma factors and phases, and
06 asked for a phase satisfying the two-frequency test.
\Cref{gz:prop:resonance}(c) supplies one, and whether a coherent galaxy zeta
exists for it is open.
\end{question}

\begin{question}[Combined independence]\label{gz:q:combined}
Are dilation-and-shift families $\zD^{(j)}(k(X+a))$ algebraically independent
over $E_X^\Gamma$, or over $E_X^\Gamma(\log X)$, which also contains
$\LSt(X)$? Neither \cref{gz:thm:indep-main} nor
\cref{gz:thm:gamma-zeta-independence} implies this (\cref{gz:sec:indep-compare}).
\end{question}

\begin{question}[Canonical versus transported values \src{04, 05, 06}]\label{gz:q:compat}
An elementary model has exactly the right finite zeros, but an arbitrary
embedding does not identify its values with Hahn lifts. Which hypotheses on
the embedding, or on a restricted analytic substructure, force equality, and
how far into infinite sectors does it extend? The restricted-analytic
results of \cite{vdDE,CE,EK} provide structure but do not settle this. A
first version \src{04}: characterize extensions that preserve restricted
analytic germs, the surreal exponential, the Dirichlet chart on $\Tp$ and a
global functional equation. Existence alone would be meaningful, and
uniqueness should not be assumed, in view of flat Gamma and Hurwitz
modifications.
\end{question}

\begin{question}[Quantitative escape \src{04, 06}]\label{gz:q:heat}
Under RH and simplicity, define in a transferred model the least, or an
infimum, scale at which nonreal zeros appear for $\tau=-\eta$. How does it
relate to $\eta$, to classical small gaps and to the collision mechanism?
\Cref{gz:thm:escape} proves only that every ordinary bound is exceeded.
\end{question}

\begin{question}[Selecting a phase \src{01, 02, 03}]\label{gz:q:preferred}
Which compatibility conditions select a preferred phase, or exclude all the
coefficient-functional phases, beyond the group law and the local Taylor law?
Candidates include compatibility with a scalar derivation, with automorphisms
or with simplicity structure. For which classes of infinite imaginary
arguments does $(E_\chi(-s\log n))_n$ have a provable joint support
certificate (\cref{gz:prop:wing,gz:prop:wingfail}), and is such a phase
compatible with a scalar-derivation chain rule \src{03}? Any proposed
uniqueness theorem must explain why it excludes the explicit families of
\cref{gz:prop:families} \src{01}. For the robustness questions of Ehrlich and
Kaplan, see \rl{trigonometry:per:sec:ek}.
\end{question}

\begin{question}[Arithmetic restrictions on deformations \src{04}]\label{gz:q:arith}
Which arithmetic restrictions on coherent perturbations of $\xi$ exclude
off-line splitting without assuming simplicity? Constant perturbations show
that symmetry alone is not enough. An Euler product, an explicit formula or
compatible Dirichlet data are candidates, and each must be formulated on a
domain where the operation exists.
\end{question}

\begin{question}[Extending the constructions \src{01, 02, 03}]\label{gz:q:extend}
(a) Inverse theorems for other Dirichlet series $1+a_NN^{-s}+\cdots$, with
weights $\log_N(n/N)$, complex leading coefficients, missing indices and
cancellations \src{01}. (b) Generalized Dirichlet series with well-ordered
real frequency supports, keeping a leading-form classification; arbitrary
frequency sets can destroy well-ordering or finite convolution fibres
\src{02}. (c) Dirichlet characters and Euler factors of finite degree, with a
completed functional equation compatible with their infinite divisors and
phase factors \src{03}. (d) Coefficient fields built from several Gamma or
Barnes scales, adjoining a new monomial only when its support coordinate is
separated from the arithmetic one \src{03}.
\end{question}

\begin{question}[Gamma on the vertical strip \src{05}]\label{gz:q:strip}
Is there a Gamma prescription on $\Vs$ with conjugation symmetry, reflection,
the recurrence and a flat modulus correction as in \cref{gz:thm:strip}(d),
together with a specified phase? At the two exceptional abscissae
$\sigma_\pm(T)$ the modulus test is silent, and nothing is claimed there
for a general prescription with Stirling moduli; $E\circ\LSt$ itself fails
reflection there too (\cref{gz:thm:strip}(e)), so such a prescription
cannot be $E\circ\LSt$.
\end{question}

\emph{Formalization.} No theorem of this report is formalized. Only imported
infrastructure has Lean-proved rows in \texttt{docs/FORMALIZATION.md}: the
substitution clauses of \rl{a:cor:complexsub}, near-one binomial powers, and
fine derivatives of evaluated ordinary power series. The sources propose similar
orders. First comes the support layer: Neumann's lemma, strong
substitution, evaluation of weighted series and the divided-difference unit
lemma. Then the valuation-unit and divisor lemmas, the local perturbation
factorization with the involution $J$, and the arithmetic algebra with
finite-prime projection and the Jacobian argument. Then the formal Bernoulli
identities for Stirling and Hurwitz. Last come the phase-period, reflection
and completed-divisor theorems as abstract consequences, and transfer in a
specified first-order language. The analytic implicit-function theorem, the
surreal exponential and every imported number-theoretic theorem should
remain named inputs, not be hidden in a computational certificate
\src{01, 02, 03, 04, 05, 06}.

\section{What this report does not claim}\label{gz:sec:nonclaims}

\subsection{For the report as a whole}
\begin{itemize}
\item It does not prove, disprove or advance the classical Riemann
 hypothesis, and it improves no classical numerical bound.
\item It constructs no canonical, unique or maximal global zeta, $\xi$ or Gamma
 function on all of $\No[i]$. Every global object here depends on a named
 phase, domain, summation rule and continuation class.
\item No theorem of the report is verified in Lean or any proof assistant; only
 the imported infrastructure noted above has Lean-proved rows. The report is
 unrefereed and AI-assisted.
\item Priority is not certified for any result. ``Derived theorem'' and
 \merge{} describe proof status, not historical novelty. No named published
 open problem is claimed to be solved.
\item The 2026 proportion results are cited as preprints, not verified, and
 nothing local depends on them.
\item The finite checks of the six suites and of the merge verify finite
 algebra and ordinary numerics. They do not verify strong summability, class
 quantification, independence, continuation, RH or any cited paper.
\item The repository audits of the sources were targeted, and so was the merge's
 comparison with the collection. An empty search is not evidence of absence.
\end{itemize}

\subsection{Non-claims of the sources}
Every limitation, disclaimer and priority caveat stated by a source is kept,
grouped by source. Symbols are renamed as in \cref{gz:tab:rename}.

\paragraph{01.}
\begin{itemize}
\item Unrefereed, not verified in Lean, priority not certified.
\item The Taylor--Stirling baseline and the finite-phase locus are
 repository precedents; the finite lifts and the Stirling extension have
 precedents (Costin--Ehrlich, Berarducci--Mantova).
\item Lagrange inversion is classical and inverse zeta has a literature
 (Kawalec). Only the package with its surreal domain and support assertions
 is claimed.
\item The lift is canonical only within germ evaluation, not by uniqueness
 among finely differentiable functions.
\item Divisor rigidity does not apply to Hahn-coefficient deformations.
\item The finite-halo RH adds nothing toward RH and says nothing at infinite
 height.
\item The finite-phase criterion is not a maximal-domain or impossibility
 theorem.
\item Stirling uniqueness holds only among normalized formal series, and flat
 modifications are not excluded.
\item $E_\lambda$ is not asserted canonical or preferred, and its fine
 derivative is not the Berarducci--Mantova derivation.
\item 01's own construction gives no Gamma reflection across the negative
 infinite half-plane. The tube construction of 03 now supplies one under a
 named phase, and the non-claim stays true of 01's construction.
\item The Dirichlet embedding is not a classical analytic-continuation claim.
\item $\zD$ and $\zp$ are glued only on disjoint domains, with no
 fine-topological identification.
\item The fibre theorem is exhaustive only in $\Tp$, and there is no
 covering-map claim.
\item The integer-power sector is a projection of the formal $u$-expansion,
 not of the evaluated normal form.
\item The left prescription $Z_{-,\lambda}$ is one-sided and not by itself a
 proved global functional equation.
\item No contour integral, proper-class sum or transfinite recurrence is
 used, and Gauss indices are ordinary.
\item The normalization at infinity is not forced by the finite function.
\item Numerics are ordinary-complex illustrations, not error bounds.
\item The Lean plan is a dependency plan only.
\item The repository audit was targeted.
\end{itemize}

\paragraph{02.}
\begin{itemize}
\item No unique Gamma or zeta on $\No[i]$: no canonical global continuation, no
 unique phase at infinite imaginary part, no maximal domain.
\item No zeta and no functional equation in the infinite-height strip.
\item The $\xip$ statement is equivalent to RH, not a proof.
\item Independence concerns infinite values, not $\zeta(2),\zeta(3),\dots$,
 and standard parts destroy the information.
\item Derivatives are coordinate derivatives. No mean value or identity theorem
 across fine components is used.
\item The finite lift is canonical only as a germ lift.
\item The injectivity of $\DD_s$ is non-effective.
\item Gamma is not defined at negative infinite arguments by 02.
\item $\zH$ is not a rearrangement of $\sum(A+n)^{-s}$. The Lerch identity does
 not exclude flat modifications, and $-1/12$ is not the strong sum
 $1+2+3+\cdots$.
\item The $E_\theta$ are exponential homomorphisms, not preferred ones. The
 list of phase prescriptions is not exhaustive, and $\Om$ is not claimed
 maximal.
\item $\zeta^\Wing_{\chi_\theta}$ depends on the phase and is not a canonical
 continuation.
\item Pong and the classical literature are acknowledged. No novelty of the
 functional-independence principle and no priority are claimed. The
 repository reading was targeted to three README files; the gamma-functions
 article source was not reviewed.
\item No proof assistant was used. The $189$ checks are finite identities.
 The normal-form and exponential facts are imported inputs.
\item All supports are set-sized, and no basis or choice is used for
 proper-class vector spaces. The phases $E_\theta$ are stated for a given
 real-linear $\theta$; the explicit ones are the coefficient functionals
 $\theta=r\,c_M$, which need no choice.
\end{itemize}

\paragraph{03.}
\begin{itemize}
\item No RH and no resolution of a named problem. Unrefereed, no proof
 assistant, priority not certified.
\item No Gamma or zeta at arbitrary infinite imaginary height, and no value at
 $\tfrac12+i\omega$.
\item Zeta on $\Tm$ is defined by the functional equation, and the exact
 Stirling sum excludes flat terms by definition. Neither is a continuation
 or uniqueness theorem.
\item \Cref{gz:prop:tube-unique} is not a surreal Bohr--Mollerup theorem.
\item No connected continuation between $\Tm$, $\OO$ and $\Tp$, no global
 identity theorem, no single entire series.
\item The Gamma values are not claimed mutually independent. Coefficients from
 all of $\SC$ are not allowed, nor is a Hahn completion containing the zeta
 series.
\item No compatibility with the Berarducci--Mantova derivation. A phase that
 commutes with a scalar derivation is flagged as an extra, possibly
 restrictive, hypothesis.
\item Hurwitz uniqueness holds only in $A^{1-s}\C(s)[[A^{-1}]]$; $-1/12$ is not
 a strong sum; flat Hurwitz sectors are not excluded.
\item The checks do not verify summability, quantification, independence or
 novelty.
\item The canonical phase is a convention, not forced by reflection. Every
 phase has infinite poles.
\end{itemize}

\paragraph{04.}
\begin{itemize}
\item No unique meromorphic zeta or Gamma on $\No[i]$. The constructions are
 not identified with each other; $\zD$ is not a continuation of $\zp$; no
 surcomplex GRH.
\item The Stirling and Hurwitz values are a selected normalization. Flat
 alternatives exist, and no uniqueness, convexity or Gauss claim is made for
 the correction.
\item No values at infinite critical heights from local data. Phases are
 underdetermined, and they do not give strong summability. Riemann--Siegel
 at infinite $T$ is not an ordinary sum.
\item The transported model depends on the ultrafilter and the embedding.
 ${}^*\N\ne\Oz$, and the poles are at internal integers only. The transferred
 RH is not a new conjecture.
\item Only quantified estimates transfer, and bounded verifications do not
 reach infinite heights. No density statement at infinitesimal
 $\sigma-\tfrac12$.
\item \Cref{gz:thm:protected} proves neither RH nor simplicity, and
 $0.417293$ must not be used in it.
\item The 2026 preprints are marked, unverified and not needed locally.
\item The escape clause of the heat theorem is conditional on RH and
 simplicity, and it gives infinite modulus only, not direction or rate.
\item Priority not established; no resolution of named open problems;
 unrefereed; no Lean. The code checks finite
 identities. The repository audit covered the catalogue and summaries only.
\end{itemize}

\paragraph{05.}
\begin{itemize}
\item No proof of RH, no numerical improvement, no canonical zeta or $\xi$ on
 $\No[i]$.
\item The lift is canonical only on finite arguments and undefined at actual
 poles; punctured pole monads are in the domain.
\item The ratio cocycle is formal and phase-free. No claim that exponentially
 small corrections vanish.
\item Bare Stirling is not an exact global log-Gamma. $\Efp$ does not define
 $\exp(iT)$ for infinite $T$. The flat correction fixes only a modulus on one
 line, and supplies neither a phase nor a global Gamma.
\item The gauge $\lambda e^{-\pp}$ carries no convexity or multiplication claim.
\item $\Tp$ misses the strip at infinite height, and there is no uniqueness of
 continuation between the finite and $\Tp$ regions.
\item \Cref{gz:thm:weak}(b) is not a counterexample to RH and not an obstruction
 to a more rigid category.
\item Elementary realizations are noncanonical; ${}^*\exp\ne\exp_{\No}$ in
 general; internal counts are not external cardinalities or Hahn sums; zeros
 are cofinal in the set-sized field $F$ (the image of ${}^*\R$), not in $\No$.
\item Transfer holds for fixed ordinary parameters only, with no shortcut at
 the endpoint $\sigma=\tfrac12+\eps$ or $N(T)\to N(\st T)$.
\item A proportion is not an all-zeros assertion.
\item Benchmarks are selected published results, not records.
\item In 05's heat theorem only the infinite-modulus clause uses RH and
 simplicity, and no rate is given.
\item RH equivalences through prime counting, derivatives, positivity or pair
 correlation transfer only after every function, sum and index set they use
 is named; no unspecified ``surreal distribution of zeros'' is used.
\item No exhaustive priority claim; targeted audit; no Lean. Finite checks do
 not certify the support lemmas, quantifiers or transfers.
\end{itemize}

\paragraph{06.}
\begin{itemize}
\item No proof of RH and no unique global zeta. Unrefereed, no proof assistant,
 no priority for every consequence. The repository audit was targeted.
\item The Gamma nonuniqueness, the convexity classification and the finite-phase
 tube are not new.
\item ``Canonical'' means the explicit Taylor--Laurent rule.
\item A strong sum is not a fine limit. The formal, local-analytic, fine and
 Berarducci--Mantova derivatives are distinct, and no chain rule for the last
 is asserted.
\item Infinitesimal residuals are not zeros.
\item The robust criterion proves neither RH nor simplicity.
\item Stable proportions are centre-indexed.
\item The results on $\Tp$ say nothing about the infinite-height strip.
\item $\zH$ is not a strong sum of $(A+n)^{-s}$. There is no global zeta by
 substituting $A=1$, and no unrestricted complex-parameter prescription
 without a named phase.
\item $\Efp(\LSt)$ is meaningful only on $\Om$; elsewhere a phase must be chosen.
\item The Gamma obstruction does not rule out every global Gamma, and the
 resonance theorem does not exclude every global zeta.
\item The functions of \cref{gz:thm:weak}(a) are not proposed completions.
\item Transported operations are not canonical ones, and finite-zero agreement
 is not equality of values.
\item ${}^*N$ is internal; big-$\bigO$ needs ordinary witnesses; no statement at
 $\sigma=\tfrac12+\eps$.
\item Platt--Trudgian is cited, not recomputed, and is not the latest record.
 06 does not infer simplicity from it. This report cites the published
 simplicity statement separately, as 05 does, and uses it only in
 \cref{gz:cor:verified}.
\item The proportion preprints are unverified. The conclusions are parametric,
 and the $\liminf$ gives only every ordinary $c$ below the accepted bound
 $c_*$.
\item No surreal primes or prime counting in a bare surreal field.
\item Finite Hermitian tests gain nothing from surreal scalars.
\item The heat escape has no scale, and its escape clause is conditional.
\item The checks are finite algebra only.
\item Workspaces are set-sized, and no single Hahn field is closed under all
 operations.
\end{itemize}

\paragraph{The merge.}
\begin{itemize}
\item \Cref{gz:prop:resonance}(c) and \cref{gz:prop:wingfail} use the axiom of
 choice, for a basis of $\R$ over $\Q(\log3/\log2)$ and for a Hamel basis.
\item Whether a coherent galaxy zeta exists for a phase without resonance is
 open (\cref{gz:q:nonresonant}).
\item The class $\DJs$ is not claimed to be maximal for \cref{gz:thm:stability}.
\item The combined dilation-and-shift independence is open.
\item The content of Pong's 2020 Addendum \cite{PongAdd} was not checked.
 Whether other classical work covers 02's dilation family was not searched.
\item At the two exceptional abscissae $\sigma_\pm(T)$ of \cref{gz:thm:strip},
 for a general prescription with Stirling moduli, no reflection obstruction
 and no compatibility is claimed. For $E\circ\LSt$ itself reflection fails
 there as well (\cref{gz:thm:strip}(e)).
\item \Cref{gz:prop:resonance}(a) covers only phases with ordinary-circle
 character values. 03's statements for arbitrary phases, 04's prescribed
 prime phases and exponentials with character values in $1+\mm$ fall under
 \cref{gz:thm:resonance} only when they have a resonance.
\item The page range of the Guth--Maynard article is as cited by source 05; the
 Crossref record consulted shows the volume, issue and DOI.
\item The merge checks corroborate finite algebra and ordinary numerics. They
 are not shipped and are not proofs.
\end{itemize}

\appendix
\crefalias{section}{appendix}\crefalias{subsection}{appendix}

\section{Provenance and reconciliation record}\label{gz:app:provenance}

\subsection{The six manuscripts}
01 and 03 were prepared with ChatGPT, 02 and 06 call themselves
AI-assisted, and 04 and 05 call themselves research articles; 01, 02, 04,
05 and 06 name Vladimir Reshetnikov as the person they were prepared for.
Five are dated 22 September 2026 and 06 September 2026. None names a human
author. Each describes itself as unrefereed and not verified in a proof
assistant. Each cites a pinned commit of this repository and a targeted
reading of it.
\begin{center}\small
\begin{longtable}{@{}lP{0.36\textwidth}P{0.14\textwidth}P{0.34\textwidth}@{}}
\toprule
No. & Package & Pin; reading & Shipped here (code and data)\\
\midrule
\endhead
01 & 31-page article, \texttt{verification.py}, recorded results, CSV of inverse coefficients, source audit, build record & \pinA; catalogue, root README, gamma-functions README and introduction, opening of the analysis article & \texttt{01-inverse-zeta-*}: suite, results, CSV (CRLF line ends, kept byte-identical), requirements, audit, build record\\
02 & 29-page article, \texttt{checks/verify.py}, \texttt{checks/results.json}, \texttt{build.sh}, audit, validation & \pinA; catalogue, main README, gamma-functions README & \texttt{02-scale-separation-*}: suite, build script, results, audit, validation, requirements\\
03 & 34-page article, \texttt{code/verify.py}, verification and build records & \pinA; root README, catalogue, gamma-functions README and introduction & \texttt{03-horizontal-tube-*}: suite, verification, build record, requirements\\
04 & 31-page article, \texttt{code/verify.py}, verification, audit, build review & \pinB; catalogue, gamma-functions and analysis READMEs & \texttt{04-rh-transfer-*}: suite, verification, audit, build review\\
05 & 28-page article, \texttt{verify\_identities.py}, verification & \pinB; catalogue, analysis opening, gamma-functions README & \texttt{05-surcomplex-fields-*}: suite, verification, requirements\\
06 & 28-page article, \texttt{code/checks.py}, \texttt{build.sh}, checks, proof status, audit, build record & \pinB; catalogue, gamma-functions, analysis and trigonometry READMEs & \texttt{06-phases-and-rh-*}: suite, build script, checks, proof status, audit, build record\\
\bottomrule
\end{longtable}
\addtocounter{table}{-1}
\end{center}
No checksum or manifest file was delivered with any package. The manuscripts
themselves (their \LaTeX{} sources, PDFs and delivery READMEs) are not shipped.
The shipped \texttt{article.tex} is this report, rewritten from 06's.

\subsection{Where each result is printed}\label{gz:app:placement}
Every numbered result of every source is printed here, once, with a
second route where marked. Source numbers are the sources' own.
\begin{center}\small
\begin{longtable}{@{}lP{0.40\textwidth}P{0.48\textwidth}@{}}
\toprule
Src & Source results & Here\\
\midrule
\endhead
01 & Thm 1.1; Def 2.1; Lem 2.2, 2.3, 2.4; Ex 2.5; the fine-topology paragraph of Section 2.4 & \cref{gz:thm:headline,gz:def:strong,gz:lem:substitution,gz:lem:realpowers,gz:lem:weightmonoid}; the paragraph after \cref{gz:def:strong}; \cref{gz:rem:fine}\\
01 & Prop 3.1; Thm 3.2; Cor 3.3; eqs.\ 3.3--3.9; Warn 3.4 & \cref{gz:prop:functorial,gz:thm:divisor,gz:thm:finiteRH}; \cref{gz:sec:germs,gz:sec:finite-identities}; \cref{gz:warn:finite-raw}\\
01 & Prop 4.1, 4.2 & \cref{gz:prop:polar}; \cref{gz:prop:families}(a) with \cref{gz:prop:Echi}\\
01 & Prop 5.1; Warn 5.2; Thm 5.3, 5.4 & \cref{gz:prop:hurwitz-exists,gz:warn:hurwitz-raw,gz:thm:lerch,gz:thm:hurwitz-identities} (second route after it)\\
01 & Thm 6.1; eqs.\ 6.4--6.5; Thm 6.2; Thm 6.3, eq.\ 6.11 & \cref{gz:thm:stirling,gz:rem:secondroute}; \eqref{gz:eq:psiSt}, \cref{gz:thm:lerch}; \cref{gz:prop:phase,gz:prop:phase-locus}; \cref{gz:prop:gammachi,gz:ex:phasedep}, \cref{gz:sec:gammachi}\\
01 & Thm 7.1; Def 7.2; Thm 7.3, 7.4 & \cref{gz:thm:dirichlet,gz:def:zD,gz:thm:euler,gz:prop:Dtaylor}\\
01 & Lem 8.1; Thm 8.2; Cor 8.3, 8.4 & \cref{gz:lem:Vunique,gz:thm:fibre,gz:cor:realinverse,gz:cor:complexfibre}\\
01 & Lem 9.1; Thm 9.2, 9.3; Cor 9.4; Ex 9.5; numerics & \cref{gz:lem:multilagrange,gz:thm:arithmetic,gz:thm:integersector,gz:cor:rightmost}; end of \cref{gz:sec:inverse}\\
01 & Thm 10.1 and remark; Sections 11--12 & \cref{gz:prop:moving,gz:rem:stale01}; \cref{gz:q:preferred,gz:q:extend,gz:q:nonresonant}, \cref{gz:sec:nonclaims}; \cref{gz:app:checks}\\
\midrule
02 & Thm 1.1; Def 2.1, Lem 2.2, Rem 2.3 & \cref{gz:thm:indep-main}; \cref{gz:def:strong,gz:lem:substitution,gz:lem:mixed}; after \cref{gz:def:strong}\\
02 & Thm 3.1, Cor 3.2, Rem 3.3; eqs.\ 4.1--4.9; Thm 4.1; eqs.\ 4.10--4.12; Warn 4.2 & \cref{gz:prop:functorial,gz:thm:divisor,gz:rem:canonical}; \cref{gz:sec:germs,gz:sec:finite-identities}; \cref{gz:cor:finite-divisors,gz:thm:finiteRH}, \eqref{gz:eq:trivialslope}; \cref{gz:sec:coefficientwise,gz:warn:finite-raw}\\
02 & Thm 5.1, Rem 5.2, Prop 5.3; Thm 6.1, Cor 6.2, Prop 6.3, Warn 6.4 & \cref{gz:thm:dirichlet,gz:rem:dirichlet-scope,gz:prop:Dtaylor,gz:thm:euler,gz:prop:dichotomy}; \cref{gz:sec:wing}\\
02 & Lem 7.1, 7.2; Thm 7.3; Cor 7.4 & \cref{gz:lem:jacobian,gz:lem:jaccriterion,gz:thm:zetajets}\\
02 & Prop 8.1; eq.\ 8.6; Prop 8.2; Ex 8.3; Section 8.4 & \cref{gz:thm:stirling,gz:thm:translation,gz:prop:phase}; after \cref{gz:prop:phase}; after \eqref{gz:eq:psiSt}\\
02 & Lem 9.1; Thm 9.2; Cor 9.3; Lem 9.4; Thm 9.5; Cor 9.6 & \cref{gz:lem:polyleading,gz:thm:fieldleading,gz:lem:logscale,gz:thm:joint}\\
02 & Warn 10.1; Thm 10.2, 10.3; Cor 10.4 & \cref{gz:warn:hurwitz-raw,gz:thm:hurwitz-identities,gz:thm:lerch}\\
02 & Prop 11.1; Thm 11.2; Prop 11.3; Section 12; Section 13, Appendix A & \cref{gz:prop:families}(c), \cref{gz:prop:Echi,gz:def:Exp}; \cref{gz:prop:gammachi,gz:ex:phasedep}(b); \cref{gz:prop:wing,gz:cor:continuum}; \cref{gz:ex:omega}, \cref{gz:sec:gammachi}, \cref{gz:sec:germs}, \cref{gz:app:checks}; \cref{gz:sec:synthesis,gz:sec:nonclaims}, and for Appendix A (dependency audit) the proofs of \cref{gz:thm:dirichlet,gz:thm:zetajets,gz:thm:fieldleading,gz:lem:logscale,gz:thm:joint}\\
\midrule
03 & Thm 1.1; Thm 1.2 & \cref{gz:part:tube}; \cref{gz:thm:indep-slow}\\
03 & Def 2.1, Lem 2.2, Warn 2.3; Section 2.3 & \cref{gz:def:strong,gz:lem:substitution,gz:rem:fine}; \cref{gz:sec:derivatives}\\
03 & Prop 3.1; Thm 3.2; eq.\ 3.1; eqs.\ 3.2--3.4 & \cref{gz:prop:functorial,gz:thm:divisor}; \cref{gz:def:lift}; \cref{gz:sec:germs}\\
03 & Lem 4.1; Def 4.2; eqs.\ 4.2--4.3; Prop 4.3; Def 4.4; Thm 4.5; Thm 4.6 & \cref{gz:prop:polar,gz:def:phase,gz:def:Exp,gz:prop:Echi}; \cref{gz:thm:periods} (definition of $\cI$ and parts (a)--(b)); \cref{gz:thm:forced}(b)\\
03 & Prop 5.1; Thm 5.2; eqs.\ 5.12--5.14; Thm 5.3 & \cref{gz:prop:gammachi,gz:thm:translation}; \eqref{gz:eq:psiSt}, \cref{gz:ex:phasedep}(c); \cref{gz:thm:stirling,gz:prop:gammachi}\\
03 & Thm 6.1, 6.2; Prop 6.3 & \cref{gz:thm:gamma-tube,gz:thm:gauss-tube,gz:prop:tube-unique}\\
03 & Thm 7.1; Lem 7.2; eq.\ 7.6; Thm 7.3 & \cref{gz:thm:dirichlet}; \cref{gz:lem:mixed} with the proof of \cref{gz:thm:dirichlet}; \cref{gz:prop:Dtaylor,gz:def:zD}; \cref{gz:thm:euler}\\
03 & Lem 8.1; Thm 8.2; Section 8.2 & \cref{gz:lem:FEfactor,gz:thm:zeta-tube,gz:cor:canonical-zeros}\\
03 & Thm 9.1; Cor 9.2; eqs.\ 9.3--9.5; eqs.\ 9.6--9.7; Prop 9.3 & \cref{gz:thm:xi-tube,gz:cor:horizontalRH,gz:prop:cancel}; \cref{gz:prop:families}(b); \cref{gz:prop:moving}\\
03 & Lem 10.1, 10.2; Thm 10.3; Lem 10.4; Thm 10.5, 10.6 & \cref{gz:lem:chebyshev,gz:lem:primeindependence,gz:thm:indep-monomial,gz:lem:sloweval,gz:thm:indep-slow,gz:thm:gamma-zeta-independence}\\
03 & Prop 11.1; Thm 11.2, 11.3; Cor 11.4; Thm 11.5 & \cref{gz:warn:hurwitz-raw,gz:thm:hurwitz-unit,gz:thm:hurwitz-identities,gz:cor:hurwitz-unique,gz:thm:lerch}\\
03 & Section 12; Prop 13.1; Sections 13--14 & \cref{gz:ex:omega}, after \cref{gz:ex:phasedep,gz:thm:gamma-zeta-independence}; \cref{gz:ex:canonical-gamma,gz:cor:canonical-zeros}, after \cref{gz:prop:cancel,gz:cor:horizontalRH}; \cref{gz:thm:lerch}; \cref{gz:prop:wing}; \cref{gz:sec:synthesis,gz:sec:nonclaims}\\
\midrule
04 & Section 1; Lem 2.1; Def 2.2; Prop 2.3; Warn 2.4 & \cref{gz:sec:question,gz:sec:domains}; \cref{gz:lem:substitution,gz:def:lift,gz:prop:functorial}; \cref{gz:sec:derivatives,gz:rem:canonical}\\
04 & Thm 3.1; Cor 3.2; Ex 3.3; Cor 4.1; eqs.\ 4.1--4.11 & \cref{gz:thm:divisor}; \cref{gz:sec:germs}; \cref{gz:cor:finite-divisors}; \cref{gz:sec:germs,gz:sec:finite-identities,gz:sec:coefficientwise}, \eqref{gz:eq:trivialslope}\\
04 & Def 5.1, Thm 5.2; Lem 5.3; Thm 5.4; Cor 5.5; Section 5.3 & \cref{gz:def:RHfin,gz:thm:finiteRH,gz:lem:splitting,gz:thm:twosign}; \cref{gz:thm:stability}, \cref{gz:cor:criteria}(a); \cref{gz:rem:shadow,gz:cor:scale}\\
04 & Prop 6.1; Cor 6.2; Thm 6.3 & \cref{gz:prop:conservation,gz:thm:protection}; \cref{gz:rem:allh}, \cref{gz:cor:criteria}(a); \cref{gz:thm:protected}\\
04 & Thm 7.1; Prop 7.2; Thm 7.3 & \cref{gz:thm:dirichlet,gz:rem:dirichlet-scope,gz:prop:Dtaylor,gz:thm:euler}\\
04 & eq.\ 8.1; Thm 8.1; eq.\ 8.8; Warn 8.2; Prop 8.3; eq.\ 8.9 & \cref{gz:sec:stirling}; \cref{gz:thm:hurwitz-identities,gz:thm:lerch,gz:rem:secondroute}; \cref{gz:thm:lerch}; \cref{gz:warn:hurwitz-raw,gz:prop:flathurwitz,gz:prop:phase}\\
04 & Prop 9.1; Sections 9.2--9.3 & \cref{gz:prop:primephases}; \cref{gz:sec:galaxies,gz:sec:weak}\\
04 & Lem 10.1; Section 10.2; Thm 10.2; Prop 10.3; Sections 10.4--10.5 & \cref{gz:prop:embedding}; \cref{gz:sec:transported-divisors}; \cref{gz:thm:RHstar}(a),(b); \cref{gz:thm:RHstar}(d); \cref{gz:warn:transported,gz:sec:quantifiers}\\
04 & Section 11 & \cref{gz:sec:quantifiers,gz:prop:mty,gz:sec:counting,gz:sec:proportions,gz:sec:density,gz:sec:verified}\\
04 & eqs.\ 12.1--12.4; Thm 12.1; Prop 13.1, 13.2; Section 14; Appendices A, B & \cref{gz:sec:heat}; \cref{gz:thm:escape}; \cref{gz:prop:hermitian,gz:prop:margin}; \cref{gz:sec:synthesis}; \cref{gz:thm:protection,gz:cor:scale,gz:ex:omega}, \cref{gz:sec:gammachi}; \cref{gz:app:checks}\\
\midrule
05 & Lem 2.1; Prop 2.2; Prop 3.1; eqs.\ 3.2--3.9 & \cref{gz:lem:substitution,gz:warn:finite-raw,gz:rem:fine,gz:prop:functorial}; \cref{gz:sec:germs,gz:sec:finite-identities}\\
05 & Thm 4.1; Cor 4.2; Thm 4.3; Section 4.1 & \cref{gz:thm:divisor,gz:cor:finite-divisors,gz:thm:finiteRH}; \cref{gz:sec:germs}\\
05 & Def 5.1; Lem 5.2; Thm 5.3, 5.4; Rem 5.5 & \cref{gz:def:DJ,gz:lem:implicit,gz:thm:protection,gz:thm:stability}; \cref{gz:cor:criteria}(b); after \cref{gz:cor:scale}\\
05 & Section 6.1; Prop 6.1; eq.\ 6.4; Section 6.3, eq.\ 6.5; Thm 6.2, eqs.\ 6.6--6.12 & \cref{gz:prop:polar,gz:thm:stirling,gz:prop:phase}; \cref{gz:sec:nonunique}(2); \cref{gz:thm:cocycle,gz:thm:translation,gz:rem:cocycle-gloss}\\
05 & Thm 7.1; Cor 7.2 & \cref{gz:thm:strip}\\
05 & Thm 8.1; Def 8.2; Cor 8.3 and the expansion after it; Warn 8.4 & \cref{gz:thm:dirichlet,gz:def:zD,gz:thm:euler,gz:prop:Dtaylor,gz:rem:flat-dirichlet}\\
05 & Section 9.1; Thm 9.1; Section 9.3 & \cref{gz:part:infinite} (opening) and \cref{gz:sec:weak}; \cref{gz:thm:weak}(b); \cref{gz:sec:weak} (second subsection)\\
05 & Section 10.1; Prop 10.1; Warn 10.2; Thm 10.3; Prop 10.4; Section 10.3 & \cref{gz:sec:structure}; \cref{gz:prop:embedding,gz:warn:transported}; \cref{gz:thm:RHstar}(a); \cref{gz:thm:RHstar}(d); \cref{gz:sec:quantifiers,gz:sec:counting}, \cref{gz:thm:RHstar}(c)\\
05 & Section 11.1; eqs.\ 11.1--11.2, Prop 11.1; eq.\ 11.3, Cor 11.2; eqs.\ 11.4--11.7; eqs.\ 11.8--11.10, Prop 11.3; Section 11.7 & \cref{gz:sec:quantifiers}; \cref{gz:prop:mty}; \cref{gz:sec:verified,gz:cor:verified}; \cref{gz:sec:proportions,gz:cor:stableproportion,gz:sec:counting}; \cref{gz:sec:density,gz:prop:density}; \cref{gz:sec:positivity}\\
05 & eqs.\ 12.1--12.4; Thm 12.1; Section 13; Appendix A & \cref{gz:sec:heat}; \cref{gz:thm:escape}; \cref{gz:sec:synthesis}; \cref{gz:sec:germs,gz:ex:omega}\\
\midrule
06 & Lem 2.1; Section 2.3 & \cref{gz:lem:substitution}; \cref{gz:sec:sumnotions}\\
06 & Def 3.1; Prop 3.2; Thm 3.3; Cor 3.4; Prop 3.5; eq.\ 3.8 & \cref{gz:def:lift,gz:prop:functorial,gz:thm:divisor,gz:cor:finite-divisors}; \cref{gz:sec:finite-identities,gz:sec:coefficientwise}\\
06 & Def 4.1; Thm 4.2; Thm 4.3; Thm 4.4; Cor 4.5, Ex 4.6; Section 4.4 & \cref{gz:def:RHfin}; \cref{gz:thm:finiteRH}; \cref{gz:thm:protection}; \cref{gz:thm:stability} with \cref{gz:cor:criteria}(c); \cref{gz:cor:scale}; \cref{gz:sec:protected}, \cref{gz:cor:stableproportion}\\
06 & Thm 5.1; Def 5.2; Thm 5.3; Prop 5.4; eq.\ 5.13 & \cref{gz:thm:dirichlet,gz:def:zD,gz:thm:euler,gz:prop:Dtaylor,gz:rem:flat-dirichlet}\\
06 & eq.\ 6.1; Prop 6.1; eqs.\ 6.4--6.5; Thm 6.2; Ex 6.3; Warn 6.4; eq.\ 6.12 & \eqref{gz:eq:LSt}; \cref{gz:thm:stirling,gz:prop:gammachi}, 06's route \cref{gz:rem:secondroute}; \cref{gz:prop:phase}; \cref{gz:thm:hurwitz-identities,gz:thm:lerch}; \cref{gz:thm:lerch}; \cref{gz:warn:hurwitz-raw}; \cref{gz:sec:nonunique}(3)\\
06 & eqs.\ 7.1--7.2; Thm 7.1; Prop 7.2; eq.\ 7.5; coherent continuation; Thm 7.3; eq.\ 7.8; Rem 7.4 & \cref{gz:def:Exp}; \cref{gz:thm:forced}(a); \cref{gz:prop:twofreq}; \eqref{gz:eq:resonantpowers}; \cref{gz:def:coherent,gz:thm:resonance}; after it; \cref{gz:rem:resonance-scope,gz:prop:resonance}\\
06 & Thm 8.1; Prop 9.1; Thm 9.2 with the remark on ${}^*\Gamma$; Section 9.3 & \cref{gz:thm:weak}(a); \cref{gz:prop:embedding}; \cref{gz:thm:RHstar}, \cref{gz:sec:transported-divisors}; \cref{gz:sec:quantifiers}\\
06 & Prop 10.1; eqs.\ 10.1--10.9; Cor 10.2; Prop 10.3; Section 10.5 & \cref{gz:cor:halofree}; \cref{gz:sec:partial}; \cref{gz:cor:stableproportion,gz:prop:hermitian}; \cref{gz:sec:positivity}\\
06 & Lem 11.1; Thm 11.2; Section 12; Appendices A--B & \cref{gz:lem:heatpersist,gz:thm:escape}; \cref{gz:sec:synthesis}; \cref{gz:app:checks,gz:app:notation}\\
\bottomrule
\end{longtable}
\addtocounter{table}{-1}
\end{center}

\subsection{Where the merge had to choose}
\begin{itemize}
\item \emph{Base text and notation.} 06 is the base; 03 is the base of
 \cref{gz:part:tube}. Notation follows \cref{gz:sec:conventions}. $\SC$,
 $\Pi$, $\Exp$ and $\cI(\Psi)$ are taken from the collection, and every
 source symbol that collided is renamed (\cref{gz:tab:rename}).
\item \emph{Duplicates printed once.} Substitution (all six); divisor rigidity
 (all six, with the analysis report); finite RH (01, 02, 04, 05, 06); the
 Dirichlet realization and Euler identities (all six); Stirling recurrence
 and Gauss (01, 02, 03, 05, 06, the recurrence also by 04, and the
 gamma-functions report); the
 finite-phase criterion (all six, and the gamma-functions report); the
 Hurwitz identities (01, 02, 03, 04, 06); the embedding lemma and transferred
 RH (04, 05, 06); forced poles (03, 06); weak completions (05, 06); the heat
 theorem (04, 05, 06); the MTY, counting, proportion and density transfers
 (04, 05, 06).
\item \emph{Second routes kept.} The Hurwitz route to the Stirling identities
 (\cref{gz:rem:secondroute}); 01's operator proof of distribution after
 \cref{gz:thm:hurwitz-identities}; 06's sign choice for the stability
 converse (\cref{gz:cor:criteria}); local conservation (04) beside the
 involution argument (05, 06); the two independence mechanisms, which concern
 different families. The germ-level proofs of the Hurwitz shift and
 derivative identities in 01, 02 and 06 are not reprinted; they differ from
 the printed proof only in working with evaluated germs instead of
 $\C(s)$-formal series.
\item \emph{Statements chosen.} For forced poles, 06's statement about arbitrary
 functions as part (a), with 03's local form as part (b). For stability, the class $\DJs$
 contains all three source classes, with 04's one-$\eta$ converse. On Platt
 and Trudgian, the published simplicity statement is cited, with 06's
 abstention kept as a non-claim. Journal versions are cited where they exist.
\item \emph{Results added.} \Cref{gz:lem:evalhom},
 \cref{gz:prop:resonance}, the strip-wide form of \cref{gz:thm:strip},
 \cref{gz:prop:phase-locus,gz:prop:wing} (general form),
 \cref{gz:prop:wingfail,gz:cor:wingfibre}, the Gauss failure in
 \cref{gz:sec:nonunique}(2), and \cref{gz:thm:stability} as one theorem.
 After the post-merge audit: clause (e) of \cref{gz:thm:strip} and the
 constant-term step in the proof of \cref{gz:thm:stirling}. Further
 \merge{} steps complete proofs inside \cref{gz:lem:Vunique},
 \cref{gz:lem:primeindependence} and \cref{gz:q:nonresonant}.
\item \emph{A correction to the reconciliation plan.} The plan stated the
 strip defect as an obstruction at every $\sigma\in(0,1)$. The proof shows
 two exceptional abscissae $\sigma_\pm(T)$ where the modulus test is silent,
 and the theorem is stated with them. Clause (e) shows that the Stirling
 Gamma $E\circ\LSt$ itself fails reflection at every point of the strip,
 so the exceptions concern only a general prescription with Stirling moduli.
\end{itemize}

\subsection{Errors found in the sources and how they were fixed}\label{gz:app:errors}
No source contains a fatal error or a false main statement. Severity codes:
G, a proof gap in a true statement; S, a statement needing a hypothesis or
qualifier; W, wording; A, attribution or bibliography; F, framing that becomes
false in the merged context.
\begin{center}\small
\begin{longtable}{@{}rP{0.16\textwidth}lP{0.30\textwidth}P{0.34\textwidth}@{}}
\caption{Errors and fixes.}\label{gz:tab:errors}\\
\toprule
No. & Source, location & Sev. & Issue & Fix, and where\\
\midrule
\endfirsthead
\toprule
No. & Source, location & Sev. & Issue & Fix, and where\\
\midrule
\endhead
1 & 01 Prop 4.2 and \S4 & A/F & $E_0$ is Ehrlich--Kaplan's $\Exp$ and $E_\lambda$ is \rl{e:thm-twisted}; 01 calls $\Efp$ ``the canonical exponential'' & identified and credited; $\Efp$ renamed (\cref{gz:sec:expconv,gz:sec:canonical})\\
2 & 01 remark after Thm 10.1 & F & ``nothing canonical decided at $-\omega$''; ``a simultaneous FE would require an additional construction'' & with $\Exp$ the zero is present; 03 supplies the construction; no FE forces it (\cref{gz:rem:stale01})\\
3 & 01 Lem 8.1, existence & G & evaluation of weighted series used as a strongly additive ring homomorphism without a lemma & \cref{gz:lem:evalhom}, used in \cref{gz:lem:Vunique}\\
4 & 01 proof of Thm 6.2 & W & remainder omits $y^3/(6x^2)$ and $-y/(12x^2)$ & corrected expansion in \cref{gz:prop:phase}\\
5 & 01 Thm 6.2 ``defined iff'' & F & reads as ``undefined elsewhere'' once $\Exp$ exists & \cref{gz:prop:phase-locus}\\
6 & 01 Prop 5.1; text after Thm 6.3 & W & ``$1/(s-1)$ interpreted separately''; ``agreement is exact'' on disjoint domains & reworded (after \cref{gz:prop:hurwitz-exists}; after \cref{gz:prop:gammachi})\\
7 & 01 bibliography & A & Costin--Ehrlich cited as forthcoming; Gessel published; Ehrlich--Kaplan missing & journal data in the references\\
8 & 02 Lem 2.2(b) & G & the ``words'' argument treats $\eps^\alpha$ as monomials & $\supp\eps^\alpha\subseteq\langle S\rangle$ (\cref{gz:lem:mixed})\\
9 & 02 Prop 5.3 & W & ``$h$ sufficiently small'' is unnecessary & every infinitesimal $h$ (\cref{gz:prop:Dtaylor})\\
10 & 02 Lem 9.1 & G & positivity of the tail valuation needs one line & $q_n/q_{n_0}\prec x^{-D}$ (\cref{gz:lem:polyleading})\\
11 & 02 eq.\ 8.6 & G & joint summability in $(1/x,\delta)$ asserted & joint evaluation (\cref{gz:thm:translation})\\
12 & 02 Thm 11.2 & S & the functional $\ell$ depends on $x$ & stated (\cref{gz:ex:phasedep}(b))\\
13 & 02 Warn 6.4 & W & ``zero-freeness confined to $\Tp$'' against its own Prop 11.3 & ``for the raw $\zD$'' (\cref{gz:sec:wing})\\
14 & 02 throughout & A/F & $E_0=\Exp$ not identified; ``no canonical choice'' & canonical in the integer-part sense, not unique (\cref{gz:sec:expconv})\\
15 & 02 notation & W & $\Gamma_0$ and $\ell$ collide with the collection & renamed (\cref{gz:tab:rename})\\
16 & 02, 03 bibliographies & A & Pong and Costin--Ehrlich as preprints & journal data; Pong's Addendum listed\\
17 & 03 Lem 10.2 & G & works in $F(x)[[t]]$ although the objects live in $F[[x,t]]$ & the local ring $R=F(x)\cap F[[x]]$ (\cref{gz:lem:primeindependence})\\
18 & 03 Lem 10.2 framing & A & presented as self-contained; it is Pong's Thm 5.5 for $F=\C$ & credited; new content identified (after \cref{gz:lem:primeindependence})\\
19 & 03 \S2.3 & G & fine derivative of germs argued in outline & \rl{b:realization}, \rl{b:germalgebra} (\cref{gz:sec:derivatives})\\
20 & 03 Thm 4.5, phases, $\phi_\lambda$, $Z_\phi$ & A & repository results not credited & credited (\cref{gz:sec:phases})\\
21 & 03 bibliography & A & Ehrlich--Kaplan as arXiv & JSL 86 (2021) and erratum\\
22 & 04 \S5.3, \S14.1 & W & exact and shadow RH ``inequivalent in strength'' & ``formally stronger'' (\cref{gz:rem:shadow})\\
23 & 04 Prop 8.3; 05 \S6.3 & S & flat gauges break Gauss, unstated & stated and proved (\cref{gz:sec:nonunique})\\
24 & 04 \S9 and domain table & S & omits the extension of the Dirichlet chart under $\Exp$ & \cref{gz:prop:wing}\\
25 & 04 proof of Thm 8.1 & G & operator regrouping sketched & finite coefficient sums (\cref{gz:thm:hurwitz-identities})\\
26 & 04 Lem 2.1, Def 2.2 & G & transfer from the workspace to $\No[i]$ asserted & \cref{gz:lem:substitution} stated on $\SC$; \rl{a:cor:complexsub}\\
27 & 04 Thm 12.1(ii) & S & the nonreal zero is unconditional & hypotheses clause by clause (\cref{gz:thm:escape})\\
28 & 04 \S11.4 & A & AF's headline is $2/3$; $0.6725$ is their Montgomery--Taylor-window variant; the closed form $C_0$ (the Montgomery--Taylor constant) is stated in Lamzouri's Theorem 1.1 & \cref{gz:sec:proportions}\\
29 & 04 bibliography; 04 \S7 & A/W & Guth--Maynard and MTY as preprints; ``divergent-at-most-ordinary-points'' & journal data; phrase dropped\\
30 & 05 \S6.2 & W & ``usually leaves'' & ``always leaves'' (\cref{gz:prop:phase})\\
31 & 05 Cor 8.3 & W & $\zD'$ used before definition & defined first (\cref{gz:prop:Dtaylor})\\
32 & 05 Thm 9.1 & S & ``another such function'' fails if $F$ vanishes on $U$ & qualified (\cref{gz:thm:weak}(b))\\
33 & 05 Thm 5.3 & G & purely imaginary displacement coefficients unproved & the involution $d^*=-\bar d$ (\cref{gz:thm:protection})\\
34 & 05 Thm 6.2, eq.\ 6.12 & S & the Gamma-ratio gloss holds only off the negative real direction & \cref{gz:rem:cocycle-gloss}\\
35 & 05 Prop 3.1 & W & conjugation stated without $\bar f$ & \cref{gz:prop:functorial}\\
36 & 05 \S6 versus \S8 & W & ``right infinite'' for two sets & $\Sinf$ and $\Tp$ (\cref{gz:sec:domains})\\
37 & 05 README & W & ``no LaTeX warnings''; its build has one duplicate-destination warning & disclosed here; this build uses \texttt{pageanchor=false} around the title page\\
38 & 05 & A & $\Exp$ and repository labels not mentioned; Thm 4.1 duplicates \rl{c:p4:liftzeros}; MTY cited for the counting formula, which MTY only quotes & credited (\cref{gz:sec:collection,gz:thm:divisor,gz:sec:counting})\\
39 & 06 Rem 7.4, last sentence & G & ``an infinite period does not imply a resonance'' unproved & proved (\cref{gz:prop:resonance}(c))\\
40 & 06 Thm 7.3; Rem 7.4 & W & the theorem is stated for $\Exp$ and extended to resonant phases only in a remark; which phases have resonances is left open & general statement with proof (\cref{gz:thm:resonance}); existence and failure of resonances (\cref{gz:prop:resonance})\\
41 & 06 Thm 7.1 & A & cites the trigonometry report for infinite periods but not \rl{trigonometry:thm:infiniteperiods}; the explicit pole requirement is not stated & label cited; explicit pole requirement (\cref{gz:thm:forced})\\
42 & 06 \S5.3 & W & ``a genuine new domain result'' & ``a result on a new domain; no priority'' (after \cref{gz:thm:euler})\\
43 & 06 proofs of Thm 9.2, 11.2 & G & real constants not in the language & constants for all reals (\cref{gz:sec:structure})\\
44 & 06 Thm 11.2, second paragraph & W & canonical-lift clause reads as unconditional & ``under RH and simplicity'' (\cref{gz:thm:escape}(i))\\
45 & 06 \S7.3 & W & coherent continuation called ``the'' repository notion & it extends \rl{c:p4:def-HU} (\cref{gz:def:coherent})\\
46 & 06 bibliography & A & Guth--Maynard without status & \emph{Ann.\ of Math.}\ 203 (2026)\\
47 & 06 \S10.4 & W & one-sided count against Guth--Maynard's two-sided & 05's sentence added (\cref{gz:sec:density})\\
48 & 06 Prop 6.1 & S & $\Gamma_{E,\rm St}$ satisfies the recurrence but not reflection on $\Vs$, unstated; by \cref{gz:thm:strip}(e) it fails at every point of $\Vs$, for every $E$ extending $\Efp$, including $\sigma_\pm(T)$ & \cref{gz:thm:strip}(e) and the remark after it\\
49 & 03, 06 (external) & -- & Costin--Ehrlich eqs.\ (82)--(83) print $B_{2n}n!/((2n)!x^{n+1})$ in the extracted text & not used; only the framework is cited\\
50 & 01, 02, 04, 06 bibliographies & A & Rodgers--Tao, MTY, Platt--Trudgian, Costin--Ehrlich partly as preprints & journal data in the references\\
51 & 06 \S6.1 (a doubt raised in the review of 06) & -- & $\LSt$ prescribed at points $\eps+iT$ & stated as a prescription only (\cref{gz:sec:stirling})\\
52 & 03 (framing, \S14.1) & -- & the reflection obstruction is listed as a contribution; forced poles follow at once from the existence of infinite periods, and tube uniqueness from the definition of $g_\chi$ on $\Tm$ & said so (after \cref{gz:thm:forced}; \cref{gz:part:tube}; after \cref{gz:prop:tube-unique})\\
53 & 05 \S3.1; 06 after Def 3.1 & W & a nonzero infinitesimal displacement from a pole has a value ``usually'' of infinite modulus & ``always'', since the value is $g(0)\eps^m(1+\eta)$ with $m\le-1$ (after \cref{gz:def:lift})\\
\bottomrule
\end{longtable}
\end{center}
The sources' statements about the collection were checked against the
current tree and are accurate at their pins (\cref{gz:sec:collection}).
Item 53 was found by the post-merge audit below.

\subsection{Post-merge audit corrections}\label{gz:app:audit}
After the merge, three independent read-only audits checked the report: one
for fidelity to the six sources, one for non-claims, notation and
citations, and one for mathematical correctness. They found no false
theorem, no counterexample to a \merge{} result and no lost source result.
Their findings were applied as follows.
\begin{itemize}
\item \emph{Mathematics.} Clause (e) of \cref{gz:thm:strip}: for every $E$
 extending $\Efp$, $E\circ\LSt$ fails reflection at every point of $\Vs$,
 including $\sigma_\pm(T)$; the statements that ``nothing is claimed'' there
 are restricted to a general prescription with Stirling moduli. The
 relation between exact and shadow stability is stated correctly
 (\cref{gz:rem:shadow}). The source classes are subclasses of $\DJs$ but
 not nested. The constant term in the proof of Gauss in
 \cref{gz:thm:stirling} is computed. The scope of the resonance theorem is
 narrowed to phases with a resonance; the phase of
 \cref{gz:prop:resonance}(c) is built with choice, not explicitly; its
 character values are described exactly; the intersection describing
 resonances runs over $n\ge2$. The comparison of the independence theorems
 (\cref{gz:sec:indep-compare}) is corrected. ``Four criteria'' became the
 criteria of 04, 05 and 06, and the arccosine in \cref{gz:thm:strip}(c) is
 defined.
\item \emph{Fidelity to the sources.} Restored in their source generality: 04's
 Proposition 6.1 and Corollary 6.2 for every real entire $f$
 (\cref{gz:prop:conservation}, \cref{gz:rem:allh}); 05's Theorem 9.1 for
 every positive infinite $T$ (\cref{gz:thm:weak}(b)); 02's Warning 10.1 on
 all of $\SA$ (\cref{gz:warn:hurwitz-raw}); 02's Lemma 9.4 with all its
 conditions and proof steps (\cref{gz:lem:logscale}). Printed: 05's
 baseline recurrence, 01's multi-index formula and the derivation of its
 headline terms, 01's third research direction, the converse clause of 02's
 Remark 2.3, 05's pair-correlation caveat and 01's reworded Gamma gluing
 sentence (item 6).
\item \emph{Credits and provenance.} 06 cites the trigonometry report, and its
 Remark 7.4 anticipates the extension of its Theorem 7.3. The proof routes
 of the Stirling identities, the two parts of \cref{gz:thm:forced}, the
 values in \cref{gz:ex:canonical-gamma} not printed by 03, 05's gloss of
 $R_{1/2,0}$, the Montgomery--Taylor constant, and further credits to 02, 04
 and 06 are attributed correctly. The dates, authorship and reading lists of
 the sources are stated as they are. A dropped second route is disclosed.
 References to internal review notes, an unused phantom label, a quotation
 that no source wrote and two paraphrases printed as quotations were
 removed or replaced by verbatim quotations. Two collection labels were
 corrected (\rl{found:eq:omnific}; the holomorphic scope of
 \rl{c:p4:def-lift}).
\item \emph{Pointers.} Equation numbers taken from review notes were checked
 against the source PDFs and corrected (02's (8.6); 05's (6.4), (6.11) and
 (6.12); 03's (9.3)--(9.7); 04's (4.1)--(4.11) and (8.8)). Several rows of
 \cref{gz:app:placement} were corrected or completed, and rows 6, 11, 17, 28,
 34, 40, 41, 46, 48, 51 and 52 of \cref{gz:tab:errors} were corrected.
\item \emph{Notation.} Symbols reserved in \cref{gz:sec:conventions} now keep
 one meaning; the local objects that had reused them were renamed, as
 listed in the last row of \cref{gz:tab:rename}. In particular $\kappa_0$
 is the line proportion, the de Bruijn--Newman constant is
 $\Lambda_{\rm DN}$, and \cref{gz:thm:protected,gz:cor:stableproportion}
 use their own constants $c$, $c_*$.
\item \emph{Non-claims and caveats.} Restored at the point of use and in
 \cref{gz:sec:nonclaims}: the two exceptions and clause (e) of the strip
 theorem at headline level; that tube reflection and the functional equation
 are definitions on $\Tm$; that forced poles lie on the real axis; that
 \cref{gz:cor:verified} uses the published Platt--Trudgian simplicity
 statement; that only two clauses of \cref{gz:thm:escape} are conditional;
 that the Pong Addendum was not checked; and the dropped non-claims of 02,
 03, 04 and 05.
\end{itemize}
The corrected statement of \cref{gz:thm:strip}(e) was checked numerically for
this revision (\cref{gz:app:checks}).

\section{Verification suites and checks}\label{gz:app:checks}
The six delivered suites are shipped byte-identical in \texttt{code/}, with
their recorded outputs in \texttt{data/}. They check finite polynomial,
rational and truncated-series identities, and in 01 also ordinary-complex
numerics. They do not check strong summability, class quantification,
independence of infinite families, continuation, RH, or any cited paper.
\begin{center}\small
\begin{longtable}{@{}lrP{0.40\textwidth}P{0.30\textwidth}@{}}
\toprule
Suite & Checks & Content & Output behaviour\\
\midrule
\endhead
01 & 61 & inverse sector (37 multisets, 23 support values, the $V$ equation, log formula, 9 table entries, 30 integer-sector identities); Hurwitz and Stirling coefficients; 5 branches at $u=10^{-8}$, 100 digits & rewrites \texttt{verification-results.json} and \texttt{inverse-coefficients.csv} (unprefixed) next to itself\\
02 & 189 & Stirling recurrence and Gauss; finite shift; Hurwitz difference; Lerch coefficients; Bernoulli values; M\"obius and Euler-log coefficients; Jacobian samples & rewrites \texttt{results.json} next to itself\\
03 & 1670 & Bernoulli translation and distribution; Gauss coefficients; Hurwitz identities and values; log-Gamma bridge; arithmetic identities to $n\le256$; a Jacobian block; sine products & default \texttt{verification-rerun.json} in the working directory; refuses an existing path\\
04 & 350 & Stirling shift defect; Lerch coefficients; Bernoulli values; M\"obius and Euler-log to $128$; root motion; splitting; heat polynomials & \texttt{--output} is mandatory\\
05 & 296 & Stirling coefficients; recurrence defect; ratio cocycle; $B_{\rm odd}(\tfrac12)=0$; shifted real parts; convolution identities; Ingham checks & default \texttt{verification-rerun.json} in the working directory\\
06 & 221 & Stirling shift; Hurwitz shift and values; Hurwitz--Stirling bridge; root deformation; Dirichlet identities; quartet symmetry & prints only, unless \texttt{--output} is given\\
\bottomrule
\end{longtable}
\addtocounter{table}{-1}
\end{center}
All six were rerun on a copy of this directory for the merge (Python 3.14.4,
SymPy 1.14.0, mpmath 1.3.0), and all passed with the recorded counts. 01's
and 02's rewritten files agreed with the shipped records except for the
recorded interpreter version (recorded runs: Python 3.13.5). The two build
scripts, \texttt{02-scale-separation-build.sh} and
\texttt{06-phases-and-rh-build.sh}, were written for the source manuscripts,
which are not shipped; the README says how to build this report instead.
Some checks are near-tautological, for example the Lerch coefficients
through $(d/ds)(s)_m|_0=(m-1)!$ and the von Mangoldt convolution.

For the reconciliation, a separate suite of 29 checks passed. It checks the
strip identity of \cref{gz:thm:strip}(a) symbolically through $T^{-14}$ and the
flat defect numerically. It also checks the value of 01's left zeta at
$-\omega$ against the formula for $\sin_\lambda$, the two Hurwitz distribution
forms, the Lerch identity, 03's (5.13) and 05's cocycle leading term, the
literature constants ($C_0$, $(1+C_0)/2$, the crossover $7/10$, $A(3/4)=5/9$,
$0.417293>5/12>0.407511$), the trivial-zero slopes, and $\xi(0)=\xi(1)=\tfrac12$.
For this report a further 22 checks passed. They check the Bernoulli
reflection behind \cref{gz:thm:strip}(a) (to $n=29$), the Stirling sum at
$T=200$ to $10^{-40}$, the two exceptional abscissae, and the agreement of
04's and 06's second-order root motion. They also check the J-symmetry of the
Taylor coefficients of $\xi$ at its first zero, the decreasing exponents of
\cref{gz:prop:wingfail}, the cocycle identities, and the three Hurwitz
normalizations. For clause (e) of \cref{gz:thm:strip}, added after the
post-merge audit (\cref{gz:app:audit}), a further check passed. It confirms
symbolically the cancellation of the inverse powers through $n=24$ and the
elementary sum of (a). With $P=\exp(\log2\pi-\pi T+i\pi(\sigma-\tfrac12))$,
the common value of $E(\LSt(z)+\LSt(1-z))$, it confirms
$P\,\Sin(\pi z)/\pi=1-e^{2\pi iz}$ to $10^{-50}$ at $T\in\{0.7,2,3,7,15\}$
and eight abscissae including $\sigma_\pm(T)$, where the factor has modulus
$1$ and argument $\pm e^{-2\pi T}(1+\eps)$. It also checks the truncated
Stirling sum at $T=30$ against (a) to $3\cdot10^{-83}$, and
$P/(\Gamma(z)\Gamma(1-z))=1-e^{2\pi iz}$ with the classical Gamma. These
suites are corroboration only. They are not shipped and are not proofs.

\section{Notation index}\label{gz:app:notation}
\begin{center}\small
\begin{longtable}{@{}P{0.22\textwidth}P{0.70\textwidth}@{}}
\toprule
Symbol & Meaning\\
\midrule
\endhead
$\SC=\No[i]$ & surcomplex field; class-sized, with set-sized supports in every calculation\\
$\OO,\mm$; $\st$ & finite elements and infinitesimals; standard part on $\OO$\\
$\pp,\fin$; $\Pi$; $\Oz$ & purely infinite and finite parts ($\fin$ keeps infinitesimals); purely infinite reals; omnific integers $\Pi\oplus\Z$\\
$U^\#$, $c+\mm$, $b+\OO$ & halo, monad, galaxy\\
$\Sh,\Ph,\Scoef$ & strong sum, strong product, coefficientwise sum\\
$\Ex,\Lgm$; $\cis$; $\Efp$ & small exponential and logarithm; finite angle; finite-phase exponential\\
$\Psi_\chi,E_\chi,\Sin_\chi,\Cos_\chi$ & phase from a character $\chi:\Pi\to\TT(\No)$ and its functions\\
$\Exp,\Phi$ & canonical (Ehrlich--Kaplan) exponential and phase, $\chi=1$\\
$E_\lambda,E_{\lambda,M},E_\theta$ & explicit phase families (\cref{gz:prop:families})\\
$\cI(\Psi)$ & normalized period group $\{x:\Psi(2\pi x)=1\}$\\
resonance & $T\ne0$ with $\Psi(T\log n)=1$ for all ordinary $n$\\
$f^\#$ & canonical Taylor--Laurent lift\\
$\xi,\Xi,H_t$ & completed zeta; $\Xi(w)=\xi(\tfrac12+iw)$; heat family, $H_0(z)=\Xi(z/2)/8$\\
$\Phi_{\rm RT}$, $\Lambda_{\rm DN}$ & Rodgers--Tao kernel of $H_t$; de Bruijn--Newman constant\\
$\mu$, $\Lambda$ & M\"obius and von Mangoldt functions\\
$H_m$ & harmonic numbers $\sum_{j\le m}1/j$\\
$\kappa_0,\kappa_{\rm s}$ & $\liminf$ proportions of critical-line zeros, and of simple ones\\
$K_\Gamma$, $\Arith$, $\mathfrak A$, $\mathfrak M^*$ & set-sized Hahn workspace; arithmetic functions; first-order structure and its elementary extension\\
$\DJs$, $\Ssym$ & J-symmetric formal deformations; entire $h$ with both symmetries\\
$\Tp,\Tm,\Th$ & positive and negative infinite real part with finite imaginary part; horizontal tube\\
$\Wing,\Sinf,\Om,\Vs$ & wing $\Re s>\R$; Stirling half-plane; finite-phase locus; vertical strip of infinite height\\
$\DD_s$, $\zD$ & Dirichlet evaluation and Dirichlet zeta on $\Tp$\\
$\LSt,\SSt,\GSt,\psiSt$ & Stirling logarithm, its inverse-power part, real Stirling Gamma, digamma\\
$\Gamma_\chi,\Gamma_{\rm fp},\Gamma_{E,\rm St}$ & phase Gamma; finite-phase Gamma on $\Om$; 06's Gamma for any homomorphic $E$\\
$Q_a,Q_{a,b},R_{a,b}$ & translation remainder, cocycle, exponentiated cocycle\\
$\zeta_\chi,\xi_\chi,\FE_\chi$ & zeta and $\xi$ on $\Th$; reflection factor\\
$\zeta^\Wing_\chi$ & wing sum\\
$\zH,\zHh,\SA$ & Hurwitz zeta at infinite parameter, normalized form, phase-free parameter set\\
${}^*f$, ${}^*N$ & transported function, internal count\\
\bottomrule
\end{longtable}
\addtocounter{table}{-1}
\end{center}

\begin{thebibliography}{99}
\small
\raggedright
\setlength{\itemsep}{3pt}
\addcontentsline{toc}{section}{References}

\bibitem{Conway}
J. H. Conway, \emph{On Numbers and Games}, second edition, A K Peters, 2001.

\bibitem{Gonshor}
H. Gonshor, \emph{An Introduction to the Theory of Surreal Numbers}, London
Mathematical Society Lecture Note Series 110, Cambridge University Press,
1986. \href{https://doi.org/10.1017/CBO9780511629143}{doi:10.1017/CBO9780511629143}.

\bibitem{Neumann}
B. H. Neumann, On ordered division rings, \emph{Trans.\ Amer.\ Math.\ Soc.}\
\textbf{66} (1949), 202--252.
\href{https://doi.org/10.1090/S0002-9947-1949-0032593-5}{doi:10.1090/S0002-9947-1949-0032593-5}.

\bibitem{vdDE}
L. van den Dries and P. Ehrlich, Fields of surreal numbers and exponentiation,
\emph{Fund.\ Math.}\ \textbf{167} (2001), 173--188.
\href{https://doi.org/10.4064/fm167-2-3}{doi:10.4064/fm167-2-3}.

\bibitem{CE}
O. Costin and P. Ehrlich, Integration on the surreals, \emph{Adv.\ Math.}\
\textbf{452} (2024), article 109823.
\href{https://doi.org/10.1016/j.aim.2024.109823}{doi:10.1016/j.aim.2024.109823};
\href{https://arxiv.org/abs/2208.14331}{arXiv:2208.14331} (v5 used by the
sources). Cited for the extension framework and its Gamma example; its
displayed Stirling tail in eqs.\ (82)--(83) of the extracted v5 text is not used.

\bibitem{EK}
P. Ehrlich and E. Kaplan, Surreal ordered exponential fields, \emph{J.\ Symbolic
Logic} \textbf{86}(3) (2021), 1066--1115.
\href{https://doi.org/10.1017/jsl.2021.59}{doi:10.1017/jsl.2021.59}; erratum,
\emph{J.\ Symbolic Logic} \textbf{87}(2) (2022), 871;
\href{https://arxiv.org/abs/2002.07739}{arXiv:2002.07739}. Proposition 3.2
and Section 11 (Lemma 11.1, Proposition 11.2, Theorem 11.1).

\bibitem{BM}
A. Berarducci and V. Mantova, Transseries as germs of surreal functions,
\emph{Trans.\ Amer.\ Math.\ Soc.}\ \textbf{371} (2019), 3549--3592.
\href{https://doi.org/10.1090/tran/7428}{doi:10.1090/tran/7428}.

\bibitem{BMder}
A. Berarducci and V. Mantova, Surreal numbers, derivations and transseries,
\emph{J.\ Eur.\ Math.\ Soc.}\ \textbf{20} (2018), 339--390.
\href{https://doi.org/10.4171/JEMS/769}{doi:10.4171/JEMS/769}.

\bibitem{Pong}
W. Y. Pong, Applications of differential algebra to algebraic independence of
arithmetic functions, \emph{Acta Arith.}\ \textbf{172}(2) (2016), 149--173.
\href{https://doi.org/10.4064/aa8112-12-2015}{doi:10.4064/aa8112-12-2015};
\href{https://arxiv.org/abs/1504.03263}{arXiv:1504.03263}. Section 2,
Theorem 5.5, Corollary 5.6, Example 5.5.

\bibitem{PongAdd}
W. Y. Pong, Addendum to ``Applications of differential algebra to algebraic
independence of arithmetic functions'', \emph{Acta Arith.}\ \textbf{196}(3)
(2020), 325--327.
\href{https://doi.org/10.4064/aa200210-14-7}{doi:10.4064/aa200210-14-7}. Listed
for completeness; its content was not checked.

\bibitem{Gessel}
I. M. Gessel, Lagrange inversion, \emph{J.\ Combin.\ Theory Ser.\ A}
\textbf{144} (2016), 212--249.
\href{https://doi.org/10.1016/j.jcta.2016.06.018}{doi:10.1016/j.jcta.2016.06.018};
\href{https://arxiv.org/abs/1609.05988}{arXiv:1609.05988}.

\bibitem{Kawalec}
A. Kawalec, The inverse Riemann zeta function,
\href{https://arxiv.org/abs/2106.06915}{arXiv:2106.06915} (v3, 2022). Cited
as an existing inverse-zeta investigation.

\bibitem{DLMFgamma}
NIST Digital Library of Mathematical Functions, Chapter 5 (Gamma function),
\S\S5.2, 5.5, 5.7. \url{https://dlmf.nist.gov/5}.

\bibitem{DLMF511}
NIST Digital Library of Mathematical Functions, \S5.11, Asymptotic
expansions, in particular 5.11.1, 5.11.2 and 5.11.8.
\url{https://dlmf.nist.gov/5.11}.

\bibitem{DLMFzeta}
NIST Digital Library of Mathematical Functions, Chapter 25, \S\S25.2, 25.4,
25.6. \url{https://dlmf.nist.gov/25}.

\bibitem{DLMFzeros}
NIST Digital Library of Mathematical Functions, \S25.10, Zeros.
\url{https://dlmf.nist.gov/25.10}.

\bibitem{DLMF2511}
NIST Digital Library of Mathematical Functions, \S25.11, Hurwitz zeta
function, in particular 25.11.13--25.11.18 and 25.11.43.
\url{https://dlmf.nist.gov/25.11}. Bernoulli polynomials:
\url{https://dlmf.nist.gov/24.2}.

\bibitem{Titch}
E. C. Titchmarsh, \emph{The Theory of the Riemann Zeta-Function}, second edition,
revised by D. R. Heath-Brown, Oxford University Press, 1986.

\bibitem{PT}
D. Platt and T. Trudgian, The Riemann hypothesis is true up to $3\cdot10^{12}$,
\emph{Bull.\ London Math.\ Soc.}\ \textbf{53}(3) (2021), 792--797.
\href{https://doi.org/10.1112/blms.12460}{doi:10.1112/blms.12460};
\href{https://arxiv.org/abs/2004.09765}{arXiv:2004.09765}. The published
abstract states that the verified zeros are simple; arXiv v1 does not.

\bibitem{MTY}
M. J. Mossinghoff, T. S. Trudgian and A. Yang, Explicit zero-free regions for the
Riemann zeta-function, \emph{Res.\ Number Theory} \textbf{10} (2024),
article 11. \href{https://doi.org/10.1007/s40993-023-00498-y}{doi:10.1007/s40993-023-00498-y};
\href{https://arxiv.org/abs/2212.06867}{arXiv:2212.06867}.

\bibitem{PRZZ}
K. Pratt, N. Robles, A. Zaharescu and D. Zeindler, More than five-twelfths of
the zeros of $\zeta$ are on the critical line, \emph{Res.\ Math.\ Sci.}\
\textbf{7} (2020), article 2.
\href{https://doi.org/10.1007/s40687-019-0199-8}{doi:10.1007/s40687-019-0199-8}.

\bibitem{GM}
L. Guth and J. Maynard, New large value estimates for Dirichlet polynomials,
\emph{Ann.\ of Math.}\ \textbf{203}(2) (2026), 623--675 (pages as cited by
source 05).
\href{https://doi.org/10.4007/annals.2026.203.2.6}{doi:10.4007/annals.2026.203.2.6};
\href{https://arxiv.org/abs/2405.20552}{arXiv:2405.20552}. Theorem 1.2 and
equation (1.4).

\bibitem{AF}
L. Alp\"oge and R. Furman, More than two thirds of the zeta zeros are simple and
on the critical line, \href{https://arxiv.org/abs/2608.13637}{arXiv:2608.13637}
(v2, 19 August 2026). Preprint; not independently verified here.

\bibitem{Lamzouri}
Y. Lamzouri, A new proof that more than $2/3$ of the zeros of the Riemann zeta
function are simple and on the critical line,
\href{https://arxiv.org/abs/2609.02882}{arXiv:2609.02882} (v2, 8 September
2026), Theorem 1.1. Preprint; not independently verified here.

\bibitem{RT}
B. Rodgers and T. Tao, The de Bruijn--Newman constant is non-negative,
\emph{Forum Math.\ Pi} \textbf{8} (2020), e6.
\href{https://doi.org/10.1017/fmp.2020.6}{doi:10.1017/fmp.2020.6};
corrected version \href{https://arxiv.org/abs/1801.05914}{arXiv:1801.05914v5}
(3 July 2021).

\bibitem{RGamma}
\emph{Gamma functions over the surreals: a sharp convexity classification},
report in this collection, \texttt{docs/surreal/gamma-functions/}.

\bibitem{RTrig}
\emph{Trigonometry on the surcomplex plane}, report in this collection,
\texttt{docs/surcomplex/trigonometry/}.

\bibitem{RAnalysis}
\emph{Surcomplex analysis: germs, coherent Hahn sections, and infinitesimal zero
geometry}, report in this collection, \texttt{docs/surcomplex/analysis/}.

\bibitem{RFound}
\emph{Foundations for surreal and surcomplex mathematics}, report in this collection,
\texttt{docs/foundations-and-computation/foundations/}.

\end{thebibliography}
\end{document}
